% Fixed-depth predictive KL scaling for centered normalized networks.
% Compiles with plain pdflatex (run twice); uses only standard packages.
\documentclass[11pt,a4paper]{article}
\usepackage[utf8]{inputenc}
\usepackage[T1]{fontenc}
\usepackage{mathpazo}
\usepackage{amsmath,amssymb,amsthm}
\usepackage[margin=2.6cm]{geometry}
\usepackage{microtype}
\usepackage[hidelinks]{hyperref}
\setcounter{secnumdepth}{2}
\setlength{\emergencystretch}{3em}
\allowdisplaybreaks
\newtheorem{theorem}{Theorem}
\newtheorem{lemma}[theorem]{Lemma}
\newtheorem{proposition}[theorem]{Proposition}
\newcommand{\E}{\mathbb E}
\newcommand{\R}{\mathbb R}
\newcommand{\I}{\mathcal I}
\newcommand{\Pall}{P_{\mathrm{hidden}}}
\newcommand{\Pout}{P_{\mathrm{out}}}
\newcommand{\dsm}{\delta_{\mathrm{sm}}}
\newcommand{\ecut}{\eta_{\mathrm{cut}}}
\newcommand{\norm}[1]{\left\|#1\right\|}
\DeclareMathOperator{\tr}{tr}
\DeclareMathOperator{\KL}{KL}
\DeclareMathOperator{\Sym}{Sym}
\title{Predictive KL Scaling of Centered Normalized Networks\\[4pt]
\large A fixed-depth theorem in the dimension-first limit}
\author{}
\date{October 14, 2026 \quad (revised October 16, 2026)}
\hypersetup{pdftitle={Predictive KL Scaling of Centered Normalized Networks}}
\begin{document}
\maketitle
\begin{abstract}
For networks of fixed depth whose layer widths are proportional to the input dimension $d$, the average Bayes predictive KL of centered normalized tanh, sigmoid and ReLU activations has leading term $\Pall(d)/(2n)$. The proof assumes neither that the risk has a limit at fixed sample ratio nor that the posterior is locally Gaussian. A scalar Hermite condition gives tangent nondegeneracy and a global function comparison for the actual network. Tangent nondegeneracy yields the information lower bound; the global function comparison controls the information cost of a noisy auxiliary observation of the parameters, which vanishes in the limit. The main text derives the two information bounds and removes the auxiliary observation; the appendices prove the local estimates, the tangent geometry and the cross-network function comparison.
\end{abstract}

\section{Model, activation class and theorem}
\paragraph{Activations are centered and normalized.}
Fix $L\ge1$, $k_0=d$, $k_\ell/d\to\gamma_\ell\in(0,\infty)$, $\gamma_0=1$.
The layer-$\ell$ weights are $W^{(\ell)}\in\R^{k_\ell\times k_{\ell-1}}$ with i.i.d.\ standard Gaussian entries. The input $x\sim N(0,I_d)$ is independent of the parameters. Every layer uses the same fixed activation $\sigma_\ell\equiv\sigma$, with
\[
\E\sigma(G)=0,\qquad\E\sigma(G)^2=1,\qquad G\sim N(0,1).
\]
Write $W=(W^{(\ell)})_{\ell=1}^L$, $\theta=\theta(W)=(v,W)$, fix $v_i=1$, and define the network by
\begin{equation}\begin{aligned}
h^{(0)}(x)&=x,\qquad
h^{(\ell)}(x)=\sigma_\ell\!\left(\frac{W^{(\ell)}h^{(\ell-1)}(x)}{\sqrt{k_{\ell-1}}}\right),\\
F_\theta^{(L)}(x)&=\frac{v^{\mathsf T}h^{(L)}(x)}{\sqrt{k_L}}.
\end{aligned}\end{equation}
The true parameter $\theta^0=(v,W^0)$ is drawn from the prior $P_0$ above. The training set $\mathcal D_n=\{(x_\mu,y_\mu)\}_{\mu=1}^n$ is generated by
\[
y_\mu=F_{\theta^0}^{(L)}(x_\mu)+\sqrt\Delta z_\mu,\qquad
z_\mu\sim N(0,1),\quad\Delta>0
\]
Training and test inputs are i.i.d.; the noises are independent of each other, of the parameters and of all inputs. The Gaussian channel is
\[
\Pout(y\mid\lambda)=N(y;\lambda,\Delta).
\]
The learner uses the true prior and channel; its posterior is
\[
P(d\theta\mid\mathcal D_n)\propto P_0(d\theta)
\prod_{\mu=1}^n\Pout\!\left(y_\mu\mid F_\theta^{(L)}(x_\mu)\right).
\]
Only the hidden matrices are learned and the readout $v$ is fixed; the number of hidden parameters is
\[
\Pall(d)=\sum_{\ell=1}^Lk_{\ell-1}k_\ell,\qquad
c=\lim_d\frac{\Pall(d)}{d^2}
=\sum_{\ell=1}^L\gamma_{\ell-1}\gamma_\ell.
\]
The posterior predictive density and the average predictive KL are
\begin{equation}\begin{aligned}
p_{\mathcal D_n}(y\mid x)
&=\E_{\theta\mid\mathcal D_n}\Pout\!\left(y\mid F_\theta^{(L)}(x)\right),\\
\epsilon_{\log,d,n}
&=\E_{\theta^0,\mathcal D_n,x}\KL\!\left(
\Pout\!\left(\cdot\mid F_{\theta^0}^{(L)}(x)\right)
\,\middle\|\,p_{\mathcal D_n}(\cdot\mid x)\right).
\end{aligned}\end{equation}
Write $\epsilon_{\log,d}(\alpha):=\epsilon_{\log,d,\lfloor\alpha d^2\rfloor}$.
In the technical parts we write $p=\Pall(d)$ and regard $W$ as a $p$-dimensional vector of hidden-weight coordinates; gradients, Hessians, Gaussian integrals and OU couplings act only on these coordinates. We write $\theta'=\theta(V)$ for the parameter built from a second set of weights. Under prior averages the true hidden weights $W^0$ are abbreviated $W$; in the appendices $F_W=F_{\theta(W)}^{(L)}$ and $F=F_W$ both denote the actual network output unless the homogeneous comparator is explicitly indicated.

All constants may depend on the fixed depth, the activation, the width ratios and the noise, but not on $d,n$.

\subsection{Scalar activation conditions}
\paragraph{Let $\mathcal A_{\rm bal}$ be the union of the following two classes.}
Let $\operatorname{He}_m$ be the probabilists' Hermite polynomials,
\[
\mu_r(s)=\E[\sigma(sG)\operatorname{He}_r(G)],\qquad
b_m(s)=\frac{\mu_m(s)}{m!},\qquad
\sigma(z)=\sum_{r\ge0}\frac{\mu_r(1)}{r!}\operatorname{He}_r(z).
\]
Since $\sigma_\ell\equiv\sigma$, the Hermite coefficients of every layer are $\mu_{\ell,r}=\mu_r(1)$.
\emph{Smooth odd branch:} $\sigma$ is bounded, odd, of class $C^\infty$, with every fixed-order derivative bounded, and there is a fixed odd $m\ge5$ with
\begin{equation}\label{eq:scalar}
b_m(1)\neq0,\qquad\text{equivalently}\qquad
a_m:=\E[\sigma(G)\operatorname{He}_m(G)]/\sqrt{m!}=\sqrt{m!}\,b_m(1)\neq0.
\end{equation}
\emph{Piecewise-linear branch:} let $c_0=\sqrt{2/\pi}$ and
\begin{equation}\label{eq:pl}
\sigma(z)=az+b(|z|-c_0),\qquad
b\ne0,\qquad a^2+b^2(1-2/\pi)=1.
\end{equation}
This branch uses $m=6$ and requires $a_6\neq0$. Both branches are globally Lipschitz; the conditions concern a single real function and involve no Fisher, posterior or information-limit hypotheses.

\begin{theorem}\label{thm:main}
For the model above and any fixed $\sigma\in\mathcal A_{\rm bal}$,
\begin{equation}\label{eq:main}
\boxed{\lim_{\alpha\to\infty}\limsup_{d\to\infty}
\left|\alpha \epsilon_{\log,d}(\alpha)-\frac c2\right|=0.}
\end{equation}
The class contains centered normalized tanh of any fixed positive slope, the logistic sigmoid and centered normalized ReLU.
\end{theorem}
The limit $d\to\infty$ is taken first, then $\alpha\to\infty$. Existence of the risk limit at fixed $\alpha$ is not required. No claim is made for uncentered models, layerwise mixtures of the two branches, growing depth, or sharp upper bounds for non-Gaussian channels.

\section{Membership of the three activations}
\paragraph{One nonzero high-order Hermite coefficient suffices.}
The proof uses only $a_m\neq0$ and the parity relations $a_{m\pm1}=0$. Variations of row radii are controlled by the uniform concentration of the prior row radii around $1$ (Appendix~\ref{tg:section}); no condition on $b_m'$ is needed.

\subsection{tanh and sigmoid}
\paragraph{tanh with a fixed positive slope.}
Let $a>0$, $s_a^2=\E\tanh^2(aG)$, $\sigma_a(z)=\tanh(az)/s_a$. Oddness gives centering and the definition gives unit variance. Since $\tanh'=1-\tanh^2$, every fixed-order derivative is a polynomial in $\tanh$ times a constant, hence bounded.
Using the scalar identity
\[
\tanh(az)=2\int_0^\infty\frac{\sin(2atz)}{\sinh(\pi t)}\,dt,
\]
and
$\E[e^{iuG}\operatorname{He}_m(G)]=(iu)^m e^{-u^2/2}$, for odd $m=2r+1$ we get
\begin{equation}\label{eq:tanh-hermite}
b_m(s)=\frac{2(-1)^r}{m!s_a}
\int_0^\infty\frac{(2ast)^m e^{-(2ast)^2/2}}{\sinh(\pi t)}\,dt.
\end{equation}
The integral is strictly positive, so $(-1)^rb_m(s)>0$. Substituting $u=2ast$, it suffices to study
\[
H_s(u)=\frac1{2as\sinh(\pi u/(2as))}.
\]
With $v=\pi u/(2as)>0$,
\[
\partial_s\log H_s(u)=\frac{v\coth v-1}{s}>0.
\]
The last inequality follows from $v\cosh v-\sinh v>0$, which vanishes at zero and has derivative $v\sinh v>0$. On compact sets of positive scales Gaussian decay justifies differentiation under the integral, so $(-1)^rb'_m(s)>0$; this monotonicity is not needed later. In particular $(-1)^rb_m(1)>0$, and $m=7$ gives \eqref{eq:scalar}. To check the integral identity directly, use
\[
\frac1{\sinh(\pi t)}=2\sum_{j=0}^\infty e^{-(2j+1)\pi t}.
\]
For $x=az$, termwise integration gives
\[
2\int_0^\infty\frac{\sin(2xt)}{\sinh(\pi t)}\,dt
=\sum_{j=0}^\infty\frac{8x}{(2j+1)^2\pi^2+4x^2}.
\]
The interchange of sum and integral is dominated by $|\sin(2xt)|\le2|x|t$ and the summable Laplace integrals on the right. By the classical product formula
\[
\cosh x=\prod_{j=0}^\infty\left(1+\frac{4x^2}{(2j+1)^2\pi^2}\right),
\]
whose logarithmic derivative converges uniformly on real compact sets to the series above, the series equals $\tanh x$.

For $\ell_b(z)=(1+e^{-bz})^{-1}$,
\[
\ell_b(z)-\tfrac12=\tfrac12\tanh(bz/2).
\]
Hence centered normalized sigmoid equals $\sigma_{b/2}$, and its membership follows from the tanh computation.

\subsection{ReLU and the piecewise-linear branch}
\paragraph{Normalization and Hermite coefficients.}
Let $\varphi$ be the standard Gaussian density,
\[
\mu_R=\frac1{\sqrt{2\pi}},\qquad
v_R=\frac12-\frac1{2\pi},\qquad
\sigma_R(z)=\frac{z_+-\mu_R}{\sqrt{v_R}}.
\]
Since $z_+=(z+|z|)/2$, this corresponds to \eqref{eq:pl} with
$a=b=1/(2\sqrt{v_R})$. More generally, for \eqref{eq:pl} direct integration gives
\[
b_6(s)=\frac{bs\varphi(0)}{120},\qquad
b'_6(s)=\frac{b\varphi(0)}{120}\ne0.
\]
In particular $b_6(1)\neq0$, i.e.\ $a_6\neq0$.
For instance, using $\operatorname{He}_6(z)=z^6-15z^4+45z^2-15$ and
$\E|G|^{2j+1}=2^jj!\sqrt{2/\pi}$, one gets
$\E[|G|\operatorname{He}_6(G)]=6\varphi(0)$; then divide by $6!$.
All odd coefficients of degree $j\ge3$ vanish, so the degree-five and degree-seven neighbours of degree six vanish. For the smooth odd branch all even coefficients vanish at every scale.

Convolve $|z|$ with an even smooth kernel, then center and normalize under the Gaussian measure to obtain $\sigma_{\dsm}$. The derivative error is concentrated on $|z|\lesssim\dsm$; normalization adds a slope error $O(\dsm)$. Hence, on compact sets of scales,
\begin{equation}\begin{aligned}
\E(1+G^2)|\sigma_{\dsm}'(sG)-\sigma'(sG)|^2&\le C\dsm,\\
\E(1+G^4)|\sigma_{\dsm}''(sG)|^2&\le C/\dsm.
\end{aligned}\end{equation}
The second bound comes from a boundary layer of height $O(\dsm^{-1})$: its square times the length $O(\dsm)$. These scalar estimates are used to smooth auxiliary gradients only; the true likelihood is kept.

\section{Prior moments and local estimates}
\paragraph{The following moment bounds use only the prior.}
Let $q_\ell=\norm{h^{(\ell)}}^2/k_\ell$, $q_0=\norm X^2/d$. For fixed $r>0$ we need
\begin{equation}\label{eq:moments}\begin{aligned}
\E|q_\ell-1|^r&\le C_r d^{-r/2},\qquad
\sup_d\E|F_\theta^{(L)}(X)|^r<\infty,\\
\sup_d\E\norm{\nabla_WF_\theta^{(L)}(X)}^r&<\infty.
\end{aligned}\end{equation}

\subsection{Norm fluctuations and the output mean}
\paragraph{The conditional mean at finite width is kept.}
Set $\mu(q)=\E\sigma(\sqrt qG)$, $v(q)=\E\sigma(\sqrt qG)^2$. Lipschitz continuity and linear growth give the global bounds
\[
|\mu(q)|\le C|q-1|,\qquad |v(q)-1|\le C|q-1|,
\quad \E|\sigma(\sqrt qG)|^r\le C_r(1+q^{r/2}).
\]
Conditionally on the previous layers, the current rows are independent. If $Y_i$ are conditionally independent and centered, then in the expansion of $(\sum_iY_i)^{2s}$ terms with an index appearing once vanish, and the remaining terms have at most $s$ distinct indices. Hence
\[
\E\left|k^{-1}\sum_iY_i\right|^{2s}
\le C_s k^{-s}\times\text{(corresponding fixed-order moment bound)}.
\]
Induction with Jensen and linear growth gives $\sup_d\E q_\ell^{2s}<\infty$, and the same counting gives
\[
\E|q_\ell-v(q_{\ell-1})|^{2s}\le C_s d^{-s},
\quad
\E|q_\ell-1|^{2s}\le C_sd^{-s}+C_s\E|q_{\ell-1}-1|^{2s}.
\]
Induct from the chi-square moments of the input. The last layer decomposes exactly as
\[
F_\theta^{(L)}=\sqrt{k_L}\mu(q_{L-1})+
\frac1{\sqrt{k_L}}\sum_i\{h_i^{(L)}-\mu(q_{L-1})\}.
\]
The $2s$-th moment of the first term is at most $C_sk_L^s d^{-s}=O(1)$; by the same row counting the second term is $O(1)$. This gives the output moments and keeps the collective mean of the piecewise-linear branch.

\subsection{Gradients, Hessians and the OU remainder}
\paragraph{Normalized derivative moments are dimension-free.}
Let $A_\ell=W^{(\ell)}/\sqrt{k_{\ell-1}}$. All fixed-order moments of its operator norm are uniformly bounded. The forward and backward recursions give
\[
\beta_L=\boldsymbol1/\sqrt{k_L},\quad
\lambda_\ell=\sigma'(u^{(\ell)})\odot\beta_\ell,
\quad\beta_{\ell-1}=A_\ell^T\lambda_\ell,
\]
\[
\nabla_{W^{(\ell)}}F=\lambda_\ell(h^{(\ell-1)})^T/\sqrt{k_{\ell-1}}.
\]
All norms are controlled by fixed polynomials in the normalized input norm and finitely many operator norms, which gives the gradient moments. For the smooth branch the Hessian is the sum of the activation term
\[
J_{u,\ell}^T\operatorname{diag}(\beta_\ell\odot\sigma''(u^{(\ell)}))J_{u,\ell}
\]
and mixed weight--feature terms. For one ordering of directions, the squared HS norm of the latter is
\[
\frac{\norm{\lambda_\ell}^2}{k_{\ell-1}}
\norm{J_{h,\ell-1}}_{\rm HS}^2
\le\norm{\lambda_\ell}^2\norm{J_{h,\ell-1}}_{\rm op}^2.
\]
Hence $\E\norm{\nabla^2F}_{\rm HS}^2\le C$, with no extra factor $\Pall(d)$.

For the piecewise-linear branch the smoothing estimates are
\begin{equation}\label{eq:smoothing}\begin{aligned}
\E\norm{\nabla F_{\dsm}-\nabla F}^2&\le C\dsm,\\
\E\norm{\nabla^2F_{\dsm}}_{\rm HS}^2&\le C/\dsm,\qquad
\E\norm{\nabla F_{\dsm}}^4\le C.
\end{aligned}\end{equation}
At network level these require controlling boundary integrals weighted by backpropagated quantities; summing scalar errors column by column does not suffice. See Lemma~\ref{an:smoothing} in Appendix~\ref{an:section}.

Write $\theta(t)=\theta(W(t))$ and let $W(t)=\sqrt{1-t}W+\sqrt tZ$ be the prior-preserving OU coupling. The Gaussian Hermite expansion gives
\begin{equation}\label{eq:ou-rem}\begin{aligned}
F_{\theta(t)}^{(L)}-F_\theta^{(L)}&=\sqrt t\,(\nabla_W F_\theta^{(L)})^TZ+E_t,\\
\E\norm{E_t}_{L_X^2}^2&\le Ct^{1+\delta},\qquad
\delta=\begin{cases}1,&\text{smooth branch},\\1/2,&\text{piecewise-linear branch}.
\end{cases}
\end{aligned}\end{equation}
This stationary remainder includes the cancellation between the OU radial drift and boundary terms. It applies to the prior OU pair above, not to arbitrary posterior pairs.

\section{Geometric core: local tangents and global comparison}
\paragraph{The tangent spectrum gives the lower bound; the function comparison gives the upper bound.}
Write $\psi_j=\operatorname{He}_j/\sqrt{j!}$, $a_j=\E\sigma(G)\psi_j(G)$ and set
\[
N=m^L,\qquad S=1+m+\cdots+m^{L-1},\qquad B=a_m^S\ne0.
\]
Build the polynomial network with homogeneous activation $z^m/\sqrt{m!}$ in every layer, the original weights and the original readout; let $C_N(W)\in\Sym^N\R^d$ be its top input-degree coefficient. Define
\[
H_m(W)=I_N(C_N(W)),\qquad T_W=BH_m(W).
\]
Here $I_N$ is the multiple Wiener--It\^o integral, with
$\langle I_N(C),I_N(D)\rangle_{L_X^2}=N!\langle C,D\rangle$.
All tensors carry the globally symmetric Hilbert norm; grouped tensor norms cannot replace the actual function norm.

\begin{lemma}[Actual tangent bridge and homogeneous bulk spectrum]\label{lem:bridge}
For $L\ge2$ there is an orthogonal projection $P_W$ on the input space $L^2$ such that
\begin{equation}\label{eq:bridge}
\E\norm{P_W\Pi_NJ_W-B\nabla H_m(W)}_{\rm HS}^2\to0.
\end{equation}
With probability tending to one, $pB^2(\nabla H_m)^*\nabla H_m$ has all eigenvalues at least a fixed $\kappa>0$ except for $O(d)$ homogeneous gauge directions. For $L=1$ the radial/tangential ridge lemma gives the same spectral conclusion.
\end{lemma}
See Theorem~\ref{tg:actual-bridge} and Theorem~\ref{tg:homogeneous-bulk} in Appendix~\ref{tg:section}.

\begin{lemma}[Coupling-uniform comparison of actual functions]\label{lem:cross}
If $\Pi_d$ is the Gaussian parameter prior, then
\begin{equation}\label{eq:cross}
e_d:=\sup_{\pi\in\Gamma(\Pi_d,\Pi_d)}
\E_\pi|\langle T_W,F_{\theta'}^{(L)}-T_V\rangle_{L_X^2}|\longrightarrow0.
\end{equation}
The bound holds for every coupling with the prescribed prior marginals; $W,V$ need not be independent.
\end{lemma}
\begin{proof}[Structure of the proof]
First approximate the true input chaos of order $N$ under each prior. Let
\[
P_N(z)=\sum_{j=1}^N\frac{a_jz^j}{\sqrt{j!}},\quad
Q_N(z)=\sum_{j=1}^N\frac{\beta_jz^j}{\sqrt{j!}},
\quad\beta_j=\E[(\sigma(G)^2-1)\psi_j(G)].
\]
Let $R_{t,N}[A;P_N]$ be the function realized via $I_N$ by the \emph{formal degree-$N$ coefficient} of the depth-$t$ prefix with root polynomial $A$ and lower polynomial $P_N$. It is not the full order-$N$ projection of the untruncated polynomial network.
For the smooth odd branch,
\[
\Pi_NF_\theta^{(L)}=R_{L,N}[P_N;P_N]+o_{L^2(W,X)}(1).
\]
For the piecewise-linear branch let $\mu=bc_0/2$, $U_N=[Q_N-P_N^2]_{\le N}$; then
\begin{equation}\label{eq:rawclosure}\begin{aligned}
\Pi_NF_\theta^{(L)}={}&R_{L,N}[P_N;P_N]\\
&+\mu\sum_{t=1}^{L-1}\sqrt{k_L/k_t}\,R_{t,N}[U_N;P_N]
+o_{L^2(W,X)}(1).
\end{aligned}\end{equation}
$U_N$ may contain odd-degree terms, so the parity of ordinary piecewise-linear nodes cannot be imposed on it. The collective mean and quadratic restoration in \eqref{eq:rawclosure} are derived in Appendix~\ref{app:closure}.

With degree and depth fixed, only finitely many ordered degree trees and physical slot assignments remain. A lowest-cut induction for each branch shows: when paired with the balanced teacher, at depth $L$ only the all-$m$ tree survives, and every source with positive root degree at depth $t<L$ is $o(1)$. Each node contributes $a_m$ and there are $S$ nodes, so the surviving term is exactly $T_V$. The cut induction and the piecewise-linear quadratic exception are in Appendix~\ref{app:cuts}.

The finite-tree estimates hold deterministically on two separately defined marginal good events; bad events are controlled by fixed high moments, uniformly over couplings. If the actual approximation error is $E_d(V)$, then
\[
\E_\pi|\langle T_W,E_d(V)\rangle|
\le(\E\norm{T_W}^2)^{1/2}(\E\norm{E_d(V)}^2)^{1/2}=o(1).
\]
This uses only the marginals. Since $T_W$ lies in the chaos of order $N$ and all other input degrees are orthogonal to it, \eqref{eq:cross} follows.
\end{proof}

\subsection{From function comparison to parameter matching}
\paragraph{Polarization.}
For a coupling with the correct marginals let
$D=\E\norm{T_W-T_V}^2$, $B_F=\E\norm{F_\theta^{(L)}-F_{\theta'}^{(L)}}^2$.
Applying \eqref{eq:cross} to the coupling, its transpose and the two diagonal couplings gives
\[
|\E\langle T_W-T_V,F_\theta^{(L)}-F_{\theta'}^{(L)}\rangle-D|\le4e_d.
\]
Hence
\[
D\le\sqrt{DB_F}+4e_d\le(D+B_F)/2+4e_d,
\quad D\le B_F+8e_d.
\]
The homogeneous global matching theorem gives a finite family of measurable maps $\mathcal G_d$ with
\begin{equation}\label{eq:matching}\begin{aligned}
\log|\mathcal G_d|&=O(d\log d)=o(p),\\
\frac1p\E\min_{g\in\mathcal G_d}\norm{W-gV}^2
&\le C\E\norm{H_m(W)-H_m(V)}^2+o_d(1)\\
&\le C\E\norm{F_\theta^{(L)}-F_{\theta'}^{(L)}}^2+o_d(1).
\end{aligned}\end{equation}
The errors are uniform over all such couplings. Appendix~\ref{app:matching} gives the layerwise matching that preserves the error exponent. For even $m$ the maps may include row signs and need not be symmetries of the actual activation. Below, $gV$ serves only as a Gaussian mixture center; $F_{\theta(gV)}^{(L)}=F_{\theta'}^{(L)}$ is not required.

\section{From the actual spectrum to the lower bound}
\paragraph{The Fisher matrix of the Gaussian channel.}
Let $a_W(x)=\nabla_WF_\theta^{(L)}(x)$,
\[
A_W=\E_Xa_Wa_W^T,\qquad\I_W=A_W/\Delta,
\qquad D_d(r)=\E\frac1p\tr\{r(p\I_W+rI)^{-1}\}.
\]
By orthogonal projection, $A_W\succeq(P_W\Pi_NJ_W)^*(P_W\Pi_NJ_W)$. The squared singular-value perturbation inequality and Lemma~\ref{lem:bridge} give, for fixed $0<u<\kappa$,
\begin{equation}\begin{aligned}
\E\frac{\#\{\lambda_j(pA_W)\le u\}}p
\le{}&\Pr(E_d^c)+O(d/p)\\
&+\frac{\E\norm{P_W\Pi_NJ_W-B\nabla H_m}_{\rm HS}^2}
{(\sqrt\kappa-\sqrt u)^2}\longrightarrow0.
\end{aligned}\end{equation}
The noise only divides the threshold by $\Delta$. Splitting the regularized trace at the threshold gives
\begin{equation}\label{eq:deficit}
D_*(r):=\limsup_dD_d(r)\le Cr,\qquad0<r\le1.
\end{equation}
Exact zero directions form a fraction $O(d)/p=o(1)$ and do not affect the bound.

Appendix~\ref{app:information} derives from the mixture identity for the actual likelihood and Gaussian integration by parts the information inequality
\begin{equation}\label{eq:lower-smooth}
\epsilon_{\log,d,n}\ge\frac{p[1-D_d(r)]^2}{2(n+Cp/r^2)}.
\end{equation}
Taking $r=\alpha^{-1/4}$ and first $d\to\infty$,
\[
\alpha\liminf_d\epsilon_{\log,d}(\alpha)
\ge\frac{c[1-D_*(\alpha^{-1/4})]^2}
{2(1+Cc\alpha^{-1/2})}\longrightarrow\frac c2.
\]
For the piecewise-linear branch only the test field is smoothed; the true likelihood is kept. By \eqref{eq:smoothing},
\begin{equation}\label{eq:lower-pl}
\epsilon_{\log,d,n}\ge
\frac{p[1-D_d(r)-C\sqrt{\dsm}/r-C\sqrt{\dsm/r}]_+^2}
{2\{n(1+C\sqrt{\dsm/r})^2+Cp/(r^2\dsm)\}}.
\end{equation}
With $r=\alpha^{-1/8}$, $\dsm=\alpha^{-1/2}$ the additional errors are
\[
\sqrt{\dsm}/r=\alpha^{-1/8},\quad
\sqrt{\dsm/r}=\alpha^{-3/16},\quad
(\alpha r^2\dsm)^{-1}=\alpha^{-1/4}.
\]
By \eqref{eq:deficit}, for both branches
\begin{equation}\label{eq:lower-final}
\liminf_{\alpha\to\infty}\alpha\liminf_d\epsilon_{\log,d}(\alpha)\ge c/2.
\end{equation}
The sharp upper bound is proved independently below, so the lower bound needs no Tauberian argument.

\section{Coarse information bound and posterior contraction in function space}
\paragraph{A coarse contraction rate for posterior function distances.}
Write $I_{d,n}=I(\theta^0;\mathcal D_n)=I(W;\mathcal D_n)$. Bayes' rule and conditional independence of the samples give
\begin{equation}\label{eq:increment}
\epsilon_{\log,d,n}=I_{d,n+1}-I_{d,n},\qquad
\epsilon_{\log,d,n+1}\le \epsilon_{\log,d,n}.
\end{equation}
The second relation holds because conditioning reduces entropy: the true conditional noise entropy does not change with the sample size, while the predictive entropy decreases.

\subsection{Coarse risk bound}
\paragraph{A noisy parameter observation controls the cumulative information.}
Write $\theta_t=\theta(V_t)$. Let $V_t=\sqrt{1-t}W+\sqrt tU$ with $U$ an independent standard Gaussian. The Hermite OU identity and the gradient energy bound give
\[
\E\norm{F_\theta^{(L)}-F_{\theta_t}^{(L)}}_{L_X^2}^2\le Ct.
\]
The chain rule and the Gaussian label law with reference mean $F_{\theta_t}^{(L)}$ give
\begin{equation}\begin{aligned}
I_{d,n}&\le I(W;V_t)+I(W;\mathcal D_n\mid V_t)\\
&\le\frac p2\log(1/t)+\frac n{2\Delta}
\E\norm{F_\theta^{(L)}-F_{\theta_t}^{(L)}}_{L_X^2}^2\\
&\le\frac p2\log(1/t)+Cnt.
\end{aligned}\end{equation}
Taking $t=(1+2Cn/p)^{-1}$ yields
\[
I_{d,n}\le\frac p2\log(1+2Cn/p)+p/2.
\]
By \eqref{eq:increment}, $n\epsilon_{\log,d,n}\le I_{d,n}$, so
\begin{equation}\label{eq:coarse}
\limsup_d \epsilon_{\log,d}(u)\le C\frac{\log u}{u}
\quad(u\text{ large}).
\end{equation}
This bound uses only the prior gradient energy, not the leading term to be proved.

\subsection{From KL to function distance to matching distance}
\paragraph{Posterior samples and prior marginals.}
Given $\mathcal D_n$, draw $W,V$ independently; after averaging over the data each has law $\Pi_d$. With $B_X=F_\theta^{(L)}(X)-F_{\theta'}^{(L)}(X)$, the Gaussian Hellinger comparison gives
\begin{equation}\label{eq:hellinger}
\E[1-e^{-B_X^2/(8\Delta)}]\le2\epsilon_{\log,d,n}.
\end{equation}
The Hellinger distance satisfies the triangle inequality in the $L^2$ space of square-root densities. Using the posterior predictive density as an intermediate point and $(x+y)^2\le2x^2+2y^2$, the mean squared distance between two posterior samples is at most four times the single-sample squared distance. The latter is controlled by the KL, and the affinity of two Gaussians with equal variance is $e^{-B_X^2/(8\Delta)}$.

For fixed $q>2$, on $B_X^2\le8\Delta$ use
$1-e^{-x}\ge(1-e^{-1})x$ to bound the truncated second moment by $\le C\epsilon_{\log,d,n}$. On the complement, \eqref{eq:hellinger} controls the probability, \eqref{eq:moments} controls the moment of order $q$, and H\"older gives
\[
\E[B_X^2\boldsymbol1_{B_X^2>8\Delta}]
\le(\E|B_X|^q)^{2/q}\Pr(B_X^2>8\Delta)^{1-2/q}
\le C_q\epsilon_{\log,d,n}^{1-2/q}.
\]
Hence
\[
\E\norm{F_\theta^{(L)}-F_{\theta'}^{(L)}}_{L_X^2}^2
\le C_q(\epsilon_{\log,d,n}+\epsilon_{\log,d,n}^{1-2/q}).
\]
Applying \eqref{eq:matching} to this posterior coupling, define
\[
r_{d,n}=\frac1p\E\min_{g\in\mathcal G_d}\norm{W-gV}^2,
\qquad r_*(u)=\limsup_d r_{d,\lfloor ud^2\rfloor}.
\]
Taking $q=2/\eta$, $0<\eta<1$, \eqref{eq:coarse} gives
\begin{equation}\label{eq:rate}
r_*(u)\le C_\eta\left(\frac{\log u}{u}\right)^{1-\eta}.
\end{equation}
The moment order is fixed and does not grow with the dimension.

\section{Sharp predictive bound with a noisy parameter oracle}
\paragraph{A reference point coupled through the Gaussian prior.}
Keep $V_t=\sqrt{1-t}W+\sqrt tU$, $0<t\le1/2$. Equivalently,
$W=\sqrt{1-t}V_t+\sqrt t\xi$ with $\xi$ independent of $V_t$. Let
\[
a(x)=\mathsf g(V_t,x),\quad A=\E_Xa(X)a(X)^T,
\quad\lambda=z/p,\quad K=A+\lambda I.
\]
where the hidden-weight gradient is evaluated at a given point,
\[
\mathsf g(A,x):=\left.\nabla_{\widetilde W}
F_{\theta(\widetilde W)}^{(L)}(x)\right|_{\widetilde W=A}.
\]
In the experiment where $V_t$ is known, use the population ridge tangent predictor
\begin{equation}\label{eq:estimator}
\widehat f_j(x)=F_{\theta_t}^{(L)}(x)+a(x)^TK^{-1}
\frac1j\sum_{i=1}^j a(X_i)(Y_i-F_{\theta_t}^{(L)}(X_i)).
\end{equation}
This predictor depends only on the known $V_t$, the training data and the input law. It only needs to be a valid predictor (the population matrix need not be computed); its risk bounds the Bayes risk.

\subsection{Exact bias and noise terms}
\paragraph{Conditionally on $(W,V_t)$ the samples remain independent.}
Write $g=F_\theta^{(L)}-F_{\theta_t}^{(L)}$, $h=\E_X[a(X)g(X)]$ and let
\[
(Pg)(x)=a(x)^TK^{-1}h,\qquad
\ell_z(x)=a(x)^TK^{-1}AK^{-1}a(x).
\]
Computing the covariance of the sample mean and averaging over an independent test input gives
\begin{equation}\label{eq:ridge-exact}\begin{aligned}
\E[\norm{\widehat f_j-F_\theta^{(L)}}_{L_X^2}^2\mid W,V_t]
={}&\norm{(I-P)g}_{L_X^2}^2
+\frac\Delta j\tr[(AK^{-1})^2]\\
&+\frac1j\{\E_X[\ell_z(X)g(X)^2]
-h^TK^{-1}AK^{-1}h\}.
\end{aligned}\end{equation}
The negative part of the last term may be dropped. Each eigenvalue in the noise trace is $(s/(s+\lambda))^2\le1$, so
\[
\frac\Delta j\tr[(AK^{-1})^2]\le\frac{\Delta p}{j}.
\]
For a Gaussian plug-in predictor the KL equals the mean squared error divided by $2\Delta$. The Bayes predictive minimizes the average log-loss, so the noise contributes with coefficient $p/(2j)$, not merely $Cp/j$ with an unknown constant.

\subsection{The two remainders}
\paragraph{Ridge regularization controls both remainders.}
$P$ is a positive self-adjoint contraction on $L_X^2$. By \eqref{eq:ou-rem},
$g=\sqrt t\,a^T\xi+E_t$. Conditionally on $V_t$, $\xi$ is isotropic, so
\[
\E_\xi\norm{(I-P)(\sqrt t\,a^T\xi)}^2
=t\tr[A\lambda^2(A+\lambda I)^{-2}]
\le tp\lambda/4=tz/4.
\]
Adding the contraction bound for $E_t$,
\[
\E\norm{(I-P)g}^2\le C(t^{1+\delta}+tz).
\]
Also, by $s/(s+\lambda)^2\le1/(4\lambda)$,
$\ell_z(x)\le p\norm{a(x)}^2/(4z)$.
Write $\theta_{\rm rot}(s)=\theta(W_{\rm rot}(s))$. Along the Gaussian rotation path $W_{\rm rot}(s)=\cos s\,V_t+\sin s\,Z$ the velocity is a standard Gaussian independent of $W_{\rm rot}(s)$. If $\sin^2\omega=t$, Jensen and the Gaussian fourth moment give
\[
\E|g(X)|^4
\le\omega^3\int_0^\omega3\E\norm{\mathsf g(W_{\rm rot}(s),X)}^4\,ds
\le Ct^2.
\]
Hence Cauchy--Schwarz gives
$\E[\norm a^2g^2]\le Ct$, and the signal sampling term is at most $Cpt/(jz)$.
Substituting into \eqref{eq:ridge-exact}, the average predictive KL given the auxiliary observation $V_t$, denoted $\epsilon^V_{\log,d,j}$, satisfies
\begin{equation}\label{eq:oracle}
\epsilon^V_{\log,d,j}\le\frac p{2j}
+C(t^{1+\delta}+tz)+\frac{Cpt}{jz}.
\end{equation}
The estimate is an unconditional average over the true prior OU pair; no Taylor expansion at posterior points is used.

\section{Information cost of the oracle and the limit}
\paragraph{The information cost of the auxiliary observation.}
Define $J_j(t)=I(W;V_t\mid \mathcal D_j)$. Given the data, draw an independent posterior sample $W'$ and form, with the finite matching family, the density
\[
Q_{W'}=\frac1{|\mathcal G_d|}\sum_{g\in\mathcal G_d}
N(\sqrt{1-t}\,gW',tI).
\]
Every mixture component is dominated by $|\mathcal G_d|Q_{W'}$, so
\begin{equation}\begin{aligned}
\KL(N(\sqrt{1-t}W,tI)\|Q_{W'})
\le\log|\mathcal G_d|+\frac{1-t}{2t}\min_g\norm{W-gW'}^2.
\end{aligned}\end{equation}
The relative-entropy decomposition splits this average KL into the conditional mutual information and a nonnegative reference error; since $W'$ is conditionally independent,
\begin{equation}\label{eq:oracle-cost}
J_j(t)\le\log|\mathcal G_d|+\frac{p\,r_{d,j}}{2t}.
\end{equation}
The optimal $g$ is used only to construct the bound; the cost of the auxiliary information is accounted for by this information bound.

\subsection{Exact chain identity over a sample block}
\paragraph{Block differences cancel the unobservable information.}
Since $V_t-W-\mathcal D_n$ is a Markov chain,
\[
I(W;\mathcal D_n\mid V_t)=I(W;\mathcal D_n)-I(W;V_t)+J_n(t).
\]
Subtracting gives, for any $j_0<n$, the exact identity
\begin{equation}\label{eq:remove}
\sum_{j=j_0}^{n-1}\epsilon_{\log,d,j}
=\sum_{j=j_0}^{n-1}\epsilon^V_{\log,d,j}+J_{j_0}(t)-J_n(t).
\end{equation}
If $T=F_\theta^{(L)}$ is the random function, then given $T$ the data $\mathcal D_j$ are independent of the parameter fibre, so
$J_j(t)=I(T;V_t\mid \mathcal D_j)+I(W;V_t\mid T)$.
The last term cancels in \eqref{eq:remove}.

Monotonicity of the risk, $J_n\ge0$, \eqref{eq:oracle} and \eqref{eq:oracle-cost} give the finite-dimensional inequality
\begin{equation}\label{eq:block}\begin{aligned}
(n-j_0)\epsilon_{\log,d,n}\le{}&\log|\mathcal G_d|+\frac{p\,r_{d,j_0}}{2t}
+\frac p2\sum_{j=j_0}^{n-1}\frac1j\\
&+C(n-j_0)(t^{1+\delta}+tz)
+\frac{Cpt}{z}\sum_{j=j_0}^{n-1}\frac1j.
\end{aligned}\end{equation}

\subsection{Taking limits}
\paragraph{All small parameters are fixed during the width limit.}
Fix $\vartheta\in(0,1)$ and set
$n=\lfloor\alpha d^2\rfloor$, $j_0=\lfloor(1-\vartheta)\alpha d^2\rfloor$.
Then $(n-j_0)/d^2\to\vartheta\alpha$, the harmonic sum tends to
$\log(1/(1-\vartheta))$, and the finite matching entropy divided by $d^2$ vanishes. \eqref{eq:block} gives
\begin{equation}\label{eq:limit-upper}\begin{aligned}
\alpha\limsup_d\epsilon_{\log,d,n}\le{}&
\frac c{2\vartheta}\frac{r_*((1-\vartheta)\alpha)}t
+\frac c{2\vartheta}\log\frac1{1-\vartheta}\\
&+C\alpha(t^{1+\delta}+tz)
+\frac{Cc}{\vartheta}\frac tz\log\frac1{1-\vartheta}.
\end{aligned}\end{equation}
Take $\eta=1/8$, $t=\alpha^{-3/4}$, $z=\alpha^{-1/2}$. Even with the weaker piecewise-linear exponent $\delta=1/2$ all errors vanish:
\begin{align*}
r_*/t&=O_\vartheta(\alpha^{-1/8}(\log\alpha)^{7/8}),\\
\alpha t^{3/2}&=\alpha^{-1/8},\qquad
\alpha tz=\alpha^{-1/4},\qquad t/z=\alpha^{-1/4}.
\end{align*}
Letting first $\alpha\to\infty$ and then $\vartheta\downarrow0$, since
$\vartheta^{-1}\log(1/(1-\vartheta))\to1$, we obtain
\[
\limsup_{\alpha\to\infty}\alpha\limsup_d\epsilon_{\log,d,n}\le c/2.
\]
Combining with \eqref{eq:lower-final} and the identity
\begin{equation}\begin{aligned}
\limsup_d|\alpha \epsilon_{\log,d,n}-c/2|
=\max\{&\alpha\limsup_d\epsilon_{\log,d,n}-c/2,\\
&c/2-\alpha\liminf_d\epsilon_{\log,d,n}\},
\end{aligned}\end{equation}
This proves Theorem~\ref{thm:main}. The argument keeps the prescribed order of the width and sample-ratio limits and never differentiates the leading term of the cumulative information.


\appendix
\section{Analytic estimates for fixed-depth networks and domains of integration by parts}\label{an:appendix}\label{an:section}\label{app:local}
\paragraph{At fixed depth and proportional width, moments and local derivatives are controlled uniformly in the dimension.}
All expectations are over independent standard Gaussian inputs and weights; posterior replica pairs are never treated as Gaussian couplings here. Fix the depth $L$, let $k_0=d$, and $cd\le k_\ell\le Cd$ with fixed $0<c<C<\infty$. We use
\[
 u^{(\ell)}=W^{(\ell)} h^{(\ell-1)}/\sqrt{k_{\ell-1}},\quad
 h^{(\ell)}=\sigma_\ell(u^{(\ell)}),\quad h^{(0)}=X,\qquad
 F=k_L^{-1/2}\boldsymbol1^Th^{(L)}.
\]
Each activation is fixed, globally Lipschitz, with $\mathbb E\sigma_\ell(G)=0$, $\mathbb E\sigma_\ell(G)^2=1$. Constants may depend on $L,c,C$, the activations and the specified finite moment orders, but not on $d$. The smoothing constants below are also independent of $0<\dsm\le\delta_{\mathrm{sm},*}$, and so are the dimension thresholds for $\dsm$ in this range. Every moment order is fixed before the dimension threshold is chosen.

\subsection{Moments of the output and hidden norms}
\begin{lemma}\label{an:moments}
Let $q_\ell=\|h^{(\ell)}\|^2/k_\ell$. For every fixed $r>0$,
\[
 \mathbb E q_\ell^r\le C_r,\qquad
 \mathbb E|q_\ell-1|^r\le C_r d^{-r/2},\qquad
 \mathbb E|F|^r\le C_r.
\]
These hold for networks of fixed depth mixing the activations above. Moreover, for every fixed $r>0$ one has $\sup_{\ell,i,d}\mathbb E|h_i^{(\ell)}|^r<\infty$; products of fixed powers of finitely many features have uniformly bounded absolute expectations, even when the features come from different layers or input copies, provided each copy has the prescribed Gaussian marginal.
\end{lemma}
\begin{proof}
Write $b(q)=\mathbb E\sigma(\sqrt qG)$ and $v(q)=\mathbb E\sigma(\sqrt qG)^2$. By $|\sigma(z)|\le M|z|+M_0$, for all $q\ge0$,
\[
 \mathbb E|\sigma(\sqrt qG)|^s\le C_s(1+q^{s/2}),\quad
 |b(q)|\le C|q-1|,\quad |v(q)-1|\le C|q-1|. \tag{A1}\label{an:scalar}
\]
The last bound follows from the difference of squares and
$|\sqrt q-1|(\sqrt q+1)=|q-1|$; the middle one uses $b(1)=0$. These are global bounds; no bad-variance event is removed.

First record a direct moment expansion for sample means. If conditionally independent, conditionally centered $Y_i$ satisfy
$\mathbb E[|Y_i|^s\mid\mathcal F]\le C_s(1+q^{as})$, $2\le s\le2m$, then
\[
 \mathbb E\left[\left|N^{-1}\sum_iY_i\right|^{2m}\middle|\mathcal F\right]
 \le C_mN^{-m}(1+q^{2ma}). \tag{A2}\label{an:sample}
\]
Expanding the $2m$-th power, terms containing an index exactly once vanish. The remaining terms have at most $m$ distinct indices; there are at most $C_mN^j$ terms with $j$ distinct indices, each with conditional moment at most $C_m(1+q^{2ma})$ in absolute value. Dividing by $N^{2m}$ gives the bound.

Given the previous layer, the current preactivations are independent $N(0,q_{\ell-1})$. Jensen first gives
$\mathbb E[q_\ell^{2m}\mid\mathcal F_{\ell-1}]\le C_m(1+q_{\ell-1}^{2m})$, so induction over the fixed depth gives the positive moment bounds. For
$Y_i=(h_i^{(\ell)})^2-v_\ell(q_{\ell-1})$
apply \eqref{an:sample} and then \eqref{an:scalar} to get
\[
 \mathbb E|q_\ell-1|^{2m}\le C_md^{-m}
       +C_m\mathbb E|q_{\ell-1}-1|^{2m}.
\]
The initial step follows from $q_0=d^{-1}\sum_iX_i^2$ and the same expansion, giving $O(d^{-m})$; hence the bound holds for all hidden layers.
Finally, conditionally on the first $L-1$ layers,
\[
 F=\sqrt{k_L}\,b_L(q_{L-1})+
 k_L^{-1/2}\sum_i\{h_i^{(L)}-b_L(q_{L-1})\}.
\]
The conditional moment of order $2m$ of the second term is controlled by \eqref{an:sample} as $C_m(1+q_{L-1}^m)$; the moment of order $2m$ of the first term is at most
$C_mk_L^m\mathbb E|q_{L-1}-1|^{2m}=O(1)$.
Taking $2m\ge r$ and using monotonicity of moments gives the output moments. Conditioning on the previous layers and using \eqref{an:scalar} gives single-feature moments; moments of fixed products follow from H\"older's inequality, without independence of the factors. This decomposition keeps the conditional mean of the all-ones readout.
\end{proof}

\subsection{Dimension-free bounds for forward and backward passes and smooth Hessians}
Let $A_\ell=W^{(\ell)}/\sqrt{k_{\ell-1}}$, $B_\ell=\|A_\ell\|_{\rm op}$. Below, $P$ denotes a nonnegative polynomial of fixed degree in $1+\|X\|/\sqrt d$ and finitely many $B_\ell$; normalized operator norms of a second independent set of Gaussian matrices are also allowed. Every fixed-order moment of every such $P$ is uniformly bounded. To check this, take $1/4$-nets of at most $9^m,9^r$ points on the unit spheres of $\mathbb R^m,\mathbb R^r$. For a Gaussian $m\times r$ matrix $G$, the net approximation of bilinear forms gives
\[
 \|G\|_{\rm op}\le2\max_{u,v\text{ in the net}}|u^TGv|,
 \quad
 \mathbb P\{\|G\|_{\rm op}^2\ge8[\log2+(m+r)\log9+s]\}\le e^{-s}.
\]
Integrating this tail bound and using bounded width ratios gives all required moments.

\begin{lemma}\label{an:derivatives}
For Lipschitz activations the weak gradient satisfies $\mathbb E\|\nabla_WF\|^r\le C_r$ for every fixed $r>0$. If moreover every activation is $C^2$ with bounded second derivative, then
\[
 \mathbb E\|\nabla_W^2F\|_{\rm HS}^2\le C.
\]
In particular normalized tanh and centered normalized sigmoid satisfy these bounds.
\end{lemma}
\begin{proof}
The forward recursion gives $\|h^{(\ell)}\|/\sqrt{k_\ell}\le P$. Define
\[
 \beta_L=k_L^{-1/2}\boldsymbol1,\quad
 \lambda_\ell=\operatorname{diag}(\sigma_\ell'(u^{(\ell)}))\beta_\ell,\quad
 \beta_{\ell-1}=A_\ell^T\lambda_\ell.
\]
By bounded slopes, $\|\beta_\ell\|+\|\lambda_\ell\|\le P$, and
\[
 \nabla_{W^{(\ell)}}F=\lambda_\ell(h^{(\ell-1)})^T/\sqrt{k_{\ell-1}},
 \qquad \|\nabla_WF\|^2\le P. \tag{A3}\label{an:gradient}
\]
The network is locally Lipschitz on each finite-dimensional space; the chain rule holds at points of differentiability and represents the weak derivative. Alternatively, smooth the activations layer by layer and identify the same weak gradient by the bounds above and weak compactness.

Let $J_{u,\ell},J_{h,\ell}$ be the Jacobians with respect to all weights. From
\[
 \|J_{u,\ell}\|_{\rm op}\le B_\ell\|J_{h,\ell-1}\|_{\rm op}
                  +\|h^{(\ell-1)}\|/\sqrt{k_{\ell-1}},\qquad
 \|J_{h,\ell}\|_{\rm op}\le M\|J_{u,\ell}\|_{\rm op},
\]
both operator norms are at most $P$. The second-order chain rule on the computation graph writes the output Hessian exactly as a sum over layers of two kinds of terms. The activation term is
\[
 J_{u,\ell}^T\operatorname{diag}(\beta_\ell\odot\sigma_\ell''(u^{(\ell)}))J_{u,\ell};
\]
whose Hilbert--Schmidt norm is at most
$\|J_{u,\ell}\|_{\rm op}^2\|\sigma_\ell''\|_\infty\|\beta_\ell\|$.
The bilinear form of the mixed term is
\[
 k_{\ell-1}^{-1/2}\lambda_\ell^T
 \{U^{(\ell)}Dh^{(\ell-1)}[V]+V^{(\ell)}Dh^{(\ell-1)}[U]\}.
\]
For one ordering of directions its squared Hilbert--Schmidt norm is exactly
$\|\lambda_\ell\|^2\|J_{h,\ell-1}\|_{\rm HS}^2/k_{\ell-1}$,
which is at most $\|\lambda_\ell\|^2\|J_{h,\ell-1}\|_{\rm op}^2$ because the Jacobian has rank at most $k_{\ell-1}$. The transposed ordering is identical. Summing over the fixed number of layers and using the moment bounds for $P$ concludes.
\end{proof}
By the two lemmas and Cauchy--Schwarz, for every fixed $Q>0$ also
\[
 \mathbb E[(1+|F|^Q)\|\nabla F\|^2]\le C_Q,
 \qquad
 \mathbb E[\|\nabla F\|^2\boldsymbol1_{\{|F|>M\}}]\le C_QM^{-Q}. \tag{A4}\label{an:weightedmoments}
\]

\subsection{The centered piecewise-linear branch: negative moments and weighted gate events}
Consider, for fixed real $a,b$, the activation
\[
 \sigma(z)=az+b(|z|-c_0),\qquad c_0=\sqrt{2/\pi},
 \qquad a^2+b^2(1-2/\pi)=1. \tag{A5}\label{an:pl}
\]
If $b=0$, then $\sigma(z)=\pm z$ is linear and does not belong to $\mathcal A_{\rm bal}$: it is unbounded and $a_m=0$ for all $m\ge2$. The network then depends only on products of the weight matrices, its effective dimension is $O(d)$, and the conclusion of the main theorem fails. Below $b\ne0$; neither $a\ne0$ nor nonvanishing of both $a+b$ and $a-b$ is required, so the even case and ReLU-type activations with one zero slope are included. Constants may depend on the fixed $(a,b)$; uniformity as $b\to0$ is not claimed.

\begin{lemma}\label{an:negative}
There are $\eta_0,C_0>0$ such that for all $r\ge0$, $0<\eta\le\eta_0$,
\[
 \mathbb P\{\sigma(rG)^2\le\eta\}\le C_0\sqrt\eta.
\]
If $Q=N^{-1}\sum_{i=1}^N\sigma(rG_i)^2$, then for every fixed $p>0$ there is $N_0(p)$ with
$\sup_{r\ge0,N\ge N_0(p)}\mathbb E Q^{-p}<\infty$.
\end{lemma}
\begin{proof}
Write $c=bc_0\ne0$. On the two half-lines the activation equals $s_\pm z-c$ with $s_\pm=a\pm b$. If a slope vanishes, that half-line takes the constant value $-c$, which does not contribute to the small-ball event when $\sqrt\eta<|c|/2$. On a half-line with nonzero slope the small-ball event lies in an interval centred at the nonzero root $c/s_\pm$ of length $2\sqrt\eta/|s_\pm|$; half-lines without an admissible root are discarded. After shrinking $\eta_0$, each interval has a fixed positive distance from zero. For $|z|\ge\delta>0$ the density $r^{-1}\phi(z/r)$ is bounded uniformly in $r>0$, since with $v=|z|/r$ it equals $v\phi(v)/|z|$. For $r=0$ the event is empty. This gives the small-ball bound.

The event $Q\le\eta$ requires at least $N/2$ terms to be at most $2\eta$; conditional independence and subset counting give
$\mathbb P(Q\le\eta)\le(C_1\eta^{1/4})^N$.
Choosing $\eta_0$ smaller so that $C_1\eta_0^{1/4}<1/2$,
\[
 \mathbb E Q^{-p}\le\eta_0^{-p}
       +pC_1^N\int_0^{\eta_0}s^{N/4-p-1}\,ds
 \le\eta_0^{-p}+\frac{p\eta_0^{-p}2^{-N}}{N/4-p}
\]
uniformly bounded for $N>4p$. Conditioning layer by layer gives fixed-order negative moments for all hidden norms and for the mean squared norms with one coordinate deleted. Negative moments of the input norm follow from
\[
 \mathbb E q_0^{-p}=\prod_{j=1}^p\frac d{d-2j}\le2^p\quad(d\ge4p)
\]
first for integer $p$, then for other orders by monotonicity.
\end{proof}

The dependence between backpropagated quantities and forward representations is handled through conditional covariances. Fix a layer $\ell<L$ and write $P_>$ for the downstream norm polynomial. Conditionally on all variables up to this layer (possibly together with a coupled representation $h'$), the downstream weighted covariance of the current network is
\[
 M(h)=\mathbb E[P_>\beta_\ell\beta_\ell^T\mid h^{(\ell)}=h].
\]
The integrand uses only the current network and downstream Gaussian norms, not $h'$. Right-multiplying the first downstream matrix by any orthogonal matrix preserves its Gaussian law and its norms; a coupled Gaussian noise matrix, if present, is rotated simultaneously. Hence $M$ is orthogonally equivariant, and the stabilizer of a fixed $h$ gives
\[
 M(h)=c_\perp(I-hh^T/\|h\|^2)+c_\parallel hh^T/\|h\|^2.
\]
The backward norm bound does not depend on the input $h$, so $\operatorname{tr}M(h)\le C$ and, for $k=k_\ell\ge2$,
\[
 M(h)_{ii}\le C\left(k^{-1}+\frac{h_i^2}{\|h\|^2}\right). \tag{A6}\label{an:covariance}
\]
For the last layer use $\beta_{L,i}^2=1/k_L$ directly. For any $P$ in the forward bounds, split each monomial into previous-layer, current-layer and downstream norm factors and apply this bound.

For the current-layer dependence, delete row $i$ and write $W^{(-i)}$; from
\[
 \|W/\sqrt{k_{\ell-1}}\|_{\rm op}
 \le\|W^{(-i)}/\sqrt{k_{\ell-1}}\|_{\rm op}
       +\|W_i\|/\sqrt{k_{\ell-1}},\qquad
 \|h\|^{-2}\le\left(\sum_{j\ne i}h_j^2\right)^{-1},
\]
the current norm polynomial is bounded by a product of a polynomial in the row-deleted variables and a polynomial in the normalized row length. Conditionally on the previous layers, the row-deleted variables are independent of row $i$. The previous lemma and H\"older's inequality give
\[
 \mathbb E\left[\frac{P_{\rm del}}{\sum_{j\ne i}h_j^2}
                  \middle|\text{previous layers}\right]\le C/k. \tag{A7}\label{an:rowdelete}
\]
This uses fixed-order negative moments and operator-norm moments, not a replacement of the random denominator by its mean.

Given the previous layer, the two preactivations of any pair of Gaussian rows can be written as
$u=\sqrt qG$, $v=\sqrt r(\cos\omega G+\sin\omega H)$.
The event of different signs is a pair of sectors of total angle $2\omega$ in the $(G,H)$ plane. Polar integration, moments of the remaining independent Gaussian coordinates, and $h_i^2\le C(1+qG^2)$ give, for any fixed polynomial $Q$ in the normalized row lengths,
\[
 \mathbb E[Q\boldsymbol1_{uv<0}\mid\text{previous layers}]\le C_Q\omega,
 \quad
 \mathbb E[Qh_i^2\boldsymbol1_{uv<0}\mid\text{previous layers}]\le C_Q(1+q)\omega. \tag{A8}\label{an:angular}
\]
Combining \eqref{an:covariance}--\eqref{an:angular} and summing over $i$ gives the weighted gate-event bound; outside the angle only a previous-layer norm polynomial remains. At no step is $\beta_{\ell,i}$ assumed independent of the current gates.

\subsection{Gradient error and Hessian energy of the smoothed network}
Take a fixed even $C^\infty$ probability density $\rho$ supported in $[-1,1]$. Define
\[
 \widetilde\sigma_{\dsm}=\rho_{\dsm}*\sigma,
 \quad m_{\dsm}=\mathbb E\widetilde\sigma_{\dsm}(G),
 \quad s_{\dsm}^2=\operatorname{Var}(\widetilde\sigma_{\dsm}(G)),
 \quad \sigma_{\dsm}=(\widetilde\sigma_{\dsm}-m_{\dsm})/s_{\dsm}.
\]
By Lipschitz continuity, $|m_{\dsm}|+|s_{\dsm}-1|\le C\dsm$. Choose a fixed $\delta_{\mathrm{sm},*}>0$ with $s_{\dsm}\ge1/2$; then
\[
 \begin{gathered}
 |\sigma_{\dsm}(z)-\sigma(z)|\le C\dsm(1+|z|),\qquad
 |\sigma_{\dsm}'(z)-\sigma'(z)|\le C\boldsymbol1_{|z|\le\dsm}+C\dsm,\\
 |\sigma_{\dsm}''(z)|\le C\dsm^{-1}\boldsymbol1_{|z|\le\dsm}.
 \end{gathered} \tag{A9}\label{an:scalar-smooth}
\]
All first derivatives are uniformly bounded.

\begin{lemma}\label{an:smoothing}
Build $F_{\dsm}$ with the same weights. There are $d_0,C,\delta_{\mathrm{sm},*}>0$ such that for all $d\ge d_0$ and all $0<\dsm\le\delta_{\mathrm{sm},*}$,
\[
 \mathbb E\|\nabla F_{\dsm}\|^4\le C,\qquad
 \mathbb E\|\nabla F-\nabla F_{\dsm}\|^2\le C\dsm,
 \qquad \mathbb E\|\nabla^2F_{\dsm}\|_{\rm HS}^2\le C/\dsm.
\]
This covers the whole $b\ne0$ branch of \eqref{an:pl}.
\end{lemma}
\begin{proof}
Forward induction and \eqref{an:scalar-smooth} give
$\|h_{\dsm}^{(\ell)}-h^{(\ell)}\|/\sqrt{k_\ell}\le\dsm P$, while the forward and backward bounds are independent of $\dsm$. Hence the fourth gradient moment follows from \eqref{an:gradient}.

First, the negative moments hold uniformly for the smoothed network. Near zero the original activation stays a fixed distance away from zero. For $\delta_{\mathrm{sm},*}$ small enough, the smoothed renormalized activation stays away from zero on this fixed neighbourhood. Outside it the convolution equals the original affine piece (the kernel is even), up to an $O(\dsm)$ shift and renormalization; nonzero slopes stay away from zero, and the constant of a zero-slope piece stays away from zero. Hence the small-ball preimages remain finitely many intervals of length $O(\sqrt\eta)$ away from the origin, with uniform constants. The proof of Lemma~\ref{an:negative} then gives all required negative moments of the smoothed features and their row-deleted norms.

Preactivations of the same layer in the original and smoothed networks form a conditionally Gaussian pair. If their previous-layer representations are $h,h_{\dsm}$, the angle satisfies
\[
 \omega\le C\frac{\|h-h_{\dsm}\|}{\sqrt{\|h\|\|h_{\dsm}\|}}
 \le C\dsm P(q q_{\dsm})^{-1/4}.
\]
using $1-\langle h,h_{\dsm}\rangle/(\|h\|\|h_{\dsm}\|) \le\|h-h_{\dsm}\|^2/(2\|h\|\|h_{\dsm}\|)$. Hence, for any fixed previous-layer polynomial $P_0$, $\mathbb E P_0\omega\le C\dsm$. In the first layer the two preactivations coincide and the angle is zero.

By \eqref{an:scalar-smooth}, the squared slope difference is at most
\[
 C\boldsymbol1_{uv<0}+C\boldsymbol1_{|v|\le\dsm}+C\dsm^2.
\]
The first event is controlled by \eqref{an:angular}. For the second, write $v=\sqrt rV$; the standardized original preactivation is $\tau V+\sqrt{1-\tau^2}H$. Gaussian integration gives
\[
 \mathbb E[Q\boldsymbol1_{|v|\le\dsm}]\le C_Q\dsm/\sqrt r,
 \quad
 \mathbb E[Qh_i^2\boldsymbol1_{|v|\le\dsm}]\le C_Q(1+q)\dsm/\sqrt r. \tag{A10}\label{an:boundary}
\]
If $\dsm/\sqrt r\ge1$, use Gaussian moments directly; otherwise integrate the density over $|V|\le\dsm/\sqrt r\le1$ and then over $H$ and the remaining coordinates. Row-length polynomials do not change the interval-length factor. The same holds for the smoothed features.

Concretely, let $\mathcal F_{\ell-1}^{\rm pair}$ contain the previous-layer variables of both networks,
$q=\|h^{(\ell-1)}\|^2/k_{\ell-1}$, $r=\|h_{\dsm}^{(\ell-1)}\|^2/k_{\ell-1}$.
First condition on all rows of the current layer, integrate out the downstream part and use \eqref{an:covariance} for the original network;
then replace the denominator by $\sum_{j\ne i}h_j^2$ and separate the row-deleted variables from row $i$ via \eqref{an:rowdelete}.
Integrating row $i$ with \eqref{an:angular} and \eqref{an:boundary} gives the explicit conditional bound
\[
 \mathbb E\left[P\sum_{i=1}^{k_\ell}\beta_{\ell,i}^2
 |\sigma'(u_i^{(\ell)})-\sigma_{\dsm}'(u_{\dsm,i}^{(\ell)})|^2
 \middle|\mathcal F_{\ell-1}^{\rm pair}\right]
 \le C P_{<}(1+q)\left(\omega+\frac{\dsm}{\sqrt r}+\dsm^2\right).
 \tag{A10a}\label{an:conditionalgate}
\]
Here $P_{<}$ is another fixed-degree previous-layer norm polynomial generated by $P$.
Each row bound carries $1/k_\ell$, which cancels when summing over rows.
The last layer gives the same bound with the deterministic $\beta_{L,i}^2=1/k_L$.
Using the angle bound, the negative moments of the input and hidden norms, and taking expectations over the previous layers,
\[
 \mathbb E\left[P\sum_i\beta_{\ell,i}^2
 |\sigma'(u_i^{(\ell)})-\sigma_{\dsm}'(u_{\dsm,i}^{(\ell)})|^2\right]
 \le C_P\dsm. \tag{A11}\label{an:weightedgate}
\]
Here $\beta$ is the single downstream gradient of the original network; its covariance depends only on the original $h$, and no claim is made that the cross-covariance of the two networks has the same form. Decompose along the backward recursion
\[
 \lambda_{\dsm,\ell}-\lambda_\ell
 =D_{\dsm,\ell}(\beta_{\dsm,\ell}-\beta_\ell)
   +(D_{\dsm,\ell}-D_\ell)\beta_\ell.
\]
Inducting from the common readout, squaring the fixed finite sum and using \eqref{an:weightedgate}, each backward difference has squared expectation $O(\dsm)$, also after multiplication by any required $P$. The gradient block formula and the forward difference $\dsm P$ give the gradient error bound.

Finally use the exact Hessian decomposition of Lemma~\ref{an:derivatives}. The squared energy of the mixed terms is uniformly bounded. The squared Hilbert--Schmidt norm of the activation term is at most
\[
 \frac P{\dsm^2}\sum_i\beta_{\dsm,\ell,i}^2
                       \boldsymbol1_{|u_{\dsm,i}^{(\ell)}|\le\dsm}.
\]
Applying \eqref{an:covariance} for the smoothed network, integrating over the boundary layer with \eqref{an:boundary}, and controlling denominators by row-deleted negative moments, its expectation is $O(\dsm^{-1})$. The first-layer $q_0^{-p}$ follows from the explicit chi-square negative-moment formula. Summing over the fixed number of layers concludes. Only finitely many fixed-order moments are used, so $d_0$ can be chosen once, independently of $\dsm$.
\end{proof}
This proves approximation of derivatives; it does not assume $\mathbb E|F-F_{\dsm}|^2=O(\dsm)$. The smoothed mean of the all-ones readout may accumulate; the test field smooths only gradients and never replaces the output in the true channel.

\subsection{Tangent remainder of the prior OU coupling}
\begin{lemma}\label{an:ou}
Let $W(t)=\sqrt{1-t}\,W+\sqrt tZ$, where $Z$ is an independent standard Gaussian vector of all parameters and $0<t\le1/2$. Then
\[
 \mathbb E|F(W(t),X)-F(W,X)-\sqrt t\,DF(W,X)[Z]|^2
 \le\begin{cases}Ct^2,&\text{smooth branch},\\Ct^{3/2},&\text{piecewise-linear branch}.
 \end{cases}
\]
That is, $O(t^{1+\delta})$ holds with $\delta=1$ and $\delta=1/2$ respectively. Constants and dimension thresholds do not depend on $t$.
\end{lemma}
\begin{proof}
Condition on $X=x$, expand $F$ in the orthonormal Hermite chaos, let $c_N(x)$ be the energy of order $N$, and write $a_t=\sqrt{1-t}$. The generating function gives the Mehler identity
$\mathbb E H_\nu(W)H_\mu(W(t))=\boldsymbol1_{\nu=\mu}a_t^{|\nu|}$,
and $\partial_iH_\nu=\sqrt{\nu_i}H_{\nu-e_i}$. Hence
\[
 \mathbb E\|\nabla F(W(t))-\nabla F(W)\|^2
   =2\sum_{N\ge1}N(1-a_t^{N-1})c_N,
\]
\[
 \mathbb E|R_t|^2
  =\sum_{N\ge1}\{2(1-a_t^N)+tN-2tNa_t^{N-1}\}c_N.
\]
The cross term in the second identity can also be obtained by Gaussian integration by parts in $Z$. For $N\ge2$, convexity gives
$a_t^N=(1-t)^{N/2}\ge1-Nt/2$; for $N=1$, $2(1-a_t)-t=(1-a_t)^2\le t^2$. Hence
\[
 \mathbb E|R_t|^2\le t^2\mathbb E\|\nabla F\|^2
             +t\mathbb E\|\nabla F(W(t))-\nabla F(W)\|^2. \tag{A12}\label{an:hermite}
\]
These identities hold first for finite Hermite sums and then by truncation in Gaussian $W^{1,2}$; Lemma~\ref{an:derivatives} guarantees the domain.
For the smooth branch $\sum_NN(N-1)c_N=\mathbb E\|\nabla^2F\|_{\rm HS}^2\le C$, so the squared gradient difference is $O(t)$.
For the piecewise-linear branch, insert $\nabla F_{\dsm}$ for any $\dsm\le\delta_{\mathrm{sm},*}$ and use the common prior marginals of the OU pair and the smoothing lemma:
\[
 \mathbb E\|\nabla F(W(t))-\nabla F(W)\|^2
 \le C\dsm+Ct/\dsm.
\]
Taking $\dsm=\sqrt t$ (if $\sqrt t>\delta_{\mathrm{sm},*}$, use the second gradient moment directly and enlarge the constant) gives $O(\sqrt t)$. Substituting into \eqref{an:hermite} completes the proof.
The full OU remainder is kept; the radial drift cannot be separated and assumed uniformly bounded in the dimension.
\end{proof}

\subsection{Gaussian divergence, weak derivatives and truncation}
\begin{lemma}\label{an:sobolev}
Let $g:\mathbb R^p\to\mathbb R^p$ be in Gaussian $W^{1,2}$, possibly with values depending on independent external $L^2$ variables. Define
$\delta g=\operatorname{div}g-W\cdot g$. Its closed extension satisfies
\[
 \mathbb E(\delta g)^2
 =\mathbb E\|g\|^2+\mathbb E\sum_{i,j}(\partial_i g_j)(\partial_jg_i)
 \le\mathbb E\{\|g\|^2+\|Dg\|_{\rm HS}^2\}. \tag{A13}\label{an:divergence}
\]
\end{lemma}
\begin{proof}
First take compactly supported smooth fields. Integrating by parts twice coordinatewise against the Gaussian density and using
$\partial_i(W_jg_j)=\boldsymbol1_{i=j}g_j+W_j\partial_i g_j$, the cross terms combine into the identity; $2uv\le u^2+v^2$ applied to each term of the derivative sum gives the inequality. In general, first apply a radial cutoff $\chi_R$ with $0\le\chi_R\le1$, equal to one on the ball, $\|\nabla\chi_R\|_\infty\le C/R$, then convolve on compact sets. The Gaussian density is smooth with positive upper and lower bounds on compact sets, so ordinary convolution approximation also gives Gaussian $W^{1,2}$ approximation. Tails are controlled by square integrability of $g,Dg$, and the cutoff derivative term by $C\|g\|/R$. Hence compactly supported smooth fields converge in this Sobolev norm; the inequality makes their divergences Cauchy in $L^2$, with limit the distributional divergence. The identity passes to the limit. External variables are handled by approximating with simple functions and integrating.
\end{proof}

For the regularized test field let $f$ be the true output, $a=\nabla f$ the true gradient and $a_{\dsm}$ the auxiliary smoothed gradient; in the smooth case take $a_{\dsm}=a$. Let $j(f)$ be the channel Fisher weight, $\mathcal I_{\dsm}=\mathbb E_X[j(f)a_{\dsm} a_{\dsm}^T]$, $M_{\dsm}=(\mathcal I_{\dsm}+\lambda I)^{-1}$, with $\lambda>0$ fixed in this paragraph. At each fixed dimension and fixed $\dsm>0$, local Lipschitz continuity of the network, the moment bounds above and the weak product rule allow differentiating $\mathcal I_{\dsm}$ and exchanging the derivative with the $X$ integral; the derivative contains the $j'(f)a$ term, which cannot be omitted. The weak chain rule for the inverse is
\[
 \partial_iM_{\dsm}=-M_{\dsm}(\partial_i\mathcal I_{\dsm})M_{\dsm}.
\]
Under the channel bounds $j\le J_0$, $j'^2/j\le L_0$ of the main text,
\[
 \|D(\sqrt{j(f)}a_{\dsm})\|_{\rm HS}^2
 \le2J_0\|\nabla^2f_{\dsm}\|_{\rm HS}^2
       +\tfrac12L_0\|a\|^2\|a_{\dsm}\|^2,
\]
and the right side is integrable by this section. The local exchange is justified by truncating to compact sets in $W$ and dominating with fixed input moments; no pointwise global bound on the Hessian in $W$ is needed.

After absorbing the mixture density, the field of the main text is $g_t=\sqrt{q_t}\,b_t=\psi_t c$ with $c=M_{\dsm} a_{\dsm}$. Its weak derivative is exactly
\[
 Dg_t=\psi_t Dc+(\partial_u\psi_t)c a^T.
\]
where $h_t=\partial_u\log q_t$. The weak product rule also produces a half-score term that must be kept:
\[
 \delta g_t=\sqrt{q_t}\,\delta b_t+\tfrac12h_t\sqrt{q_t}\,b_t\cdot a,
 \qquad
 \delta b_t+h_tb_t\cdot a=q_t^{-1/2}\delta g_t+\tfrac12h_tb_t\cdot a.
\]
As long as the scalar mixture integrals and matrix estimates of the main text guarantee
$N_0=\int(1-t)\mathbb E\int\|g_t\|^2<\infty$ and
$N_1=\int(1-t)\mathbb E\int\|Dg_t\|_{\rm HS}^2<\infty$,
Lemma~\ref{an:sobolev} applies directly in the $L^2$ space of the external variables $(X,y,t)$.

Finally use $b_{t,R}=\chi_R b_t$ for the original field. For fixed $0<t<1$, positivity of the mixture density and the product rule for the smooth channel on compact sets make it a local Sobolev field; the true piecewise-linear $f$ is continuous and locally Lipschitz, and the field uses $a_{\dsm}$ to avoid differentiating discontinuous gates. For compactly supported fields integrate by parts against the joint density first. Removing the cutoff produces the error $\nabla\chi_R\cdot b_t$, whose integrated squared norm under the joint mixture measure is at most $C N_0/R^2$ and hence vanishes. The auxiliary divergence term is dominated by \eqref{an:divergence} and $N_0+N_1$; the training-score term by conditional centering and the finite Fisher quadratic form of the main text. The target side and the estimation-error side are dominated by the finite trace bound and the finite squared-error bound, so Cauchy--Schwarz gives the limit. Alternatively restrict to $t\in[\eta,1-\eta]$ first and then let $\eta\downarrow0$; the weighted integral bounds above are exactly the dominating functions for removing the endpoint cutoff. For Gaussian regression and binary channels the fixed-dimensional likelihood is bounded, so Fubini is valid in the compactly supported step.

All of the above concern the Gaussian prior and independent Gaussian inputs; the OU remainder requires the specified independent Gaussian coupling and does not apply to arbitrary posterior replica pairs.

\section{Full-parameter tangent bridge and homogeneous bulk spectrum}\label{tg:section}\label{app:bridge}
\paragraph{Scope and scaling.}
This section concerns only tangent geometry in the Gaussian input space, not the existence of risk limits; the statistical conclusions for the predictive KL are proved in the main text. Fix the depth $L$ and let $k_0=d$, $k_\ell/d\to\gamma_\ell\in(0,\infty)$,
\[
 h^{(0)}=X\sim N(0,I_d),\quad h^{(\ell)}=\sigma(W^{(\ell)} h^{(\ell-1)}/\sqrt{k_{\ell-1}}),\quad
 F_W=k_L^{-1/2}\sum_i h_i^{(L)},\quad p=\sum_{\ell=1}^Lk_{\ell-1}k_\ell.
\]
All weights are independent standard Gaussians. $J_WU=D_WF_W[U]$ takes values in $L^2(N(0,I_d))$, so
$\|J_W\|_{\rm HS}^2=\mathbb E_X\|\nabla_WF_W\|^2$ and $A_W=J_W^*J_W$.
We use the normalized probabilists' Hermite polynomials $\psi_r=\mathrm{He}_r/\sqrt{r!}$, $a_r=\mathbb E\sigma(G)\psi_r(G)$.
The scalar hypothesis is one of the following:
\begin{enumerate}
\item $\sigma$ is bounded, smooth, odd, with unit Gaussian variance and every fixed-order derivative bounded; a fixed odd $m\ge5$ satisfies $a_m\ne0$.
\item $\sigma(z)=[az+b(|z|-\sqrt{2/\pi})]/v$ with $v^2=a^2+b^2(1-2/\pi)$, $b\ne0$; take $m=6$.
\end{enumerate}
In both cases $a_0=a_{m-1}=a_{m+1}=0$. The nonvanishing coefficients of tanh and centered sigmoid are proved in the scalar part; this section does not use them as an unproved spectral hypothesis. Write
$N=m^L$, $S=1+m+\cdots+m^{L-1}$, $B=a_m^S$.
Let $H_m$ be the top-input-chaos tangent operator of the homogeneous network obtained by replacing the activation with $z^m/\sqrt{m!}$. In this section $H_m$ is the parameter derivative $D_WH_m(W)$ of the homogeneous function of the main text.

\subsection{Symmetric tensors, multiplicity frames and collision pruning}
For $T\in\operatorname{Sym}^r\mathbb R^n$ we use the inherited Hilbert tensor norm, and $I_r$ is the Wick map, with
\[
 \mathbb E I_r(T)I_r(U)=r!\langle T,U\rangle,\qquad
 \|C_r(\textstyle\sum_{|\alpha|=r}c_\alpha x^\alpha)\|^2
 =\sum_{|\alpha|=r}\frac{\alpha!}{r!}|c_\alpha|^2.\tag{TG.1}
\]
At fixed degree the tensor norm is equivalent to the plain sum of squared coefficients, with dimension-free constants. All symmetrizations below average over all slots; colored tree tensor spaces are never used.

Let unit vectors $w_1,\ldots,w_k\in\mathbb R^n$ have mutual coherence at most $\delta$. Fix $q$ and for $|r|=q$ define
\[
 P_r=\operatorname{Sym}\bigotimes_jw_j^{\otimes mr_j},\quad
 Z_{r,i}=\operatorname{Sym}\bigotimes_jw_j^{\otimes(mr_j-\mathbf1_{j=i})},\quad
 D_{r,i}(u)=\operatorname{Sym}(Z_{r,i}\otimes u),\quad u\perp w_i,
\]
where $r_i>0$. Let $\mathcal F$ be the direct-sum frame of pure and defect tensors, and $\mathcal F^-$ the frame with one slot deleted. There are positive $c,C$ depending only on $m,q$ such that
\begin{equation}\label{tg:frame}
 cI\preceq\mathcal F^*\mathcal F\preceq CI,\qquad
 cI\preceq(\mathcal F^-)^*\mathcal F^-\preceq CI
\end{equation}
provided $\delta+\sum_{h=1}^qk^h\delta^{mh-2}=o(1)$. The same holds when the row lengths lie in a fixed positive interval.

The proof uses the exact pairing formula
\[
 \left\langle\operatorname{Sym}\bigotimes_{a=1}^Ru_a,
 \operatorname{Sym}\bigotimes_{a=1}^Rv_a\right\rangle
 =\frac1{R!}\sum_{\pi\in\mathfrak S_R}\prod_a\langle u_a,v_{\pi(a)}\rangle.\tag{TG.2}
\]
The pure diagonal is $\prod_j(mr_j)!/R!$. With $n_j=mr_j-\mathbf1_{j=i}$, the defect diagonal on $w_i^\perp$ is
\[
 \frac{\prod_jn_j!}{R!}\left(I+\sum_jn_jw_jw_j^*\right)+O_{\rm op}(\delta).
\]
whose fixed positive lower bound comes from the identity. Blocks with different labels and the same finite support are $O(\delta)$: without mismatch the required integer relation
$m(r-r')=e_i+e_b-e_j-e_a$ forces $r=r'$ when $m\ge5$; the remaining cases with different defect slots are cancelled by $u\perp w_i$. If the target has $h$ new labels, at least $mh$ of its new-direction slots must be matched and at most two are absorbed by the defects, so the block norm is at most $C\delta^{mh-2}$. Each row has at most $Ck^h$ such blocks. The block Schur test gives \eqref{tg:frame}; the deleted frame is identical with smaller losses.

Assume also $\max_{i,a}|w_{ia}|\le\kappa$. Let $C_r$ be the orthogonal projection onto coordinate monomials with at least one repeated coordinate, and $M_a$ the projection onto degree-$mq-1$ monomials containing coordinate $a$. Then
\begin{align}
 \|C_{mq}\mathcal F\|_{\rm op}^2+\|C_{mq-1}\mathcal F^-\|_{\rm op}^2
 &\le C\kappa^2\left(1+\sum_{h=1}^qk^h\delta^{mh-4}\right),\label{tg:collision}\\
 \max_a\|M_a\mathcal F^-\|_{\rm op}^2
 &\le C\kappa^2\left(1+\sum_{h=1}^qk^h\delta^{mh-2}\right).\label{tg:marked}
\end{align}
For the first bound, in the unsymmetrized space let $D_{uv}$ force coordinates $u,v$ to coincide. These are commuting diagonal projections, so $0\preceq C_r\preceq\sum_{u<v}D_{uv}$. The four-vector contractions on two specified slots are of only
\[
 \sum_aw_{ia}w_{ja}w_{la}w_{ta},\quad
 \sum_au_aw_{ia}w_{ja}w_{la},\quad
 \sum_au_aw_{ia}v_aw_{ja}
\]
three kinds, bounded respectively by $\kappa^2,\kappa^2\|u\|,\kappa^2\|u\|\|v\|$. The remaining mismatch is at least $mh-4$, and the block Schur test applies again. No coordinate flatness of the defect vectors is needed. The second bound replaces two repeated slots by one specified coordinate slot, gaining $w_{ia}w_{ja}\le\kappa^2$ and losing at most two new-label slots. The proof controls the whole collision operator; the blocks of $\mathcal F^*C_r\mathcal F$ need not carry the same power of the coherence.

For Gaussian rows normalized by the input width, simultaneously
\[
 \max_i|\|w_i\|-1|+\delta=O(\sqrt{\log d/d}),\qquad
 \kappa^2=O(\log d/d),\qquad k\asymp d.
\]
The complementary events can be given any fixed negative power of probability. The perturbation in \eqref{tg:frame} and the right sides of \eqref{tg:collision}--\eqref{tg:marked} all vanish; the worst collision term, at $m=5,h=1$, is $O(d^{-1/2}(\log d)^{3/2})$.

\subsection{Joint tangent spectral lower bound for the homogeneous network}
\paragraph{Homogeneity gauge directions.}
Let $P_\ell(Y)$ be the downstream polynomial in the layer-$\ell$ features as formal variables, of degree $q_\ell=m^{L-\ell}$; $B_{\ell,i}=\partial_iP_\ell$. Write $w_{\ell,i}=W^{(\ell)}_{i:}/\sqrt{k_{\ell-1}}$. Deleted frames and the chain rule give
\[
 \|C(B_{\ell,i})\|^2\asymp
 \sum_jw_{\ell+1,ji}^2\|C(B_{\ell+1,j})\|^2\asymp k_\ell^{-1}.\tag{TG.3}
\]
The recursion starts at $B_{L,i}=k_L^{-1/2}$ and its last step uses the squared column norms $\asymp k_{\ell+1}/k_\ell$. The frame bounds hold for arbitrary coefficients, so no independence between backpropagated coefficients and row directions is assumed.

Take the linear section $\mathcal S_W=\{U:U^{(\ell)}_{i:}\perp W^{(\ell)}_{i:},\ \ell<L\}$, with no restriction on the top layer. If $R$ is the perturbation from higher layers, the joint perturbation at the current layer is
\[
 R(z^{\odot m}/\sqrt{m!})+
 \frac{m}{\sqrt{m!}}\sum_i B_{\ell,i}(z^{\odot m}/\sqrt{m!})z_i^{m-1}
       \frac{U^{(\ell)}_{i:}\cdot Y}{\sqrt{k_{\ell-1}}}.
\]
The first term lies in the pure frame and the second in the non-radial defect frame. Keeping cross-layer terms and applying \eqref{tg:frame} jointly gives
\[
 \|C(\dot F)\|^2\asymp\|C(R)\|^2+
 k_{\ell-1}^{-1}\sum_i\|C(B_{\ell,i})\|^2\|U^{(\ell)}_{i:}\|^2.
\]
At the top layer, row perturbations split into radial and perpendicular parts, each row having a single radial pure label. Iterating downwards,
\[
 \|H_mU\|^2\asymp\sum_\ell\frac{\|U_\ell\|_F^2}{k_{\ell-1}k_\ell}
 \asymp p^{-1}\|U\|^2\qquad(U\in\mathcal S_W).\tag{TG.4}
\]
The exact zero directions are
\[
 G_\ell=\operatorname{diag}(t_\ell)W^{(\ell)}-mW^{(\ell)}\operatorname{diag}(t_{\ell-1}),
 \quad t_0=t_L=0.
\]
They come from internal positive diagonal rescalings, of dimension $g=\sum_{\ell<L}k_\ell$; checking the first nonzero $t_\ell$ layer by layer shows the parametrization is injective. Removing radial parts row by row and passing the compensation to the next layer gives a unique decomposition $U=V+G$, $V\in\mathcal S_W$. Lower bounds on row lengths and upper bounds on column norms make this normalized map bounded, so
$\operatorname{dist}(U,\mathcal G_W)\le\|V\|\le C\operatorname{dist}(U,\mathcal G_W)$.
\begin{theorem}[Homogeneous bulk spectrum]\label{tg:homogeneous-bulk}
With probability tending to one, $pH_m^*H_m$ has exactly $g=O(d)$ zero eigenvalues and all others lie in a fixed interval $[c_*,C_*]\subset(0,\infty)$. In particular $\|H_m\|_{\rm HS}^2=O(1)$.
\end{theorem}
The proof also gives the pointwise inequality
$c_*p^{-1}\operatorname{dist}(U,\mathcal G_W)^2\le\|H_mU\|^2\le C_*p^{-1}\operatorname{dist}(U,\mathcal G_W)^2$, covering all parameter blocks.

\subsection{Recursive projections and transport of deletion modules}
\paragraph{The projections are orthogonal projections in input space.}
For $1\le\ell<L$ let $Q_\ell$ keep the square-free monomials of degree $q_\ell$ in the formal neuron coordinates, and $Q^-_{\ell,i}$ keep the square-free monomials of degree $q_\ell-1$ not containing $i$. $S_{\ell+1}$ denotes the substitution $y_j\mapsto(w_{\ell+1,j}\cdot y)^m/\sqrt{m!}$.
Starting from $T_{L-1}=\operatorname{ran}Q_{L-1}$, $Z_{L-1,i}=\operatorname{ran}Q^-_{L-1,i}$, define recursively
\begin{align*}
 Z_{\ell,i}&=\operatorname{ran}\left[(B_j)_j\mapsto Q^-_{\ell,i}
     \sum_j(w_{\ell+1,j}\cdot y)^{m-1}S_{\ell+1}B_j\right],\quad B_j\in Z_{\ell+1,j},\\
 T_\ell&=\operatorname{ran}\left[(A,(B_j\otimes v_j)_j)\mapsto Q_\ell
  \left(S_{\ell+1}A+\sum_j(w_{\ell+1,j}\cdot y)^{m-1}(v_j\cdot y)S_{\ell+1}B_j\right)\right],
\end{align*}
where $A\in T_{\ell+1}$, $v_j\perp w_{\ell+1,j}$, extended linearly from simple tensors. Each module carries the inherited norm and an orthonormal basis. Square-freeness and avoidance of the marked label prevent repeated coefficient labels, so these maps are restrictions of \eqref{tg:frame}, and \eqref{tg:collision}--\eqref{tg:marked} preserve their positive lower bounds.
The bottom layer uses the same pure/defect substitution without pruning physical input coordinates:
\[
 \mathcal F_{\rm deep}(A,(B_i\otimes u_i)_i)
 =S_1A+\sum_i(w_{1,i}\cdot x)^{m-1}(u_i\cdot x)S_1B_i,\qquad u_i\perp w_{1,i}.
\]
It has a fixed frame lower bound, so
\[
 P_{\rm deep}=\mathcal F_{\rm deep}(\mathcal F_{\rm deep}^*\mathcal F_{\rm deep})^{-1}\mathcal F_{\rm deep}^*
 \tag{TG.5}
\]
is an orthogonal projection, set to zero when the frame event fails. Transport to the input chaos by $I_N/\sqrt{N!}$. All actual row radii are kept in the substitution.

The full index definition of the prefix projections is as follows. Fix $1\le s\le L-1$ and set
\[
 T_s^{[s]}=\operatorname{ran}Q_s,\qquad
 Z_{s,i}^{[s]}=\operatorname{ran}Q^-_{s,i}\quad(1\le i\le k_s).
\]
For $\ell=s-1,s-2,\ldots,1$ define successively
\begin{align*}
 Z_{\ell,i}^{[s]}&=\operatorname{ran}\left[(B_j)_{j=1}^{k_{\ell+1}}\mapsto
 Q^-_{\ell,i}\sum_{j=1}^{k_{\ell+1}}(w_{\ell+1,j}\cdot y)^{m-1}S_{\ell+1}B_j\right],
 \\[-2pt]&\hspace{1em} B_j\in Z_{\ell+1,j}^{[s]},\\
 T_\ell^{[s]}&=\operatorname{ran}\left[(A,(B_j\otimes v_j)_j)\mapsto Q_\ell
 \left(S_{\ell+1}A+\sum_{j=1}^{k_{\ell+1}}(w_{\ell+1,j}\cdot y)^{m-1}
 (v_j\cdot y)S_{\ell+1}B_j\right)\right],
 \\[-2pt]&\hspace{1em} A\in T_{\ell+1}^{[s]},\quad v_j\in w_{\ell+1,j}^{\perp}.
\end{align*}
The deletion modules are indexed by $1\le i\le k_\ell$; the domain of the second map is
$T_{\ell+1}^{[s]}\oplus\bigoplus_j(Z_{\ell+1,j}^{[s]}\otimes w_{\ell+1,j}^{\perp})$.
The domain of the bottom map is $T_1^{[s]}\oplus\bigoplus_i(Z_{1,i}^{[s]}\otimes w_{1,i}^{\perp})$,
with the same formula as $\mathcal F_{\rm deep}$; the orthogonal projection $P_{\le s}$ is defined through its Gram inverse.
For $s=1$ the intermediate recursion is empty, but all pure and physical defect blocks of the bottom layer are included.
Let $\mathcal E_{\le s}$ be the intersection of the events on row lengths, coherences, maximal coordinates, matrix norms and the frame bounds above for the first $s$ layers;
$P_{\le s}$ is set to zero only on its complement. The event for the full projection is the intersection of all required prefix events.
Since $T_s\subset T_s^{[s]}$, $Z_{s,i}\subset Z_{s,i}^{[s]}$ and the layerwise linear maps preserve inclusions,
$\operatorname{ran}P_{\rm deep}\subset\operatorname{ran}P_{\le s}$, including the convention on failure events.
Hence $P_{\rm deep}\preceq P_{\le s}$, and the latter does not depend on $W^{(s+1)},\ldots,W^{(L)}$ at all.
This is needed for the later conditional orthogonality over higher-layer weights.

The diagonal norm of the square-free pure frame is $(m!)^q/(mq)!$, and that of the deleted frame is $(m-1)!(m!)^{q-1}/(mq-1)!$. With the off-diagonal block bounds, the pruning bounds and convergence of row radii, each full Gram equals $\lambda I+o_{\rm op}(1)$, $\lambda>0$. The defect Gram is not scalar, but its extra rank lies in the span of finitely many labels. After inserting physical parameters $e_a/\sqrt{k}$, its squared HS cost is
\begin{equation}\label{tg:rankcost}
 \frac{Ct}{k}\sum_{r,i}\|\beta_{r,i}\|^2,
\end{equation}
where $t$ is the fixed support size and $\beta$ the source coefficients; this must be used after the total coefficient energy is bounded.

To see that deletion modules lose no width, let $H=\bigoplus_jH_j$ be the upper deletion coefficient space and $V=\bigoplus_jV_j\subset H$ the target module. The pruned deletion map $D_i$ with mark $i$ satisfies
$D_i^*D_i=\lambda I+E_i$, $\sup_i\|E_i\|\le\eta_d\to0$. If $P_i$ projects onto $D_iV$, then for $x_i=(w_{ji}C_j)_j$
\begin{align}
 \sum_i\|P_iD_ix_i\|^2
 &\le C\sum_i\left(\|\Pi_Vx_i\|^2+\eta_d^2\|x_i\|^2\right)\notag\\
 &\le C\sum_j\|w_j\|^2\left(\|\Pi_{V_j}C_j\|^2+\eta_d^2\|C_j\|^2\right).\label{tg:delete-transport}
\end{align}
The first line uses the target frame lower bound and $\Pi_VD_i^*D_ix_i$; the second uses $\sum_iw_{ji}^2=\|w_j\|^2$, which cannot be replaced by $\sum_i1$. This can be tensored with the raw parameter-direction space and iterated over layers. Defect targets of different layers meet pure/non-radial defect cross blocks of norm $o(1)$ at the first branching interface; same-layer matching uses the scalar isotropic blocks and \eqref{tg:rankcost}.

\subsection{Values of actual features and full-gradient tail estimates}
\paragraph{Layerwise input-chaos representation from the weight expansion.}
Given previous-layer features $h\in\mathbb R^k$, let $q=\|h\|^2/k$, $v=h/\|h\|$,
$c_r(q)=\mathbb E\sigma(\sqrt qG)\psi_r(G)$. The degree-$r$ term in an independent current row $V_j$ is $c_r(q)\psi_r(V_j\cdot v)$. For $z=z_\parallel v+z_\perp$, radial and angular Hermite orthogonality give
\begin{align}
 \mathbb E_V|D_h[c_r(q)\psi_r(V\cdot v)][z]|^2
 &=k^{-1}\left(4q(c'_r(q))^2|z_\parallel|^2+
              r c_r(q)^2q^{-1}\|z_\perp\|^2\right),\label{tg:radial}\\
 \mathbb E_V\|\nabla_V[c_r(q)\psi_r(V\cdot v)]\|^2&=r c_r(q)^2.\notag
\end{align}
The angular derivative introduces independent Gaussian coordinates orthogonal to $v$, so the cross term vanishes. Root terms of positive degree have mean zero; lower-layer gradients are orthogonal across different root rows, and root-parameter gradients occupy disjoint parameter blocks, so the all-ones readout does not enlarge these tails.

For the smooth odd branch $q\le\|\sigma\|_\infty^2$ and
\[
 \sqrt{r/q}\,c_r(q)=\langle\sigma'(\sqrt qG),\psi_{r-1}\rangle,\quad
 2\sqrt q\,c'_r(q)=\langle G\sigma'(\sqrt qG),\psi_r\rangle.
\]
The functions on the right are continuous into Gaussian $L^2$ on a compact interval containing the endpoint zero, so their image is compact. Finite Hermite projections converge strongly with operator norm at most one, hence uniformly on this compact set. Therefore
\[
 \tau_R:=\sup_q\sum_{r>R}\left(c_r(q)^2+r c_r(q)^2/q+4q(c'_r(q))^2\right)\longrightarrow0.
 \tag{TG.6}
\]
$q=0$ is defined by continuity of the right sides; no inverse-norm assumption is used.
For the piecewise-linear branch $c_r(q)=\sqrt q\,a_r$ ($r\ge1$), so the tail is controlled by
$\sum_{r>R}(r+1)a_r^2\to0$ times a fixed power of $1+q$; this converges by the first-order Gaussian Sobolev identity $\sum r a_r^2=\mathbb E|\sigma'(G)|^2<\infty$.

At fixed depth and any fixed moment order, the moments of $\|J_h\|_{\rm op}$ and $q$ are uniformly bounded. Concretely, the current parameter block has norm at most $\|\sigma'\|_\infty\|h^{(\ell-1)}\|/\sqrt{k_{\ell-1}}$; earlier parameter blocks are multiplied successively by bounded gates and $W^{(\ell)}/\sqrt{k_{\ell-1}}$, whose operator norms have uniform moments of all fixed orders. The moment inequality for conditionally independent row means and Lipschitz continuity of the variance profile near one give by induction
\[
 \|q_\ell-1\|_{L^s}=O(d^{-1/2}),\qquad
 \|\nabla q_\ell\|\le2\sqrt{q_\ell/k_\ell}\|J_{h^{(\ell)}}\|_{\rm op},
 \tag{TG.7}
\]
The induction starts from the chi-square mean $q_0=\|X\|^2/d$. Outside compact intervals, use Markov bounds with arbitrarily high fixed moments.

For a product of $t$ \emph{distinct} current rows, expand the tail error by a product telescope and apply the product rule to the gradient. Each differentiated row is conditionally independent of the other $t-1$ rows; Cauchy--Schwarz is used only over the fixed $t$ terms. The other rows have conditional second moments that are fixed polynomials in $1+q$. When \eqref{tg:radial} is assembled into $J_h$,
$\operatorname{tr}(J_hJ_h^*)\le k\|J_h\|_{\rm op}^2$ cancels the $1/k$ of the angular term, so
\begin{equation}\label{tg:product-tail}
 \mathbb E\left[\|E_R\|_{L_X^2}^2+
 \|\nabla_{W_{\le\ell}}E_R\|_{L_X^2}^2\right]\le C_t\tau_R\longrightarrow0.
\end{equation}
No total parameter count $p$ appears, and no high-power convergence of the Hermite tail of a single row is needed.

The multivariate coefficients of the fixed-degree expansion are $d_r(q)=c_r(q)q^{-r/2}$. For the odd branch, by Stein's identity
\[
 d_r(q)=\mathbb E\sigma^{(r)}(\sqrt qG)/\sqrt{r!},\qquad
 d'_r(q)=\mathbb E\sigma^{(r+2)}(\sqrt qG)/(2\sqrt{r!})
\]
they are bounded on the whole interval; for the piecewise-linear branch $d_r(q)=a_rq^{(1-r)/2}$ is bounded with its derivatives on $[1/2,2]$. Using (TG.7) to replace $d_r(q)$ by $a_r$ in finitely many terms, value and full-gradient errors vanish. Outside the interval the latter branch does not estimate inverse powers directly: the exact radial/angular metrics are $a_r^2$ and $r a_r^2$, with denominators already cancelled; conditional fixed-order Gaussian hypercontractivity, moment bounds for $q,J_h$ and H\"older make the rare-interval contribution vanish. The flattened terms are fixed polynomials with fixed moment bounds.

\subsection{Normal-ordered collisions, finite-support separation and all parameter blocks}
At a fixed row degree $r$ the exact expansion is
\[
 d_r(q)\sum_{|\alpha|=r}\sqrt{r!/\alpha!}\,k^{-r/2}h^\alpha\psi_\alpha(V).
 \tag{TG.8}
\]
First flatten, then delete repeated child labels at the same interface; the child features are still actual functions. With $t$ child slots there are only $O(k^{t-1})$ collision choices, with squared normalization $k^{-t}$. Normal-ordered multi-index Hermite orthogonality gives the sum of squares directly; for fixed differentiation weights $\alpha\mapsto\alpha-e_a$ is injective, so gradients have the same orthogonality. Fixed-degree product moments of repeated actual features are bounded, so the collision error is $O_R(k^{-1})$. Only after repeated labels are removed at every interface is each first-order Hermite edge written as a plain weight variable; raw even weight powers cannot replace normal-ordered Hermite polynomials beforehand.

The truncation order must be $R_L,R_{L-1},\ldots,R_2$. Each time a degree truncation is fixed, first flatten and delete child collisions, then apply \eqref{tg:product-tail} to products of distinct child rows of the next layer. The whole error at interface $s$ can be written as
$\sum_gb_g\Psi_g(W_{>s})E_g$, where orderings with the same edge set are merged and $\Psi_g$ are orthogonal square-free monomials in higher-layer edges. Counting upper labels gives
\[
 \sum_gb_g^2\le C_{R_{>s}},\quad
 \sum_{e>s}\sum_gb_g^2\mathbb E|\partial_e\Psi_g|^2
 =\sum_gb_g^2\#\{\text{higher-layer edges of }g\}\le C_{R_{>s}}.
\]
Hence the total error, including higher-layer derivatives, is at most $C_{R_{>s}}\tau_{R_s}$. New truncation choices do not affect earlier constants. Given $\ecut$, choose finite truncations successively so that all squared tail errors are below $C_L\ecut^2$, and only then let $d\to\infty$.

\paragraph{Finite graphs and phase indices.}
With all truncations fixed, a shape $\mathfrak t$ is an ordered tree rooted at layer $L$; each vertex at layer $\ell\ge2$
$v$ has $1\le r_v\le R_\ell$ ordered children, all at layer $\ell-1$, and layer-one vertices are actual ridge leaves.
Let $V_\ell(\mathfrak t)$ be its set of layer-$\ell$ vertices, $t_\ell=|V_\ell(\mathfrak t)|$, $t_L=1$.
These numbers depend only on the finite shape, not on the width. A label assignment is a family of injections
$\iota_\ell:V_\ell(\mathfrak t)\hookrightarrow\{1,\ldots,k_\ell\}$.
The edge $v\to u$ corresponds to $W^{(\ell)}_{\iota_\ell(v),\iota_{\ell-1}(u)}$;
global injectivity within each layer ensures that different edges use different Gaussian variables and that every non-root vertex has a single parent.
For one ordered representation the scalar term is exactly
\[
 k_L^{-1/2}\prod_{\ell=2}^Lk_{\ell-1}^{-t_{\ell-1}/2}
 \prod_{v\in\cup_{\ell\ge2}V_\ell}\frac{a_{r_v}}{\sqrt{r_v!}}
 \prod_{v\to u}W^{(\ell)}_{\iota(v),\iota(u)}
 \prod_{u\in V_1}\sigma(w_{1,\iota(u)}\cdot X).
 \tag{TG.8a}
\]
This formula expands powers by ordered children. If an unordered edge graph has several ordered representations, their coefficients are added first; there are at most
$\prod_v r_v!$ of them, a number fixed by the truncation. This gives a finite scalar sum $G_R$, and its actual parameter derivative
$J(G_R)$ is its tangent operator. Orthogonality is never claimed separately for repeated representations.

When testing interface $s\ge2$, fix a shape and keep the edges above layer $s$ with their endpoint labels; call this higher edge set $g$.
$g$ fixes the finite source support $A_g=\iota_s(V_s)$ at layer $s$ but no label at or below layer $s-1$.
Phases and normalized coefficients are defined by
\[
 \Psi_g=\prod_{\substack{v\to u\\\operatorname{layer}(v)>s}}
 W^{(\operatorname{layer}(v))}_{\iota(v),\iota(u)},\qquad
 b_g=k_L^{-1/2}\prod_{\ell=s}^{L-1}k_\ell^{-t_\ell/2}\,c_{\mathfrak t,g}.
 \tag{TG.8b}
\]
$H_g:=b_g\Psi_g$ denotes the higher-layer factor with its normalization. $c_{\mathfrak t,g}$ contains the fixed higher-layer scalar coefficients and the multiplicity of orderings, bounded by truncation constants.
The current-layer coefficient $\prod_{v\in V_s}a_{r_v}$ stays in the source polynomial coefficients.
For a fixed shape, the $\Psi_g$ with different edge sets $g$ are orthonormal.
If different shapes or degree profiles produce the same higher edge set, their finitely many source terms are first added within that $g$;
equivalently, apply Cauchy--Schwarz over the finitely many shapes and take orthogonal sums shape by shape.
This finite factor depends at most on $L,R_2,\ldots,R_L$ and not on the dimension.

Once the lower part is balanced, all lower-layer assignments sharing $g$ add up through the substitution operator acting on (TG.9):
the ordered sum of the $r_i$ children of current row $i\in A_g$ is exactly
$(w_{s,i}\cdot Y)^{r_i}/\sqrt{r_i!}$; label injectivity within the layer amounts to applying $Q_{s-1}$ to the whole product.
Lower layers use the same square-free substitution. Hence lower-layer labels are contained in a single coefficient tensor,
form no independent phases and do not enter the coefficient count of (TG.8b). For $s=1$, by the same definition all edges at layer two and above are phases,
and the source is the finite product of bottom-layer ridges rather than the polynomial in (TG.9).

\paragraph{Interfacewise errors.}
Starting from $J(G_R)$, in the bottom layer first replace the terms with source support size $q_1$ by
$a_m^{q_1}$ times the homogeneous bottom-layer value and its exact derivative, delete all other terms, and call the resulting tangent sum $K^{(1)}$.
Let $K^{(0)}=J(G_R)$. For $s=2,\ldots,L-1$, in $K^{(s-1)}$ keep the current profiles with
$r_i=m$ and $|A_g|=q_s$ and delete the others, obtaining $K^{(s)}$.
These operations replace explicit finite scalar graph terms and their derivatives; no random projection is differentiated.
Define $E_s=K^{(s-1)}-K^{(s)}$. Then
$J(G_R)-K^{(L-1)}=\sum_{s=1}^{L-1}E_s$ is an exact telescoping identity.
The bottom layer fixes the number of actual leaves, and each subsequent correct profile divides the number of upper vertices by $m$;
hence the final root degree can only be $m$, and $\Pi_NK^{(L-1)}=a_m^S H_m^{\rm inj}$.
The final comparison applies $\Pi_N$ and then $P_{\rm deep}$ to the whole telescoping identity;
$K^{(s)}$ itself is the graph derivative across all chaoses and is not mistaken for the top chaos.
Coefficient and radius errors with the same support in the bottom-layer replacement are included in $E_1$ rather than in an unstated remainder.

We now prove separation of the finite graphs. Lower-layer assignments sharing higher-edge phases must be combined and cannot be treated as orthogonal. First use $P_{\le1}$: for a fixed bottom-layer source support $A$, the degree-$N$ tensor of the actual ridge product lies in
\[
 \operatorname{Sym}^N H_A+\operatorname{Sym}(\operatorname{Sym}^{N-1}H_A\otimes H_A^\perp),
 \qquad H_A=\operatorname{span}\{w_{1,i}:i\in A\}.
\]
The coefficient operator of a physical source perturbation has scale $d^{-1/2}$. If the target has $h$ new labels, at most two pairings are absorbed by defects and the rest give $C_R\delta^{mh-2}$; after squaring and summing, at most $C_R\sum_hk_1^h\delta^{2mh-4}$. If the target is contained in the source support, at orthogonal unit rows the extra undifferentiated actual features have mean zero; the case with a single extra differentiated feature requires the target's $m-1$ coefficient, which also vanishes. When the source support equals the target, the pure coefficient is $a_m^{q_1}$ and the external defect is its exact Gaussian derivative. Physical directions inside the source support have fixed rank and squared HS cost $O(d^{-1})$. For non-orthogonal rows, finite-dimensional Gaussian integrals are locally Lipschitz in the covariance, giving $O_R(\delta)$ coefficient errors; for the piecewise-linear branch one differentiates the density and dominates by polynomial envelopes, so no second-order smoothness of the activation is needed. The maximal deviation of actual radii enters the same bound.

Assume the lower interfaces are balanced. All lower-layer assignments for one higher phase at interface $s$ must be combined into
\[
 A_A(Y)=\prod_{i\in A}(w_{s,i}\cdot Y)^{r_i}/\sqrt{r_i!},\qquad D=\sum_i r_i=mq_s.
 \tag{TG.9}
\]
whose square-free restriction is the injective child-label sum. The two-slot proof for fixed supports gives
$\|C_D C(A_A)\|^2\le C_R\kappa^2$, and the same for the defect/deletion versions; hence restoring the full (TG.9) first costs only a vanishing error. This step does not use the full frame lower bound, which fails for low-degree sources.
At orthogonal rows, pure matching requires $A=B$ and all $r_i=m$. External defect and deletion matching require
\[
 r-e_i=m\mathbf1_B-e_j.
\]
Besides the correct matching, only shifts by $m-1,m+1$, or an extra first-order node combined with the target's $m-1$, can occur; all are excluded because the adjacent Hermite coefficients vanish. New target labels give the coherence powers above; internal source directions use \eqref{tg:rankcost}, and the pairing is then transported to the input through the pure scalar Gram and \eqref{tg:delete-transport}.

All coefficient energies can be checked directly. For fixed-degree symmetric tensors one has the exact identities
\[
 \sum_a\|C_{D-1}(\partial_a A_A)\|^2=D^2\|C_D(A_A)\|^2\le C_R,
 \quad
 \sum_{i\in A,a}\|C_D(\partial_{W^{(s)}_{ia}}A_A)\|^2\le C_R.\tag{TG.10}
\]
The second identity inserts $e_a/\sqrt{k_{s-1}}$, so the dimension cancels after summation. For any frozen coefficient tensor $A$, the partial derivatives through the lower balanced substitutions satisfy
\begin{equation}\label{tg:frozen}
 \sum_{\ell<s}\|D^{\rm prefix}_{W^{(\ell)}}(\text{balanced substitution of }A)\|_{\rm HS}^2\le C_{L,R}\|A\|^2.
\end{equation}
The proof uses the defect frame upper bound at the differentiated layer, with a factor $1/k_{\ell-1}$ per parameter direction; after summing over directions, (TG.10) handles the neuron sum, and the other layers use the pure frame upper bound. These are partial derivatives at frozen $A$; in the application $A_A$ depends only on $W^{(s)}$, not on the differentiated lower-layer parameters, so this is exactly the required actual derivative.

Parameter blocks above $s$ lower the higher phases; a deleted edge is reinserted uniquely, so orthogonality persists. The current parameter block uses (TG.10), and blocks below $s$ use \eqref{tg:frozen} and deletion transport. Counting only labels above interface $s$ gives
\[
 \sum_g\mathbb E|H_g|^2\le C_R
 \left[k_L^{-1}\prod_{\ell=s}^{L-1}k_\ell^{-t_\ell}\right]
 \left[k_L\prod_{\ell=s}^{L-1}k_\ell^{t_\ell}\right]\le C_R.\tag{TG.11}
\]
Lower-layer assignments are already merged in (TG.9) and are not counted twice. Since $P_{\le s}$ depends only on the prefix, conditional orthogonality survives the projection. Indicators of future-weight events must not be inserted into this computation.

\paragraph{From per-phase operator bounds to full-parameter Hilbert--Schmidt bounds.}
We write the sums above as identities, so that a pointwise bound for one tensor is not mistaken for a bound on the whole family.
Fix a shape and a wrong interface $s$, and write its scalar source in $E_s$ as
$\sum_g b_g\Psi_g R_g(W_{\le s},X)$;
at the bottom layer $s=1$, $R_g$ may be the difference between the actual product and the chosen homogeneous product.
For $s\ge2$ it is the scalar function of the wrong current profile (TG.9) after the balanced prefix,
with the fixed coefficient $\prod_{i\in A_g}a_{r_i}$ and all lower-layer $a_m$ factors included in $R_g$.
Let
\[
 Z_g^0=P_{\le s}\Pi_NR_g,\qquad
 Z_g^{\rm cur}=P_{\le s}\Pi_ND_{W^{(s)}}R_g,\qquad
 Z_g^{\rm low}=P_{\le s}\Pi_ND_{W_{<s}}R_g.
\]
$Z_g^{\rm cur},Z_g^{\rm low}$ are operators from the corresponding full parameter spaces to the input $L^2$;
$P_{\le s}$ is never differentiated.
On the prefix good event, the pairing bounds above give three bounds uniform in $g$:
\begin{equation}\label{tg:one-phase}
 \|Z_g^0\|^2+\|Z_g^{\rm cur}\|_{\rm HS}^2+
 \|Z_g^{\rm low}\|_{\rm HS}^2\le C_R\omega_{s,d},\quad
 \omega_{s,d}=\Delta_d^2+\delta_d^2+\kappa_d^2+k_{s-1}^{-1}
 +\sum_{h=1}^{q_s}k_s^h\delta_d^{2mh-4}+\rho_{s-1,d}^2.
\end{equation}
For $s=1$, $Z_g^{\rm low}=0$ and $\rho_{0,d}=0$. The coefficient scales are as follows:
pure block coefficients are bounded by $\|C(A_{A_g})\|\le C_R$; the sum of squared coefficients over all current directions by
the second identity in (TG.10); and the sum of squared raw lower-layer defect coefficients by the first identity in (TG.10) and
\eqref{tg:frozen}. Hence each error term carries a single factor $C_R$, not the number of parameters.
New-label blocks, squared and summed over targets, give $\sum_hk_s^h\delta^{2mh-4}$;
near-orthogonal pairings with the same support give $\Delta_d^2+\delta_d^2$;
square-free restoration of the source gives $\kappa_d^2$;
internal defect directions give the $k_{s-1}^{-1}$ of \eqref{tg:rankcost}.
Prefix scalar Gram errors act on the whole coefficient vector before squaring, giving $\rho_{s-1,d}^2$;
the sum over all deletion marks is exactly \eqref{tg:delete-transport}.
The target frame lower bound converts these squared double-frame pairings into projection norms.
These estimates give \eqref{tg:one-phase} uniformly on the three full coefficient spaces, without multiplying a single-direction bound by the dimension.

Now condition on $W_{\le s}$. For the combined lower-layer and current-parameter block, exactly
\[
 \mathbb E_{W_{>s}}\left\|\sum_gb_g\Psi_g
       (Z_g^{\rm cur}\oplus Z_g^{\rm low})\right\|_{\rm HS}^2
 =\sum_gb_g^2\left(\|Z_g^{\rm cur}\|_{\rm HS}^2+\|Z_g^{\rm low}\|_{\rm HS}^2\right).
 \tag{TG.11a}
\]
This is an orthogonal sum of Hilbert-space-valued terms; the $Z_g$ need not be mutually orthogonal.
For a fixed higher-layer parameter coordinate $e=(\ell,i,j)$, let $E(g)$ be the set of higher edge variables and
$\Psi_{g\setminus e}=\prod_{f\in E(g)\setminus\{e\}}W_f$. Then
\[
 \partial_e\Psi_g=\mathbf1_{e\in E(g)}\Psi_{g\setminus e},\qquad
 \mathbb E_{W_{>s}}\left\|\sum_{g:e\in E(g)}b_g\Psi_{g\setminus e}Z_g^0\right\|^2
 =\sum_{g:e\in E(g)}b_g^2\|Z_g^0\|^2.\tag{TG.11b}
\]
Injectivity holds because $E(g)=(E(g)\setminus\{e\})\cup\{e\}$ and all participating terms contain the same fixed edge $e$.
If a higher phase has a single edge, deleting it leaves a constant, but still only one edge set contributes for a fixed parameter;
parameters of different root rows enter the HS sum through different coordinates, so no coherent cross terms arise.
Summing over all higher-layer parameters, the right side of (TG.11b) becomes
$\sum_g b_g^2|E(g)|\|Z_g^0\|^2$, where $|E(g)|=\sum_{\ell=s}^{L-1}t_\ell$ is a fixed number.

Finally, for a fixed shape the number of admissible higher labels is at most
$k_L\prod_{\ell=s}^{L-1}k_\ell^{t_\ell}$, while (TG.8b) gives
$b_g^2\le C_Rk_L^{-1}\prod_{\ell=s}^{L-1}k_\ell^{-t_\ell}$.
Hence
\[
 \sum_g b_g^2(1+|E(g)|)\le C_R.
 \tag{TG.11c}
\]
Combining \eqref{tg:one-phase} with (TG.11a)--(TG.11c), averaging first over higher layers and then over the prefix,
the full-parameter squared error of the shape on the prefix good event is $\le C_R\omega_{s,d}$.
Finite sums of profiles within one higher phase, and over all finite shapes, only add fixed Cauchy--Schwarz factors.
On bad prefix events the fixed truncated sources grow polynomially in the dimension at fixed moment order;
H\"older with a sufficiently high negative power of the failure probability gives $o(1)$.
Since $P_{\rm deep}\preceq P_{\le s}$, this bound also controls the required full projection;
future good-event indicators need not and must not be inserted into (TG.11a)--(TG.11b).

Let $E_s$ be the full gradient of the finite graphs whose first wrong interface is $s$, $\rho_{s-1,d}$ the sum of prefix Gram and pruning operator errors, and $\Delta_d$ the maximal deviation of row radii. The blockwise estimates give
\begin{equation}\label{tg:separator}
 \mathbb E\|P_{\rm deep}\Pi_NE_s\|_{\rm HS}^2
 \le C_R\left[\Delta_d^2+\delta_d^2+\kappa_d^2+k_{s-1}^{-1}
 +\sum_{h=1}^{q_s}k_s^h\delta_d^{2mh-4}+\rho_{s-1,d}^2\right]+o(1)=o(1).
\end{equation}
At the first interface the prefix error is absent. Bad prefix events are handled separately by H\"older and arbitrarily high polynomial failure exponents. After induction over the finitely many interfaces, only graphs with every node of degree $m$ and no repeated labels within a layer remain, with coefficient exactly $a_m^S$. Extra centered features at the bottom layer are removed at the level of actual features; balanced leaf counts cannot be manufactured afterwards by contracting even homogeneous polynomials.

\subsection{The collective-mean branch of piecewise-linear activations}
For the odd branch the conditional mean vanishes at every variance. For the piecewise-linear branch the mean is
$\mu(q)=c(\sqrt q-1)$, $c=b\sqrt{2/\pi}/v$. Mean factors inside finite positive-root graphs and their gradients have $L^s$ scale $O(d^{-1/2})$, because
$\|\nabla\mu(q)\|\le C\|J_h\|_{\rm op}/\sqrt k$.
The mean gradient of the all-ones root readout is of first order and must be kept. Let $\psi=\sigma^2$ and let $J_r^\psi$ be the gradient of the all-ones readout of depth $r$ with root activation $\psi$ and lower activation $\sigma$. Then
\[
 \sqrt{k_L}\mu'(q_{L-1})\nabla q_{L-1}
 =\mu'(1)\sqrt{k_L/k_{L-1}}J_{L-1}^\psi+o_{L^2}(1).\tag{TG.12}
\]
The error uses the global cancellation
\[
 \sqrt{k_L}\|[\mu'(q)-\mu'(1)]\nabla q\|
 \le C\sqrt{k_L/k}|1-\sqrt q|\,\|J_h\|_{\rm op},
\]
so no inverse moments are needed near zero variance.

$\psi$ is a sum of a positively homogeneous quadratic, a positively homogeneous linear function and a constant; its positive-degree coefficients are
$c_r^\psi(q)=qA_r+\sqrt qB_r$, with
$\sum_r(r+1)(A_r^2+B_r^2)<\infty$. Substituting into \eqref{tg:radial} gives
\[
 4q(c_r^\psi{}'(q))^2=4(\sqrt qA_r+B_r/2)^2,\qquad
 r(c_r^\psi(q))^2/q=r(\sqrt qA_r+B_r)^2.
\]
The current-weight gradient energy is $r(c_r^\psi(q))^2$. Hence positive-degree tails have the same full-gradient control as before. The fixed-degree flattened coefficients
$d_r^\psi(q)=A_rq^{1-r/2}+B_rq^{(1-r)/2}$ are smooth and bounded on $[1/2,2]$; outside it use the nonsingular metrics just given and fixed moments.

Its zeroth-order mean is
\[
 \nu(q)=\frac{(a^2+b^2)q-2b^2(2/\pi)\sqrt q+b^2(2/\pi)}{v^2},\qquad \nu'(1)=1.
\]
The same global cancellation replaces $\nu'(q)$ by $\nu'(1)$ and pushes the zeroth-order gradient down one layer. The recursion ends at the parameter-free input norm, so (TG.12) equals a finite sum of positive-degree $\psi$-root branches of depth $r=1,\ldots,L-1$, plus $o_{L^2}(1)$; the coefficients are $\mu'(1)\sqrt{k_L/k_r}$, uniformly bounded.

Truncate, flatten and remove collisions in each branch in the order above, and first balance all $\sigma$ nodes below the root. If the root layer is $r<L$, the preserved leaf count forces the root degree
$D=m^{L-r+1}=mq_r$, $q_r\ge m$. If this degree lies outside the finite truncation, the remainder is already in the tail error; otherwise the source is a single-ridge pure/defect/deletion frame of degree $D$, whereas the target has $q_r$ distinct labels of degree $m$ each.
The block pairing of the two full frames satisfies
\[
 \|\mathcal F_{{\rm bal},B}^*\mathcal F_{{\rm sing},j}\|_{\rm op}
 \le C\begin{cases}\delta^{D-m-2},&j\in B,\\\delta^{D-2},&j\notin B.\end{cases}
\]
since at most $m$ raw slots can match the same label and at most two more are absorbed by defects. With $u=\delta^{D-m-2}$, $v_0=\delta^{D-2}$, the row and column Schur test bounds the whole source frame by
\begin{equation}\label{tg:collective-angle}
 \|\mathcal F_{\rm bal}^*\mathcal F_{\rm sing}\|_{\rm op}^2
 \le C(u+k_rv_0)(k_r^{q_r-1}u+k_r^{q_r}v_0)=o(1).
\end{equation}
The dimension power of the leading term is $d^{-(m-1)q_r+m+1}$, strictly negative. The source frame has degree $D\ge m^2$, hence fixed upper and lower bounds as well. Pruning losses are $o(1)$; the lower-layer pure/deletion scalar Grams and \eqref{tg:rankcost} transport this small angle to the input.
Conditional orthogonality over root weights is not used here, since the projection depends on the same root weights. The source coefficient vector of the all-ones readout has unit squared norm; the coefficient energy summed over root-parameter directions is $O(1)$; lower-layer energies are controlled by (TG.10), \eqref{tg:frozen} and the single-ridge frame upper bound. Hence the projected HS energy of each finite collective branch vanishes, and the limit of all tails is bounded by $C_L\ecut^2$.

\subsection{The full bridge, normalized HS transfer and the order of limits}
Let $H_m^{\rm inj}$ be the homogeneous tangent operator that keeps square-free monomials at every formal interface. Coordinate projections are fixed linear maps and commute with weight differentiation. Compare the full and the pruned homogeneous networks interface by interface: higher-layer parameters differentiate coefficients and use \eqref{tg:collision}; current parameters use the defect version with an inserted direction; lower-layer parameters use the frozen-prefix bound \eqref{tg:frozen} on the lost coefficient tensor. The pure/defect frame upper bounds and (TG.10) keep all coefficient energies bounded. A finite telescoping sum gives
\begin{equation}\label{tg:homogeneous-trim}
 \mathbb E\|H_m^{\rm inj}-H_m\|_{\rm HS}^2=o(1).
\end{equation}
This step applies to even $m$ and does not mistake raw even root-weight collisions for centered terms. Outside the good events fixed-degree polynomial moments grow at most polynomially, so H\"older with a high enough failure exponent gives vanishing expectations.

\begin{theorem}[Actual full-parameter tangent bridge]\label{tg:actual-bridge}
For both activation classes, at fixed depth and proportional width, the actual orthogonal projection on input space defined in (TG.5) satisfies
\[
 \mathbb E_W\|P_{\rm deep}\Pi_NJ_W-BH_m(W)\|_{\rm HS}^2\longrightarrow0,
 \qquad A_W\succeq(P_{\rm deep}\Pi_NJ_W)^*(P_{\rm deep}\Pi_NJ_W).
\]
\end{theorem}
The proof first fixes $\ecut>0$, chooses finite truncations successively, and removes all wrong finite branches with \eqref{tg:separator} and \eqref{tg:collective-angle}, giving a limiting squared error $C_L\ecut^2$ relative to $BH_m^{\rm inj}$. Then use \eqref{tg:homogeneous-trim} and finally let $\ecut\downarrow0$. For $L=1$ no interface recursion is needed: the single-ridge pure/defect frames, finite-dimensional Gaussian coefficients and continuity in the radius give the same conclusion, with non-radial error $o(1)$ and total radial energy error $O(d^{-1})$. The PSD inequality is just the contraction property of an orthogonal projection on input space.

To check that the HS scale is the right one, let $K_d=\sqrt pP_{\rm deep}\Pi_NJ_W$, $M_d=B\sqrt pH_m$; then
$p^{-1}\mathbb E\|K_d-M_d\|_{\rm HS}^2\to0$.
The sum of squared singular-value differences is at most the squared Hilbert--Schmidt distance; for fixed $t>0$ the function
$f_t(s)=t/(s^2+t)$ on $s\ge0$ has Lipschitz constant $O(t^{-1/2})$. Cauchy--Schwarz gives
\[
 \mathbb E\left|p^{-1}\sum_{i=1}^p[f_t(s_i(K_d))-f_t(s_i(M_d))]\right|\longrightarrow0.
\]
By the bulk spectrum theorem and PSD monotonicity,
\begin{equation}\label{tg:regularized}
 D_*(t):=\limsup_d\mathbb E\frac1p\operatorname{tr}\{t(pA_W+tI)^{-1}\}
 \le\frac{t}{B^2c_*+t},\qquad
 \int_0^1D_*(t)\frac{dt}{t}\le\log(1+(B^2c_*)^{-1}).
\end{equation}
This proves the integrability condition with the dimension limit first and regularization removed afterwards. The finite-dimensional homogeneous Gram has exact gauge zeros, so no claim is made that the expected unregularized log-determinant is finite at finite dimension, nor that the HS bridge controls rescaled operator norms.

The scalar conditions are verified in the main text; \eqref{tg:regularized} gives the spectral-deficit control needed for the information lower bound.

\section{Raw tree kernels, fresh-coordinate marks and local estimates for literal residuals}
\label{fp:core-section}

This section proves only local statements that are invoked repeatedly. All spaces are real Euclidean spaces and their Hilbert
tensor products; $\|\cdot\|_{\mathrm{HS}}$ is the Hilbert tensor norm,
and $\operatorname{Flat}_{S|S^c}$ the matricization grouping the specified slots.
Layer widths lie in $[c d,Cd]$; the depth, degrees, number of colors and number of operations are fixed.
Write $\delta_d=L_d/\sqrt d$, $\rho_d=L_d/d$, where $L_d$ is a fixed power of a logarithm that may increase after finitely many
operations. All probabilistic events are defined under each Gaussian marginal separately;
once the events hold the arguments are deterministic, so the two networks may be coupled arbitrarily.

\subsection{Invariants of raw tree kernels and exact linear compression}
\label{fp:core-invariant-section}

\paragraph{Raw tree kernels.}
Start from the normalized readout vector of the original network and expand top-down along the original computation tree.
Each current mode carries the index of an original node; one expansion of mode $p$ is
\begin{equation}
 J e_a=\bigotimes_{t=1}^{q}u_a^{(t)},\qquad q\geq1,
 \qquad \|U_t\|_{\mathrm{op}}\leq C,
 \quad U_t=(u_a^{(t)})_a .
 \label{fp:core-replacement}
\end{equation}
One factor may be $e_a$, meaning that the original parent index is kept; a mode may also be stopped and kept as a
frozen leaf. Several physical factors of one parent index share that index and do not create independent marks.
Adaptive cross-network maps may act only on frozen leaves, each producing exactly one physical slot.
Nonlinear features with several output slots must first be written into the kernel as explicit tree expansions.

\refstepcounter{theorem}\paragraph{Lemma~\thetheorem\ (full lifts are bounded).}
The full lift operator of any finite sequence of such expansions is bounded. Concretely, with
$\Delta_q e_a=e_a^{\otimes q}$ one has
$J=(\bigotimes_tU_t^*)\Delta_q$, and $\Delta_q$ is an isometric embedding, so
$\|J\|\leq\prod_t\|U_t\|$.
Similarly, if $F_j:e_a\mapsto f_a^{(j)}$ are bounded, the product features with a shared index
$e_a\mapsto\bigotimes_jf_a^{(j)}$ are bounded.
In Gram form: the Gram matrix is $\mathop{\circ}_jF_j^*F_j$,
and for positive semidefinite $X\preceq bI,Y\succeq0$,
$X\circ Y\preceq b\operatorname{diag}(Y)$.

\refstepcounter{theorem}\paragraph{Lemma~\thetheorem\ (deterministic closure of kernels).}
\label{fp:core-closure}
Let $E$ have $r$ modes with
\begin{equation}
 \|E\|_{\mathrm{HS}}\leq C,\qquad
 \max_{p,a}\|E(e_a\text{ in mode }p,\cdot)\|_{\mathrm{HS}}\leq\delta_d,
 \qquad
 \|\operatorname{Flat}_{S|S^c}E\|_{\mathrm{op}}\leq\delta_d
 \quad(\varnothing\ne S\ne[r]).
 \label{fp:core-invariants}
\end{equation}
Except when $r=q=1$, one expansion preserves these bounds up to a fixed constant factor.

\emph{Proof.}
The Hilbert norm bound follows from boundedness of full lifts. Consider a nontrivial bipartition of the new kernel.
If the whole new block lies on one side, the bipartition is a nontrivial bipartition of the old kernel composed with bounded operators on both sides.
If the new block straddles both sides, temporarily keep the orthogonal parent index $a$ on both sides. The central operator is block diagonal in $a$,
each block being a matricization of the old slice $E(e_a,\cdot)$, whose operator norm is at most its Hilbert norm
$\delta_d$. Applying the row maps on both sides concludes, including the case where the old kernel has a single mode.

Coordinate slices of old modes are controlled directly by bounded lifts. If coordinate $j$ of the $t$-th factor of the new block is fixed
and $q\geq2$, keep the orthogonal parent label in the other new factors to obtain the squared norm
\[
 \sum_a |(U_t)_{aj}|^2\|E(e_a,\cdot)\|_{\mathrm{HS}}^2
 \leq\delta_d^2\|U_te_j\|^2\leq C\delta_d^2.
\]
The remaining map is bounded. If $q=1,r\geq2$, the new slice is
$E(U_1e_j,\cdot)$, controlled by the first matricization norm of the old kernel.
When $r=q=1$, a bounded matrix may concentrate a spread vector on one coordinate, so this conclusion fails.
\hfill$\square$

\paragraph{The Gaussian unary exception and raw tree instances.}
Let $A_\ell$ be independent Gaussian matrices normalized by the square root of the input width.
Every fixed consecutive product $P=A_b\cdots A_a$ satisfies
\begin{equation}
 \|P\|_{\mathrm{op}}\leq C,\qquad
 \max_{ij}|P_{ij}|\leq C_M\sqrt{\log d/d}
 \label{fp:core-product-event}
\end{equation}
with failure probability $O(d^{-M})$. The operator norm follows by multiplying layers; conditionally on
$Q=A_{b-1}\cdots A_a$, $(A_bQ)_{ij}$ is Gaussian with variance
$\|Qe_j\|^2/n_{b-1}\leq C/d$, and tail bounds with a finite union give the second bound.
The coordinate maximum of the readout vector times a consecutive Gaussian product is bounded the same way by conditioning on the last unexposed
Gaussian factor. Hence $r=q=1$ chains before the first branching can be handled separately;
after the first branching the number of modes never decreases and the previous lemma handles all later steps.
Frozen or identity steps need no probabilistic estimate. All finite raw tree kernels satisfy
\eqref{fp:core-invariants}. This requires that no adaptive cross-network matrix be inserted inside a unary chain.

\paragraph{External maps and reconstruction.}
Bounded single-slot maps applied outside each kernel mode preserve all nontrivial matricization bounds but not necessarily
the coordinate-slice bounds. After a pruning step, the kernel must be rebuilt from the original tree: stop at the original parent of the newly completed block
and place its metric on a frozen leaf. ``An arbitrary contraction of the old kernel still satisfies the slice bounds'' cannot replace this reconstruction.
Frozen leaves of different depths may coexist, provided no adaptive metric lies inside a root-to-leaf path.

\refstepcounter{theorem}\paragraph{Lemma~\thetheorem\ (exact compression of linear slots).}
\label{fp:core-linear}
Let the bottom-layer rows be $b_j$, the parent matrix $U=(u_{aj})$, and
$z_a=\sum_ju_{aj}b_j$. If a parent node has $t$ linear child slots and nonlinear child slots of degrees
$q_1,\ldots,q_h\geq2$, its unsymmetrized tensor is exactly
\begin{equation}
 z_a^{\otimes t}\otimes
 \bigotimes_{\ell=1}^h\left(\sum_j u_{aj}b_j^{\otimes q_\ell}\right).
 \label{fp:core-linear-identity}
\end{equation}
\emph{Proof.} Assign an independent summation index to each ordered child slot and use multilinearity slot by slot;
the sums over linear slots are exactly $z_a$, and nonlinear slots are unchanged. Fixed-degree symmetrization is an average over slot permutations
and can be applied after the summation. A purely unary chain is thus ordinary multiplication of normalized matrices, with no extra
$\sqrt{\text{width}}$ factor. All nonlinear factors still share $a$ and must not be replaced by independent arrays.
Coordinate masks do not commute with this summation in general; one must first show in the full physical lift that the error of deleting the mask is
$o(1)$, then restore the raw tensor and use this identity.\hfill$\square$

\subsection{Fresh coordinates and double raw marks}
\label{fp:core-fresh-section}

\paragraph{A probabilistic tool.}
If deterministic matrices $A_j:\mathbb R^u\to\mathbb R^v$ satisfy
$\max\{\|\sum_jA_j^*A_j\|,\allowbreak\|\sum_jA_jA_j^*\|\}\leq s^2$, then for independent
standard Gaussian coefficients
\[
 \mathbb P\!\left(\left\|\sum_jg_jA_j\right\|>t\right)
 \leq2(u+v)\exp[-t^2/(2s^2)].
\]
This is the finite-dimensional matrix Gaussian series bound. Its self-adjoint dilation is
$H_j=\left(\begin{smallmatrix}0&A_j^*\\ A_j&0\end{smallmatrix}\right)$,
$\|\sum_jH_j^2\|\leq s^2$; the matrix Gaussian Laplace bound
$\mathbb E\operatorname{tr}\exp(\lambda_{\rm c}\sum_jg_jH_j) \leq\operatorname{tr}\exp(\lambda_{\rm c}^2\sum_jH_j^2/2)$
with $\lambda_{\rm c}=t/s^2$ gives the upper tail, and repeating with the opposite sign gives the two-sided bound.
Below we only use the resulting $O(s\sqrt{\log d})$ bound when the dimensions are fixed powers of $d$.

\refstepcounter{theorem}\paragraph{Lemma~\thetheorem\ (fresh-color double marks).}
\label{fp:core-fresh}
Expand all unfrozen modes of a kernel through the original layer simultaneously. Among the new colors introduced by the Gaussian matrix of this layer
choose a mode $p$ and any different mode $q$. Assume the total number of modes is $N\geq3$.
On the bare kernel $E$, with arbitrarily high polynomial probability,
\begin{equation}
 \max_j\sup_{\|v\|\leq1}
 \|E(e_j\text{ in }p,v\text{ in }q,\cdot)\|_{\mathrm{HS}}
 \leq\rho_d.
 \label{fp:core-fresh-slice}
\end{equation}
The lower descendants of this layer are not yet expanded; frozen higher colors do not take part in the rotation of the new colors.

\emph{Proof.}
Conditionally on a kernel satisfying \eqref{fp:core-invariants}, set
$A_jv=E(e_j,v,\cdot)$. Summing matrix entries directly,
\[
 \left\|\sum_jA_j^*A_j\right\|
 =\|\operatorname{Flat}_{\{q\}|\{q\}^c}E\|^2\leq\delta_d^2,
 \qquad
 \left\|\sum_jA_jA_j^*\right\|
 =\|\operatorname{Flat}_{\{p,q\}|\{p,q\}^c}E\|^2\leq\delta_d^2.
\]
Nontriviality of the second bipartition uses exactly $N\geq3$. The series bound above and the
Gaussian tail bound $\|g\|\geq c\sqrt d$ show that, for a uniform vector on the sphere $u=g/\|g\|$,
the contracted norm is at most $C_M\delta_d\sqrt{\log d/d}$.
Take a union bound over the columns of an independent Haar matrix; the columns need not be independent.
Right-multiplying the Gaussian layer matrix by this orthogonal matrix preserves its law and rotates all new-color modes simultaneously.
The rotation is an isometry on unfixed modes; if $q$ is also a new color, the supremum over $v$ is unchanged.
Hence the estimate after the Haar rotation transfers back to the original coordinates. There are only finitely many kernels and modes.
\hfill$\square$

\refstepcounter{theorem}\paragraph{Corollary~\thetheorem\ (collision restoration).}
\label{fp:core-collision}
If a new color appears at least twice and there are at least three modes, its collision sector satisfies
\begin{equation}
 \|E_{\mathrm{coll}}\|_{\mathrm{HS}}^2
 \leq\sum_{p<q}\sum_j\|E(e_j\text{ in }p,e_j\text{ in }q,\cdot)\|^2
 \leq C d\rho_d^2=o(1).
 \label{fp:core-collision-bound}
\end{equation}
This tensor must be deleted when the new color is first introduced, and the error then passed through the full bounded lift.
The final error is then $o(1)$, with no dimension loss from summing over teacher indices afterwards.
Deterministic duplication of a parent index is not such a collision; it is always one shared label.

\paragraph{Dense Hilbert fibres.}
\label{fp:core-dense}
If a bare Gaussian polynomial kernel $E$ has polynomial degree at most a fixed $D$,
$\mathbb E\|E\|^2\leq C$, and the kept tuples of the chosen $r$ coordinates form a permutation orbit of size
$\geq c d^r$, then
$\mathbb E\|E_I\|^2\leq C'd^{-r}$: sum along the orbit and use distributional invariance.
The Gaussian Hilbert-valued hypercontractive inequality gives
\[
 (\mathbb E\|E_I\|^p)^{1/p}
 \leq(p-1)^{D/2}(\mathbb E\|E_I\|^2)^{1/2}.
\]
Taking $p=c_M\log d$, Markov's inequality and a union bound over polynomially many tuples give
$\max_I\|E_I\|\leq L_dd^{-r/2}$.
Several colors use the product of orbit sizes, but forced shared coordinates must be recounted along the actual orbit.
In particular, two distinct coordinates of the raw same-color coefficient $D_R$ satisfy
\begin{equation}
 \max_{a\ne b}\|D_R(e_a,e_b,\cdot)\|\leq\rho_d.
 \label{fp:core-distinct}
\end{equation}
The second-moment bound for bare kernels follows from fixed-degree expansion, fixed-order Gaussian moments and bounded full lifts;
alternatively condition layer by layer and use dimension-free fixed-order moments of Gaussian matrix norms.

\refstepcounter{theorem}\paragraph{Lemma~\thetheorem\ (double-parent pruning).}
\label{fp:core-two-parent}
A teacher leaf has $m$ physical slots touching at least two distinct parent modes of the raw coefficient $D_R$,
$R\geq3$. The correlation vector of each linear slot is
$\eta^{(p)}_{ia}=\langle a_i,w_a^{(p)}\rangle$,
$\|\eta_i^{(p)}\|_2\leq C$.
The full features of each parent mode and the nonempty residual features after deleting a specified linear slot are bounded operators.
If at least one of the two chosen parent modes has a residual, and \eqref{fp:core-distinct} and
\eqref{fp:core-fresh-slice} hold on the raw slices used, then the pairing of this configuration is
$O(\sqrt d\rho_d)+o(1)=o(1)$.
If one residual is kept and the other vanishes, the fixed coordinate must lie in a fresh mode, or another raw
rotation-invariant slice theorem must give the same directional estimate.

\emph{Proof.}
If both parent modes have residuals, first keep their orthogonal indices $a,b$. In the collision-free sector, after deleting
the two slots the source tensor has Hilbert norm at most
\[
 C\left(\sum_{a\ne b}|\eta^{(p)}_{ia}|^2|\eta^{(q)}_{ib}|^2
                 \|D_R(e_a,e_b,\cdot)\|^2\right)^{1/2}
 \leq C\rho_d.
\]
The mask with $(a,b)$ fixed is an orthogonal projection, and the lifts of the other parent modes are bounded.
If only parent mode $p$ keeps a residual, first restore the raw coefficient in the full physical lift, with error controlled by
\eqref{fp:core-collision-bound}; then contract the other mode to get
\[
 C\left(\sum_a|\eta^{(p)}_{ia}|^2
       \|D_R(e_a,\eta_i^{(q)},\cdot)\|^2\right)^{1/2}
 \leq C\rho_d.
\]
The contraction operator for the remaining slots of the teacher leaf is bounded. If the teacher coefficient is $C$, its full remaining lift
slice has norm at most $C_0\|C(e_i,\cdot)\|$. The final sum uses
$\sum_i\|C(e_i,\cdot)\|\leq C_1\sqrt d\|C\|$, which concludes.
Among several linear siblings of the same parent only the selected one may be deleted; the others belong to the nonempty residual.
If all touched parent modes are purely unary, compress exactly; if there is only one parent mode, the whole leaf merges, and
the lemma does not claim smallness.\hfill$\square$

\paragraph{Marks that do not generalize.}
For $v_a=d^{-1/2}$ and normalized Gaussian rows $u_a$, the kernel
$E_{a,j,k}=v_au_{aj}u_{ak}$ has bounded Hilbert norm and a dense maximal entry, yet
\[
 \left\|E\left(e_a,\frac{u_a}{\|u_a\|},\cdot\right)\right\|
 =d^{-1/2}\|u_a\|^2.
\]
Two visible descendants do not form two raw marks. With only two modes,
the second variance bound in \eqref{fp:core-fresh-slice} is also unavailable; for example
for $E_{ab}=d_aU_{ab}$ only $\|E\|_{\mathrm{op}}\leq C\|d\|_\infty$
is directly available, and one cannot claim that all coordinate--vector slices are $O(d^{-1})$.

\refstepcounter{theorem}\paragraph{Corollary~\thetheorem\ (expected control of raw higher-order collisions).}
\label{fp:core-raw-collision-expectation}
Let $C_R$ be the raw homogeneous coefficient of a fixed-depth, fixed-degree Gaussian polynomial network with zero constant term,
$R\geq3$, possibly a sum of finitely many profiles. Then for every fixed $p<\infty$
$\sup_d\mathbb E\|C_R\|_{\mathrm{HS}}^p<\infty$, and
$\mathbb E\|\operatorname{Coll}_R C_R\|_{\mathrm{HS}}^2=o(1)$.
For the moment bound, write each profile as a finite sum of raw tree lifts and slot permutations. Its norm is at most
a fixed constant times fixed powers of finitely many raw Gaussian operator norms; normalized Gaussian
matrices at fixed width ratios have dimension-free moments of all fixed orders. H\"older's inequality gives the
bound without independence of the factors across layers. For the raw same-color kernel, the first input Gaussian layer gives
orthogonal invariance; the fresh lemma and the collision corollary give, for any $M$, an event $\Omega_M$ whose complement
has probability $O(d^{-M})$ and on which the squared collision norm is at most $L_d/d$. Hence
\[
 \mathbb E\|\operatorname{Coll}_R C_R\|^2
 \leq L_d/d+
 (\mathbb E\|C_R\|^4)^{1/2}\mathbb P(\Omega_M^c)^{1/2}=o(1).
\]
Symmetrization involves finitely many permutations, and finite sums of profiles are handled by the triangle inequality.

\refstepcounter{theorem}\paragraph{Lemma~\thetheorem\ (singleton pinning for scalar-slot literals).}
\label{fp:core-scalar-literal}
Fix a configuration graph whose teacher vertices all have degree $m\geq3$ and whose source vertices have positive integer degrees.
The cross matrices of the physical slots come from bounded teacher and source single-slot maps and have bounded operator norm.
At least one specified literal slot lies in a source column of degree at least two, and its cross matrix
$G$ satisfies $\max|G|\leq\kappa$. All literal slots satisfy the same small-entry bound.
Then the full configuration operator $K$ has norm at most $C\sqrt\kappa$.
A column may contain several pure literals, and singleton columns may appear anywhere.

\emph{Proof.}
First work in ordered index spaces without distinctness constraints. Deleting all singleton source columns gives the kernel
$H$; let $P$ be the set of teacher positions touching these columns. The matrix $S$ of a deleted column
contributes the Gram $SS^*$ to $KK^*$; the entrywise product of these Grams at one teacher position is
$Q_p$. Since $SS^*\preceq CI$ with bounded diagonal, the Schur PSD order gives
$Q_p\preceq CI$. The remaining teacher positions carry all-ones Grams, also PSD.
Replacing these Gram factors one at a time gives
\[
 KK^*\preceq C\operatorname{Pin}_P(HH^*),\qquad
 \|K\|^2\leq C\max_{\mathbf p}\|H_{\mathbf p}\|^2,
\]
where pinning forces the two teacher tuples to agree on $P$, and $H_{\mathbf p}$ is the submatrix with these
labels fixed.
If a kept literal touches a fixed teacher position, its factor
$G_{p,a}$ is a diagonal multiplier on source indices of norm at most $\kappa$.
After deleting this slot the source column still has positive degree, and the remaining kernel is the adjoint of a composition of bounded source lifts, physical slot permutations, contractions with fixed teacher
vectors and the other teacher lifts, so $\|H_{\mathbf p}\|\leq C\kappa$.
Otherwise all kept literals touch unfixed positions. Each kept source column has degree at least two; with the teacher
tuple fixed, the row $\ell^2$ bounds and H\"older bound the absolute sum over that source index by
a constant, so absolute row sums are at most $C$. Each unfixed teacher vertex keeps all $m$ slots; with the source tuple fixed,
extract $\kappa$ at one teacher vertex touching a literal and sum the remaining $m-1\geq2$ factors
with the column $\ell^2$ bounds and H\"older, and likewise at the other teacher vertices. Hence absolute column sums
are at most $C\kappa$, and the Schur test gives $\|H_{\mathbf p}\|\leq C\sqrt\kappa$.
Distinct-index constraints are orthogonal projections on the domain and range and do not increase norms.\hfill$\square$

If all literals are singletons the lemma claims nothing, since pinning deletes
all of them; if the residual of a column is an arbitrary polynomial that does not factor into scalar single-slot factors, the lemma
does not apply either, and the full source-tensor lemmas below are used instead.

\subsection{Full source-tensor estimates for literals and polynomial residuals}
\label{fp:core-residual-section}

\refstepcounter{theorem}\paragraph{Lemma~\thetheorem\ (small diagonal factor for split hits).}
\label{fp:core-split-hit}
Let the source column be $\ell_a\otimes F_a$, with bounded full residual family $F:e_a\mapsto F_a$,
and the full lifts of the other source columns bounded and not depending additionally on $a$. The teacher leaf contains this literal and has
$t\geq1$ slots from other source columns. If the teacher row matrix is bounded and
$G_{ia}=\langle a_i,\ell_a\rangle$ has maximal entry at most $\kappa_d$, the configuration
operator has norm at most $C\kappa_d$. The physical slots of $F_a$ need not be separable linear factors.

\emph{Proof.}
Starting from the source coefficient space, keep $a$ as an orthogonal coordinate and lift only the other source columns.
Extract the specified $t$ physical slots from these columns with
$x\mapsto(\langle a_i^{\otimes t},x\rangle)_i$, keeping the output coordinate $i$.
This analysis operator is bounded since its adjoint is
$e_i\mapsto a_i^{\otimes t}=(A^*)^{\otimes t}\Delta_te_i$.
Now $(i,a)$ are both orthogonal coordinates, and multiplication by the diagonal function $G_{ia}$ has operator norm
exactly $\max|G_{ia}|\leq\kappa_d$.
Then apply the same bounded family map $F$ in each $i$ block, and contract the remaining slots of this leaf lying in $F_a$
with the corresponding $a_i$. This is a bounded operator, block diagonal in $i$, which keeps $i$ throughout.
This yields $i$ and all unconsumed physical source slots; the full analysis maps of the other teacher leaves remain bounded.
The result is an operator from full source coefficients to full teacher coefficients, with norm the product of the norms above,
$C\kappa_d$. No partial trace after summing labels occurs, and there is no loss from summing over $i$ one by one.
If $t=0$, no independent teacher coordinate can be produced before multiplying by $G$, and the proof does not apply.
\hfill$\square$

\refstepcounter{theorem}\paragraph{Lemma~\thetheorem\ (matricizations of true Gaussian residuals).}
\label{fp:core-residual-flat}
Let $F_a=\sum_bU_{ab}Q_b$, where $U$ is a Gaussian layer of the original network or a fixed consecutive product,
and $Q_b$ are positive-degree raw tree profiles, possibly with frozen literal leaves at different heights, each bounded
adaptive leaf map producing a single physical slot. Then $e_a\mapsto F_a$ has bounded operator norm, and
every nontrivial physical matricization of each $F_a$ of degree at least two has norm at most $\delta_d$,
simultaneously for all $a$.

\emph{Proof.}
For fixed $a$, rebuild the kernel of $Q$ with readout $U_{a,\cdot}$.
\eqref{fp:core-product-event} bounds the norms and coordinate maxima of all these readouts simultaneously;
the closure lemma bounds all nontrivial matricizations, and the final single-slot leaf maps preserve the bounds.
The family operator is the composition of $U^*$ with full tree lifts and leaf maps, hence bounded.
Finite residual products with the same index stay bounded through the diagonal embedding.
Independence of the leaf maps is not used, and residuals containing literals are not treated as new Gaussian networks.
\hfill$\square$

\refstepcounter{theorem}\paragraph{Lemma~\thetheorem\ (source-flattening repair with at least two modes).}
\label{fp:core-hybrid-Rge2}
Let the explicitly rebuilt source coefficient kernel $D\in H_1\otimes\cdots\otimes H_R$ satisfy
$R\geq2$, $\|D\|\leq C$ and
$\|\operatorname{Flat}_{1|\mathrm{rest}}D\|\leq\delta_d$.
The first-mode feature is
$H_a=\ell_a\otimes F_a^1\otimes\cdots\otimes F_a^s$,
all full family operators are bounded, and every nontrivial matricization of each $F_a^j$ is at most $\delta_d$.
The teacher leaf $a_i^{\otimes m}$ lies entirely in $H_a$, with one slot touching the literal, where
$G_{ia}=\langle a_i,\ell_a\rangle$ satisfies
$\max|G_{ia}|\leq\kappa_d$, $\sup_i\|G_{i,\cdot}\|_2\leq C$.
The remaining $m-1\geq1$ slots touch residuals. Other source modes are lifted by bounded maps not depending additionally on $a$.
Then for any teacher coefficient of bounded Hilbert norm the pairing of this configuration is
\begin{equation}
 O\!\left(\sqrt d\,\delta_d(\delta_d+\kappa_d)\right)=o(1),
 \qquad \kappa_d\leq L_d/\sqrt d.
 \label{fp:core-hybrid-bound}
\end{equation}

\emph{Proof.}
Fix $i$, contract all slots of the teacher leaf, and keep all uncontracted physical source slots.
Call the resulting first-mode operator $B_i$. If some residual factor is partly consumed and partly outside,
its nontrivial matricization gives $\|B_ie_a\|\leq C\delta_d|G_{ia}|$, so
$\|B_i\|_{\mathrm{HS}}\leq C\delta_d$.
Otherwise every residual factor is fully consumed or fully outside, and at least one is fully consumed.
Contracting that factor gives a vector $\eta_a$ with $\|\eta\|_2\leq C$ by boundedness of its full family operator.
The other factors have bounded row norms, so
$\|B_i\|_{\mathrm{HS}}\leq C\kappa_d$.
The two cases exhaust all assignments.

Let $T=\bigotimes_{p=2}^RT_p$; the full source residual is $S_i=(B_i\otimes T)D$.
In the mode-one matricization, the Hilbert--Schmidt product inequality gives
\[
 \|S_i\|_{\mathrm{HS}}
 \leq\|T\|\,\|B_i\|_{\mathrm{HS}}
       \|\operatorname{Flat}_{1|\mathrm{rest}}D\|_{\mathrm{op}}
 \leq C\delta_d(\delta_d+\kappa_d).
\]
All remaining source slots now lie in the Hilbert space of $S_i$; take the inner product with the remaining teacher slice directly
and use $\sum_i\|C(e_i,\cdot)\|\leq C\sqrt d\|C\|$ to obtain
\eqref{fp:core-hybrid-bound}. No final partial trace is taken.
\hfill$\square$

For removable collisions with $R\geq3$, restoration must be done first in the full lift. If an off-diagonal restriction is imposed on the raw coefficient with $R=2$,
its diagonal part need not have small Hilbert norm, but
$\|\operatorname{diag}D\|_{\mathrm{op}}\leq\|D\|_{\mathrm{op}}$,
$\|\operatorname{diag}D\|_{\mathrm{HS}}\leq\|D\|_{\mathrm{HS}}$,
so the same lemma applies separately to the raw matrix and its diagonal. Other masks cannot be handled this way.

\refstepcounter{theorem}\paragraph{Lemma~\thetheorem\ (one-mode product residuals).}
\label{fp:core-hybrid-R1}
Let $S=\sum_ad_a\ell_a\otimes\bigotimes_{j=1}^{h}F_a^j$, $h\geq2$,
$\|d\|_\infty\leq\kappa_d$, with the family, matricization and $G$ assumptions of the previous lemma.
The teacher leaf contains the indicated literal and consumes at least one non-literal slot.
Then the full remaining source tensor satisfies
\begin{equation}
 \sup_i\|S_i\|\leq C(\kappa_d\delta_d+\kappa_d^2),
 \label{fp:core-R1-bound}
\end{equation}
so the full pairing is at most $C\sqrt d(\kappa_d\delta_d+\kappa_d^2)\|C\|=o(1)$.

\emph{Proof.}
Divide the $F^j$ into fully consumed, partly consumed and fully outside factors.
Fully consumed factors give $\ell^2$-bounded scalar vectors $\eta$; partly consumed ones give vector families with row norms at most
$C\delta_d$ and bounded family operators; fully outside ones keep their bounded families.
If some factor is partly consumed and at least two factors keep outside slots, then in the Gram of the residual product one factor
has diagonal at most $C\delta_d^2$ and another satisfies $\preceq CI$.
The Schur PSD order bounds the norm of the product family by $C\delta_d$.
Its coefficient $d_aG_{ia}$ times the remaining bounded scalars has $\ell^2$ norm at most $C\kappa_d$,
which gives $C\kappa_d\delta_d$.
If only one factor is outside and it is partly consumed, $h\geq2$ guarantees a fully consumed factor $\eta$, and
the triangle inequality and Cauchy--Schwarz give
$\|S_i\|\leq C\kappa_d\delta_d\sum_a|G_{ia}\eta_a| \leq C\kappa_d\delta_d$.
If nothing is partly consumed but outside slots remain, at least one factor is fully consumed, and its $\eta$ together with
$\|d\|_\infty\max|G|$ bounds the $\ell^2$ norm of the coefficients by $C\kappa_d^2$.
If everything is consumed, the product of at least two fully consumed vectors is $\ell^1$-bounded, so the scalar is also at most
$C\kappa_d^2$. The four cases are exhaustive.\hfill$\square$

\paragraph{Permanent root index and the exact boundary case.}
\label{fp:core-spectator}
If the one-mode exception has a single residual factor, the previous lemma does not apply.
Keep the original root index $a$ permanently, start from $E_{ab}=d_aU_{ab}$ and expand only descendant modes.
By the closure lemma every rebuilt snapshot has at least two modes and all nontrivial matricizations are
$O(\delta_d)$; the root literal acts as a bounded external single-slot map.
Hence descendant hybrids with a separate full feature representation are handled by
Lemma~\ref{fp:core-hybrid-Rge2}. When two fresh active modes are present, the permanent root index
ensures at least three modes, and Lemma~\ref{fp:core-fresh} applies.
This does not license deleting or tracing out the root index.

\paragraph{Legitimate amplification by Hilbert spectators.}
\label{fp:core-amplification}
If the scalar configuration operator $K:\ell^2(J)\to\ell^2(I)$ has a norm bound and all remaining physical slots
are rebuilt through $S_I:\mathcal H\to\mathcal Q,T_J:\mathcal K\to\mathcal Q$,
$\sup_I\|S_I\|,\sup_J\|T_J\|\leq B$, then
\[
 \left|\sum_{I,J}K_{IJ}\langle S_Ic_I,T_Jb_J\rangle\right|
 \leq B^2\|K\|\left(\sum_I\|c_I\|^2\right)^{1/2}
                   \left(\sum_J\|b_J\|^2\right)^{1/2}.
\]
The proof is the full operator factorization
$\operatorname{diag}(S_I)^*(K\otimes I_{\mathcal Q})\operatorname{diag}(T_J)$.
The index dependence must be strictly separated; factors depending on both $I,J$ cannot be hidden on either side.

\paragraph{Low-order cases.}
Two-mode double-parent states cannot use the
$d^{-1}$ mark bound, and one-mode states with a single residual factor cannot use the one-mode product lemma.
In particular, knowing only that a consumed row kernel $T$ is small is not enough: take
$C_{ir}=\delta_{ir}/\sqrt d,T_{ib}=\delta_{ib}/\sqrt d,F_{br}=\delta_{br}$;
then $\|C\|_{\mathrm{HS}}=1$, $\|C\|_{\mathrm{op}}=\|T\|_{\mathrm{op}}=d^{-1/2}$,
$\|F\|=1$, but $\sum_{i,b,r}C_{ir}T_{ib}F_{br}=1$.
Partial traces taken after the raw coefficient is lost can thus recover dimension factors; the full source-tensor proofs above avoid
this operation. The factorization of the full history is constructed step by step in Section~\ref{fp:history}.

\section{Fixed input chaos of actual activations and closure of raw positive-degree polynomials}
\label{fp:actual-section}\label{app:closure}

This section studies value errors under a single network prior. Let $k_0=d$, with $L$ fixed, and
$c d\leq k_\ell\leq C d$. Write $A_\ell=W^{(\ell)}/\sqrt{k_{\ell-1}}$;
the entries of $W^{(\ell)}$ are independent $N(0,1)$, independent across layers, and $X\sim N(0,I_d)$ is independent of the weights.
Define
\[
 h^{(0)}=X,\quad h^{(\ell)}=\sigma(A_\ell h^{(\ell-1)}),\quad
 q_\ell=k_\ell^{-1}\|h^{(\ell)}\|^2,\qquad
 F_L:=F_{\theta(W)}^{(L)}=k_L^{-1/2}\sum_i h_i^{(L)}.
\]
The activation satisfies $\mathbb E\sigma(G)=0$, $\mathbb E\sigma(G)^2=1$, and is one of the following:
bounded smooth odd (every fixed-order derivative bounded), or
\begin{equation}\label{fp:actual-pl}
 \sigma(z)=a z+b(|z|-u),\qquad u=\sqrt{2/\pi},\qquad
 a^2+b^2(1-u^2)=1.
\end{equation}
Write $\psi_j=\mathrm{He}_j/\sqrt{j!}$,
$a_j=\mathbb E[\sigma(G)\psi_j(G)]$,
$b_j=\mathbb E[(\sigma(G)^2-1)\psi_j(G)]$.
$\Pi_N$ is the orthogonal projection onto the input chaos of order $N$ of $X$; for symmetric tensors
\begin{equation}\label{fp:actual-isometry}
 \mathbb E_X I_N(T)I_N(S)=N!\langle T,S\rangle.
\end{equation}
In particular, ``take the coefficient of formal degree $N$ and apply $I_N$'' below is not the same as computing the whole
raw polynomial function and applying $\Pi_N$; the latter contains contractions of higher formal degree.

\paragraph{Local moment estimates.}
The following estimates involve only single-network marginals.
They follow from Lemma~\ref{an:moments}, the conditional row estimate~\eqref{an:scalar}
and the Gaussian net tail bound in Lemma~\ref{an:derivatives}.
For a fixed number of layers and fixed $p<\infty$,
\begin{equation}\label{fp:actual-local-interface}
 \|q_\ell-1\|_{L^p(W,X)}\leq C_p d^{-1/2},\quad
 \sup_{d,\ell,i}\mathbb E|h_i^{(\ell)}|^p<\infty,\quad
 \sup_d\mathbb E\|A_\ell\|_{\rm op}^p<\infty.
\end{equation}
The same fixed product-moment bounds hold for any fixed number of actual features, with repeated labels allowed.
By the first bound with arbitrarily large $p$, the probability of $q_\ell\notin[1/2,2]$ decays faster than any polynomial.
The first-layer unit directions $v_i=A_{1,i:}/\|A_{1,i:}\|$ and radii $s_i=\|A_{1,i:}\|$
satisfy simultaneously
\begin{equation}\label{fp:actual-row-event}
 \max_i|s_i-1|+\max_{i\ne j}|\langle v_i,v_j\rangle|
 \leq C_M\sqrt{\log d/d},\qquad \|A_1\|_{\rm op}\leq C_M
\end{equation}
with failure probability at most $d^{-M}$. These also involve only a single network prior.
From these local bounds we derive the truncation, collective-mean and finite restoration steps; raw-coefficient
higher-order collision bounds are stated where they are used.

\subsection{Full input projection of finitely many actual ridge functions}
\label{fp:actual-finite-support}
Fix $t,N$, let $v_1,\ldots,v_t$ be unit vectors, $G=(\langle v_i,v_j\rangle)$,
$\|G-I\|\leq\eta$, $\max_i|s_i-1|\leq\delta$. Then
\begin{equation}\label{fp:actual-ridge-projection}
 \left\|\Pi_N\prod_{i=1}^t\sigma(s_iv_i\cdot X)
 -\sum_{\substack{n_i\geq1\\\sum_i n_i=N}}
 \frac{\prod_i a_{n_i}}{\sqrt{\prod_i n_i!}}
 I_N\!\left(\operatorname{Sym}\bigotimes_i(s_iv_i)^{\otimes n_i}\right)
 \right\|_{L^2_X}\leq C_{t,N}(\eta+\delta).
\end{equation}
For $t>N$ the approximating sum on the right is empty.

Proof. All tensors involved lie in $\operatorname{Sym}^N\operatorname{span}\{v_i\}$.
Use as a basis the tensors with $n_i\geq0$, $\sum n_i=N$,
$B_n=I_N(\operatorname{Sym}\bigotimes_i v_i^{\otimes n_i})$.
Their Gram matrix is a polynomial in $G$ and uniformly invertible near $G=I$: the linear map from an orthonormal basis
to $v_i$ has uniformly bounded largest and smallest singular values, and so do its fixed symmetric tensor powers.
The coefficient vector of the projection is the inverse Gram times
$\mathbb E[\prod_i\sigma(s_iZ_i) B_n(Z;G)]$ with $Z\sim N(0,G)$.
On a fixed compact neighbourhood of $G$, differentiating the Gaussian density in the covariance gives uniformly integrable
polynomials times the Gaussian density; the coefficients of $B_n$ also depend polynomially on $G$.
The activations grow at most linearly, so these finitely many integrals are locally Lipschitz in $G$.
Lipschitz dependence on $s_i$ follows from Lipschitz continuity of $\sigma$ and Gaussian weighted moments.
At $G=I,s_i=1$, independence gives the coefficients $\prod_i a_{n_i}/\sqrt{\prod_i n_i!}$;
$a_0=0$ removes all zero degrees. Finally, replacing $v_i^{\otimes n_i}$ by
$(s_iv_i)^{\otimes n_i}$ costs only $O(\delta)$. Equivalence of norms in finite dimension proves
\eqref{fp:actual-ridge-projection}. The proof integrates the full non-polynomial product
and therefore includes all lower-order contractions from high leaf Hermite degrees and non-orthogonal source directions.

\subsection{Finite reduction of the actual positive-root expansion: truncation order and errors}
\label{fp:actual-positive}
Given previous-layer features $h\in\mathbb R^k$, $q=\|h\|^2/k$, the current row preactivation is
$W_i\cdot h/\sqrt k$. For $q>0$ write
\[
 \sigma(\sqrt q G)=\sum_{r\geq0}c_r(q)\psi_r(G),\qquad
 c_r(q)=\mathbb E\sigma(\sqrt qG)\psi_r(G).
\]
The multivariate Hermite addition formula is
\begin{equation}\label{fp:actual-multivariate}
 \psi_r(W_i\cdot h/\|h\|)=
 \sum_{|\alpha|=r}\sqrt{\frac{r!}{\alpha!}}
 (h/\|h\|)^\alpha\prod_j\psi_{\alpha_j}(W_{ij}).
\end{equation}
$q=0$ is taken as the value of the continuous actual random variable; no moments of $q^{-1}$ are introduced.

For odd functions the conditional mean vanishes identically, and $q\mapsto\sigma(\sqrt q\,\cdot)$ is
$L^2(\gamma)$-continuous on any compact $q$ interval; tails of orthogonal projections vanish uniformly on compact sets.
Bounded activations above the first layer give $q\leq\|\sigma\|_\infty^2$.
For the piecewise-linear branch, $c_r(q)=\sqrt q\,a_r$ for $r\geq1$, and the tail sum is
$q\sum_{r>R}a_r^2$, integrable with the local moment bounds.
For $\chi=\sigma^2-1$, which appears later, the positive coefficients have the form
$c_r^\chi(q)=q A_r+\sqrt q B_r$ with
$\sum_r(A_r^2+B_r^2)<\infty$, since the two terms are Hermite coefficients of fixed scalar functions of quadratic and linear growth.
Their tails are also controlled by fixed $q$ moments.

These single-row facts suffice for products of a fixed number of distinct parent rows, without requiring single-row tails to vanish in arbitrarily
high $L^p$. Conditionally on the previous layer, distinct parent rows are independent. Writing the product difference factor by factor,
the conditional second moment of each term equals the second moment of one tail times the conditional second moments of the other factors;
truncated factors have $L^2$ norm at most that of the full factor. Finitely many Cauchy--Schwarz steps and
\eqref{fp:actual-local-interface} control the total error. Truncation orders are chosen from root to leaf,
and each new truncation need only adapt to the previously fixed number of parent rows.

With the kept $r$ fixed, replace the radial factor $c_r(q)q^{-r/2}$ by $a_r$.
On $q\in[1/2,2]$ this factor is smooth with bounded fixed-order derivatives, so the error is controlled by
$\|q-1\|_{L^p}=O(d^{-1/2})$. On the complement, for the exact finite component condition first on $h$
and use fixed-order Gaussian Hermite moment bounds, whose norms contain $c_r(q)$ and no inverse variance;
the replaced finite polynomial is controlled by fixed moments of the actual features. Arbitrarily high local moments and H\"older
make the complementary error vanish. The factor $A_rq^{1-r/2}+B_rq^{(1-r)/2}$ of $\chi$
is handled identically; this uses compact-interval estimates and moments of actual random variables, not negative moments of $q$.

Next delete repeated child labels at each interface; \emph{the coefficients are still normal-ordered Hermite and the child features are still actual functions}.
With a fixed number $T$ of child slots there are at most $C_T k^{T-1}$ collision assignments, with squared normalization
$k^{-T}$. Multi-index Hermite orthogonality in the current weights and fixed product moments of repeated actual features
give a squared expected collision error $O_T(k^{-1})$.
Orderings describing the same edge set have only multiplicities depending on the fixed $T$ and are merged before orthogonality is used.
Only after collisions are deleted does each weight variable appear once, so normal-ordered coefficients become plain products of edge variables.
For the piecewise-linear branch, the zero-degree term appearing \emph{inside} fixed positive-root products,
$c_0(q)=bu(\sqrt q-1)$, has every fixed $L^p$ norm $O(d^{-1/2})$,
so it is removed by the same product estimate; the zero-degree term at the top root is kept separately.

To see that these local errors are not amplified by the width, write the whole error at a given stage as
\begin{equation}\label{fp:actual-phase}
 \Delta=\sum_g\beta_g\Psi_g(W_{>\ell}) E_g(W_{\leq\ell},X),
 \qquad \sum_g\beta_g^2\leq C_{\rm fixed}.
\end{equation}
Here $\Psi_g$ are products over the distinct square-free edge sets already processed; descriptions of the same edge set
are merged first. Each edge carries the power $-1/2$ of an input width and the root readout carries $k_L^{-1/2}$.
Summing over the root label and each free downward child label cancels these squared normalizations exactly,
leaving only combinatorial constants of the fixed graphs, degrees and orderings, which proves the second inequality.
$\Psi_g$ depends only on higher independent Gaussian layers, so
\begin{equation}\label{fp:actual-phase-bound}
 \mathbb E\|\Delta\|_{L^2_X}^2
 \leq C_{\rm fixed}\sup_g\mathbb E\|E_g\|_{L^2_X}^2.
\end{equation}
If errors share a higher phase, the whole lower-layer product is taken as $E_g$, so shared phases are not mistaken
for independent ones. Thus, after fixing $\ecut>0$ and choosing finite root-to-leaf truncations successively,
we obtain a finite sum $F_{+,\ecut}$ of graph values with distinct labels at every interface, and
\begin{equation}\label{fp:actual-finite-error}
 \limsup_d\mathbb E\|F_+-F_{+,\ecut}\|_{L^2_X}^2
 \leq C_L\ecut^2.
\end{equation}
All normalizations, radial replacements and collision errors are controlled under single-network marginals.

Finally keep the actual first-layer ridge functions and apply \eqref{fp:actual-ridge-projection}.
For each fixed support size, \eqref{fp:actual-row-event} makes the error vanish; aggregate with
\eqref{fp:actual-phase-bound} rather than the triangle inequality over all supports.
On the complement of the first-layer good event, fixed-degree polynomial moments and actual fixed product moments grow at most polynomially in the dimension,
and a large enough failure exponent makes their contribution vanish.
If a node of a graph has degree larger than $N$, it has more than $N$ distinct subtrees; each subtree has positive-degree descendants,
so there are more than $N$ distinct first-layer leaves, and the leading term vanishes by \eqref{fp:actual-ridge-projection}.
Hence once all truncations are at least $N$, the set of nonempty leading graphs no longer changes.
With truncations fixed, letting first $d\to\infty$ and then $\ecut\to0$ gives \emph{one fixed finite family},
not a final expression that grows with the precision.

\subsection{Keeping the top-root mean: exact recursions for $\mu$ and the centered square}
\label{fp:actual-mean}
For \eqref{fp:actual-pl}, direct expansion gives
\begin{align}
 \mu(q)&=\mathbb E\sigma(\sqrt qG)=bu(\sqrt q-1),\\
 \nu(q)&=\mathbb E\sigma(\sqrt qG)^2
 =(a^2+b^2)q-2b^2u^2\sqrt q+b^2u^2,
 \qquad \nu(1)=\nu'(1)=1.
\end{align}
Let $\chi(z)=\sigma(z)^2-1$, $\mu'=\mu'(1)=bu/2$. The identity
\begin{equation}\label{fp:actual-sqrt}
 \sqrt q-1-\tfrac12(q-1)=-\tfrac12(\sqrt q-1)^2,
 \qquad |\sqrt q-1|\leq |q-1|\quad(q\geq0)
\end{equation}
and $\mathbb E|q_\ell-1|^4=O(d^{-2})$ give
\begin{align}
 \sqrt{k_L}\mu(q_{L-1})
 &=\mu'\sqrt{k_L/k_{L-1}}F_{L-1}^{\chi}+e_{L,d},
 &\mathbb E\|e_{L,d}\|_2^2&=O(d^{-1}),\label{fp:actual-mean1}\\
 \sqrt{k_t}\{\nu(q_{t-1})-1\}
 &=\sqrt{k_t/k_{t-1}}F_{t-1}^{\chi}+e_{t,d},
 &\mathbb E\|e_{t,d}\|_2^2&=O(d^{-1}).\label{fp:actual-mean2}
\end{align}
Here $F_t^\chi$ is shorthand for depth $t$ with the last output activation replaced by $\chi$, i.e. $F_t^\chi=k_t^{-1/2}\sum_i\chi((A_th^{(t-1)})_i) =\sqrt{k_t}(q_t-1)$.
In \eqref{fp:actual-mean2} the right-hand side for $t=1$ is read as
$\sqrt{k_1}(q_0-1)$, $q_0=\|X\|^2/d$.
Marking by $+$ the positive Hermite-order part in the root weights and iterating over layers gives
\begin{equation}\label{fp:actual-mean-telescope}
 F_L=F_{L,+}^{\sigma}
 +\mu'\sum_{t=1}^{L-1}\sqrt{k_L/k_t}\,F_{t,+}^{\chi}
 +\mu'\sqrt{k_L}(q_0-1)+o_{L^2(W,X)}(1).
\end{equation}
The input endpoint before the last term generally has nonzero limiting norm and must be kept; it lies exactly in the second input chaos,
so it is cancelled exactly by $\Pi_N$ only when $N>2$.
On the odd branch all conditional means vanish; set $\mu'=0$, and no centered-square recursion is needed.

\subsection{From the actual graph to masked fixed formal coefficients}
\label{fp:actual-masked-formula}
Define the fixed zero-constant polynomials
\begin{equation}\label{fp:actual-PQ}
 P_N(z)=\sum_{j=1}^N\frac{a_j}{\sqrt{j!}}z^j,
 \qquad Q_N(z)=\sum_{j=1}^N\frac{b_j}{\sqrt{j!}}z^j.
\end{equation}
Given the depth $t$, the root polynomial $A$ and lower zero-constant polynomials $p_1,\ldots,p_{t-1}$,
compose a formal polynomial network with the raw matrices $A_1,\ldots,A_t$ and the readout $k_t^{-1/2}\mathbf1$.
Denote its formal symmetric coefficient of degree $N$ by $C_{t,N}[A;p]$, and define
$R_{t,N}[A;p]=I_N(C_{t,N}[A;p])$.
If, in the root-to-leaf expansion, all tensor coefficients with repeated coordinates are deleted at every formal neuron interface,
we obtain $J_{t,N}[A;p]$. No square-free mask is imposed on the physical input coordinates of the first layer.
At depth one, $J_{1,N}=R_{1,N}$.
Different node degrees can be recorded by finitely many auxiliary marker coefficients, so this notation also covers finite linear combinations of fixed graphs.

The two preceding subsections show that, for every fixed $N>2$,
\begin{equation}\label{fp:actual-J-closure}
 \Pi_N F_L=J_{L,N}[P_N;P_N]
 +\mu'\sum_{t=1}^{L-1}\sqrt{k_L/k_t}\,J_{t,N}[Q_N;P_N]
 +o_{L^2(W,X)}(1).
\end{equation}
Here $;P_N$ abbreviates using $P_N$ at all lower layers.
The only depth-one endpoint without a higher-layer phase must be checked separately. The functions actually used are
$f=\sigma$ or $f=\chi=\sigma^2-1$, both of at most quadratic growth; their Gaussian coefficients
$\mathbb E[f(sG)\psi_N(G)]$ are locally Lipschitz in $s$ near one.
For $\sigma$ this follows from scalar Lipschitz continuity and Gaussian-weighted moments.
For $\chi$ on the piecewise-linear branch, set $v(x)=ax+b|x|$; then for $s>0$ there is the explicit identity
\[
 \chi(sx)=s^2v(x)^2-2bu\,s v(x)+b^2u^2-1.
\]
Hence for $s,t\in[1/2,2]$,
$|\chi(sx)-\chi(tx)|\leq C|s-t|(|x|^2+|x|)$; after multiplying by
$|\psi_N(x)|$ it is still dominated by an integrable Gaussian polynomial, giving the required local Lipschitz property.
On the odd branch the depth-one term uses only $\sigma$ and produces no collective $\chi$ root.
Replacing by $b_Ns^N$ (or $a_Ns^N$) gives a uniform $O(\max|s_i-1|)$ coefficient error.
The map $e_i\mapsto v_i^{\otimes N}$ has uniformly bounded operator norm: when $N\geq2$,
the absolute row sums of its Gram matrix are at most $\sum_j|\langle v_i,v_j\rangle|^2\leq\|V\|_{\rm op}^2$.
The normalized readout vector has unit norm, so the error of the shallow ridge sum also vanishes.
This completes the proof for all depth-one terms in \eqref{fp:actual-J-closure}.

\subsection{Quadratic diagonal correction and finite restoration}
\label{fp:actual-quadratic}
By the raw positive-degree coefficient collision corollary~\ref{fp:core-raw-collision-expectation},
we may use the following: for any fixed zero-constant polynomial network,
including fixed auxiliary marker components, its full raw formal coefficient of degree $r\geq3$, $D_r$, satisfies
\begin{equation}\label{fp:actual-raw-interface}
 \mathbb E\|\operatorname{Coll}_rD_r\|_{\rm HS}^2=o(1),
 \qquad \sup_d\mathbb E\|D_r\|_{\rm HS}^p<\infty\quad(p<\infty).
\end{equation}
where $\operatorname{Coll}_r$ keeps the ordered tensor terms with at least two equal coordinates.
This statement concerns \emph{full raw} coefficients; it must not be applied to masked or adaptively selected tensors.
Concretely, Lemma~\ref{fp:core-fresh} gives
$\max_i\sup_{\|v\|\leq1}\|D_r(e_i,v,\cdot)\|_{\rm HS} \leq (\log d)^C/d$, so
\[
 \|\operatorname{Coll}_rD_r\|_{\rm HS}^2
 \leq \binom{r}{2}\sum_i\|D_r(e_i,e_i,\cdot)\|_{\rm HS}^2
 \leq C(\log d)^{2C}/d.
\]
Off the high-probability event, fixed moments give \eqref{fp:actual-raw-interface}.
This does not extend to $r=2$; we compute this exception directly.

Consider an upper polynomial network of depth $t$ with input width $k$ and layer polynomials
$p_j(z)=\sum_{r\geq1}g_{j,r}z^r$. Its full quadratic coefficient $D_2$ satisfies
\begin{equation}\label{fp:actual-diag}
 \mathbb E\left\|\operatorname{Diag}D_2-
 \lambda I_k/\sqrt k\right\|_{\rm HS}^2=O(k^{-1}),\qquad
 \lambda=g_{t,2}\prod_{j=1}^{t-1}g_{j,1}^{2}\sqrt{k_t/k}.
\end{equation}
Proof: a quadratic term of a positive-degree network has exactly one quadratic node; all others are linear.
If the quadratic node is in layer $j$, the matrix after removing fixed scalars is
\[
 D_2^{(j)}=B_j^T\operatorname{diag}(v_j)B_j,\quad
 B_j=A_j\cdots A_1,\quad v_j=A_{j+1}^T\cdots A_t^T\mathbf1/\sqrt{k_t}.
\]
Write $B_j=C_jA_1$, $H_j=C_j^T\operatorname{diag}(v_j)C_j$.
The columns of $A_1$ are independent $N(0,I/k)$, so conditionally on $H_j$,
\begin{equation}\label{fp:actual-diag-gauss}
 \mathbb E(D_2^{(j)})_{ii}=\operatorname{tr}(H_j)/k,
 \qquad \mathbb E\sum_i|(D_2^{(j)})_{ii}-\operatorname{tr}(H_j)/k|^2
 =2\|H_j\|_{\rm HS}^2/k.
\end{equation}
By $\|H_j\|_{\rm HS}\leq\|C_j\|_{\rm op}^2\|v_j\|$ and the moment bounds, the latter is $O(k^{-1})$.
When $j<t$, $v_j$ is independent of $C_j$ with mean zero and covariance $I/k_j$, so
\[
 \mathbb E|\operatorname{tr}H_j|^2
 =k_j^{-1}\mathbb E\sum_a\|C_{j,a:}\|^4=O(1).
\]
the squared norm of its scalar diagonal part is also $O(k^{-1})$.
When $j=t$, $\operatorname{tr}H_t=\|C_t\|_F^2/\sqrt{k_t}$ with expectation $\sqrt{k_t}$.
The $\|C\|_F^2$ variance of a product of a fixed number of normalized Gaussian matrices is $O(1)$: adding a left factor gives
\[
 \mathbb E(\|AC\|_F^2\mid C)=\frac{k_{\rm new}}{k_{\rm old}}\|C\|_F^2,
 \quad \operatorname{Var}(\|AC\|_F^2\mid C)
 =\frac{2k_{\rm new}}{k_{\rm old}^2}\|CC^T\|_{\rm HS}^2=O_{L^1}(1),
\]
and the law of total variance gives the recursion. Restoring the fixed scalars, only a top-layer quadratic node contributes a deterministic main diagonal,
which is exactly \eqref{fp:actual-diag}. Finitely many auxiliary marker coefficients do not change the proof.

Substituting the diagonal main term at the interface $s$ into the lower output $y_i=p_s(z_i)$ gives exactly
\begin{equation}\label{fp:actual-fusion}
 \lambda k_s^{-1/2}\sum_i p_s(z_i)^2.
\end{equation}
So the correction is a positive-degree network with one fewer upper segment and top polynomial $p_s^2$, with a minus sign.
The two factors of a shared label must be multiplied first; they cannot be treated as independent sources.
The global distinctness masks of the other descendants are unchanged.
At the first layer this includes the possibly nonvanishing
$\lambda k_1^{-1/2}\sum_i A_{1,i:}^{\otimes(n_1+n_2)}$, which is exactly a shallow raw ridge coefficient.

Errors propagate through later layers: for every positive integer $q$,
$e_i\mapsto A_{\ell,i:}^{\otimes q}$ is
the restriction of $(A_\ell^T)^{\otimes q}$ to the coordinate-diagonal subspace, with norm at most
$\|A_\ell\|_{\rm op}^q$; coordinate masks and physical symmetrization are contractions.
At fixed order and fixed depth there are finitely many such operations, so by moment bounds and H\"older
the small errors in \eqref{fp:actual-raw-interface}, \eqref{fp:actual-diag} remain $o(1)$.
Error moments of order higher than two, if needed, follow from second moments by moment equivalence for Gaussian polynomials of fixed degree.

\subsection{Telescoping cancellation for simultaneous restoration of two families, and final closure}
\label{fp:actual-telescope}
Let the root polynomial be $A(z)=\sum g_{t,r}z^r$ and the lower layers $p_j$. For fixed $N>2$,
\begin{equation}\label{fp:actual-restoration}
 J_{t,N}[A;p]=R_{t,N}[A;p]
 -g_{t,2}\sqrt{k_t/k_{t-1}}\,
 R_{t-1,N}[p_{t-1}^2;p_{<t-1}]+o_{L^2(W,X)}(1)
 \quad(t\geq2).
\end{equation}
First restore the highest interface. $r=1$ has no collisions, $r\geq3$ uses \eqref{fp:actual-raw-interface},
and $r=2$ uses \eqref{fp:actual-diag}, producing the second family on the right; both families still carry lower masks.
At any lower interface $s<t-1$, the deterministic quadratic diagonal multiplier of the main family is
\[
 g_{t,2}\prod_{j=s+1}^{t-1}g_{j,1}^2\sqrt{k_t/k_s}.
\]
The second family contains the external multiplier $g_{t,2}\sqrt{k_t/k_{t-1}}$, and its root quadratic coefficient is
$[p_{t-1}^2]_2=g_{t-1,1}^2$, so the resulting multiplier is
\[
 g_{t,2}\sqrt{k_t/k_{t-1}}g_{t-1,1}^2
 \prod_{j=s+1}^{t-2}g_{j,1}^2\sqrt{k_{t-1}/k_s},
\]
equal to that of the main family with opposite sign. Both leave the same lower square-root polynomial and the same descendant mask,
so they cancel exactly. Other terms whose quadratic node is not at the current upper root have only vanishing diagonal errors.
Repeating over interfaces, all higher collisions are $o(1)$, there are a fixed number of them, and finally multiplying by
the fixed factor $N!$ in \eqref{fp:actual-isometry} proves \eqref{fp:actual-restoration}.
At each step the argument applies the collision theorem to the full raw coefficients of each family; it does not assume that masks preserve rotation invariance.

For any $f$ of at most polynomial growth, differentiating the Gaussian density in the variance gives
\[
 \left.\frac{d}{dq}\mathbb E f(\sqrt qG)\right|_{q=1}
 =\frac12\mathbb E[f(G)(G^2-1)].
\]
Hence the piecewise-linear branch satisfies
\begin{equation}\label{fp:actual-coeff-identities}
 a_2/\sqrt2=\mu',\qquad b_2/\sqrt2=\nu'(1)=1.
\end{equation}
Apply \eqref{fp:actual-restoration} to each family in \eqref{fp:actual-J-closure}.
The main depth-$L$ family produces at depth $L-1$ a root with coefficient
$-\mu'\sqrt{k_L/k_{L-1}}$ of type $P_N^2$; the collective family of depth $t\geq2$
produces at depth $t-1$ a root of the same kind with
$-\mu'\sqrt{k_L/k_t}\sqrt{k_t/k_{t-1}} =-\mu'\sqrt{k_L/k_{t-1}}$.
Each depth $1,\ldots,L-1$ therefore has exactly one square-root correction, which merges with the existing $Q_N$ root of that depth.
Let
\begin{equation}\label{fp:actual-U}
 U_N=[Q_N-P_N^2]_{\deg\leq N},\qquad
 \mathcal Q_N=R_{L,N}[P_N;P_N]
 +\mu'\sum_{t=1}^{L-1}\sqrt{k_L/k_t}\,R_{t,N}[U_N;P_N].
\end{equation}
No polynomial has a constant term, so higher formal degrees cannot be lowered by positive-degree substitution, and we may truncate at $N$.
By the error ordering and the finite restoration proof above,
\begin{equation}\label{fp:actual-final}
 \boxed{\displaystyle
 \mathbb E_W\|\Pi_NF_L-\mathcal Q_N\|_{L^2_X}^2\longrightarrow0,
 \qquad N>2.}
\end{equation}
On the odd branch all positive degrees are odd and formal interface degrees remain odd, so no quadratic diagonal arises;
the formula is the case $\mu'=0$ of \eqref{fp:actual-U}.

\paragraph{Arbitrary couplings of the prior marginals.}
If both marginals of $(W,W')$ are the Gaussian network priors above and
$\sup_d\mathbb E\|B(W)\|_2^2<\infty$, then
\[
 \left|\mathbb E\langle B(W),\Pi_NF_{W'}-\mathcal Q_N(W')\rangle\right|
 \leq(\mathbb E\|B(W)\|_2^2)^{1/2}
 (\mathbb E\|\Pi_NF_{W'}-\mathcal Q_N(W')\|_2^2)^{1/2}=o(1).
\]
This uses only marginal errors; it does not claim that higher-layer weights remain independent under an arbitrary coupling.
The closure keeps all raw families and the $U_N$ corrections; different families with the same bottom degree must be added,
and one may not keep only the single graph with degree $m$ everywhere.


\section{Ordered trees, physical slot assignments and separation of full profiles}\label{fp:history}\label{app:cuts}
\paragraph{Fix the order, then sum.}
Throughout this section fix the depth $L$, the degree $N=m^L$ and the width ratio $k_s\asymp d$. For odd activations take odd $m\ge5$; for piecewise-linear activations take even $m\ge6$. Constants may depend on these fixed parameters. $K_d$ denotes a fixed power of the logarithm; set $\delta_d=K_d/\sqrt d$, $\kappa_d=K_d/\sqrt d$. All finite-degree high-probability bounds below hold on the raw Gaussian marginal events of the source and of the teacher separately; no independence between the two networks is assumed.

An \emph{ordered profile} chooses a positive-degree monomial at every raw computation-tree node and orders the children of that node. Physical slots are the $N$ input-vector positions, ordered after full expansion. Fix a bijection from these positions to the $N$ positions of the balanced teacher; all estimates are first proved for this fixed bijection. Tensor symmetrization is a fixed finite average of at most $N!$ terms, so it may not be used first to hide a coherent contraction inside one term. Different node occurrences may choose different admissible degrees, and repeated numerical neuron labels must be distinguished from distinct tree-node occurrences.

In phase $h$ a physical slot represents a completed teacher block of height $h-1$; in phase $1$ the slots are the input positions. Each teacher leaf to be completed contains $m$ current slots. After fully contracting these $m$ slots, the source node carries a \emph{literal factor} $M_{ia}$: it keeps the source neuron label $a$, a completed teacher block label $i$, and all raw edge normalizations. It is not a new independent Gaussian edge. A literal factor produced in the current phase can only be a spectator output; it becomes an active slot only in the next phase.

\subsection{Teacher context and the lowest cut}
\paragraph{Keep the full cotensor.}
Denote the current raw homogeneous teacher coefficient by $C$. Completed output positions use the identity map; uncompleted positions use bounded homogeneous power column maps. After fixing an active teacher row $i$, the norm of all remaining physical outputs is therefore at most $K_d\|C(e_i,\cdot)\|_{\rm HS}$, and
\begin{equation}\label{fp:teacher-slices}
 \sum_i\|C(e_i,\cdot)\|_{\rm HS}
 \le \sqrt{k}\|C\|_{\rm HS}\le K_d\sqrt d.
\end{equation}
The cross metric is always placed on the source side; a teacher tensor that has already been adaptively contracted is never claimed to be a new raw homogeneous coefficient.

At each stopping node of a source cut, the raw subtree strictly below it, together with the installed outputs, is treated as one complete bounded feature map. Each literal factor strictly above the stopping height keeps its own coefficient output mode. Even when several outputs share one source label, the per-output identity copy modes are kept. Structural copies are not distinct source parents, and their diagonals may not be deleted. All nodes at the current height must stop \emph{before} the feature maps and identity copies in order to use the right-rotation estimate for a new color.

\begin{lemma}[Lowest nontrivial meeting]\label{fp:lowest-cut}
Assume that every old literal factor, before its first active use, has passed the full-cotensor test for the case with remaining factors. If all cuts containing two active source parent positions have been excluded by the corresponding analytic estimates, then every surviving new teacher-leaf contraction consumes, at exactly one raw source parent node, $m$ complete raw child factors that are linear in the current variables.
\end{lemma}
\begin{proof}
Mark the $m$ source slots of a teacher leaf to be completed. Push the marks upward in increasing raw source height; marks starting higher wait. Before a mark meets its first multiplicative branch, the expression it traverses is a product of raw normalized matrices ending in an existing metric, hence still a bounded row map. A parent node is called single-factor if, after deleting one selected linear factor, no physical factor or spectator output remains. Marks reaching the same parent are grouped first; a node reached by two marks at once is necessarily not single-factor.

If every node met is single-factor, sum its labels exactly and continue along the raw linear path. Otherwise stop at the non-single-factor node of smallest height. Every selected factor arriving there is a complete linear factor: had it crossed a nonlinear branch below, it would have stopped earlier. If the marks lie at two or more parent positions, choose the lowest non-single-factor position as the new color coordinate. Push the other, lower single path up to that height; if it meets a branch earlier, stop at that earlier branch instead. A higher active old literal factor has passed its first-use test, so it carries no remaining factor and cannot be an ancestor of another active mark; new literal factors of this phase are only spectator outputs. Rebuilding the common cut from the raw tree and stopping before all feature maps at this height, as above, gives the required new-color coordinate cut.

By assumption the mixed-parent case is excluded, so all marks meet at the same parent node. If two of them came from the same nonlinear raw child, they would have met and stopped at a strictly lower node, a contradiction. Hence they come from $m$ distinct complete raw children, each linear in the current physical variables.
\end{proof}

When the cut order is at least three, the fresh-coordinate slice bound of Lemma~\ref{fp:core-fresh} directly completes the mixed-parent test of this lemma. Precisely, the selected row contracts to $\eta_i,\zeta_i$, each with bounded $\ell^2$ norm. The non-single-factor position $a$ still has a nonempty feature after deleting one factor, so the orthogonal label of $a$ is kept. If the other position is single-factor, the norm of the full source cotensor is at most
\begin{equation}\label{fp:two-mark-vector}
 K_d\left(\sum_a|\eta_{ia}|^2
       \|E(e_a,\zeta_i,\cdot)\|_{\rm HS}^2\right)^{1/2}
 \le K_d/d.
\end{equation}
If both sides keep features, keep both labels $a,b$ and use
\begin{equation}\label{fp:two-mark-coordinate}
 K_d\left(\sum_{a,b}|\eta_{ia}\zeta_{ib}|^2
       \|E(e_a,e_b,\cdot)\|_{\rm HS}^2\right)^{1/2}
 \le K_d/d.
\end{equation}
Numerical label collisions are handled first on the raw coefficients, then the complete feature maps are applied, and only then is the result paired with \eqref{fp:teacher-slices}; the total error is at most $K_d/\sqrt d=o(1)$. Collision masks cannot be passed through vector contractions.

\subsection{Closed histories of odd profiles}
\paragraph{Parity controls the cut order.}
For odd activations all raw degrees are odd. Each raw substitution changes the number of outputs by an even number $r-1$, and a previously verified complete birth replaces $m$ factors by one, changing it by the even number $-(m-1)$. Hence every nonempty stopped subtree has an odd total number of current physical and spectator outputs, and stopping it as one mode preserves parity. The induction uses only births accepted \emph{earlier}; it does not assume that the birth being tested is complete. Consequently an expanded cut with at least two distinct active parent positions has order at least three. Every strictly higher literal output is still counted separately; structural copies count only as output modes, not as independent active parent positions.

If $q$ births have occurred in the current phase at a node of raw degree $r$, its factor list contains $q$ literal factors and $r-qm$ unconsumed raw children, so $r-q(m-1)$ factors are visible, still an odd number. If an old literal factor still carries remaining factors at its first active use, at least two factors remain. If the current teacher leaf also visits another source column, the split-hit estimate of Lemma~\ref{fp:core-split-hit} excludes the term; if it hits only this column, use Lemma~\ref{fp:core-hybrid-Rge2} when the upper source coefficient has order at least two, and the at-least-two-extra-factors estimate of Lemma~\ref{fp:core-hybrid-R1} when it has order one. Extra literal factors are also kept separately; ordinary nonlinear factors keep the raw Gaussian readout structure. Hence a surviving active old literal factor is always the full feature of its birth node. This supplies the first-use hypothesis of Lemma~\ref{fp:lowest-cut}.

The metric of a complete birth is the entrywise product of $m$ linear row kernels. If the kernels are $\rho^{(j)}$, then by the bounded $\ell^2$ norms of all rows and columns and H\"older's inequality,
\begin{equation}\label{fp:metric-schur}
 \sup_i\sum_a\left|\prod_{j=1}^m\rho^{(j)}_{ia}\right|
 +\sup_a\sum_i\left|\prod_{j=1}^m\rho^{(j)}_{ia}\right|\le K_d.
\end{equation}
so the metric has bounded absolute row and column sums. The next raw Gaussian teacher matrix $A$ satisfies $\max|A_{ij}|\le\kappa_d$, so $\max|AM|\le K_d\kappa_d$ holds even when the metric depends on both networks; no conditional independence is needed.

Each path push exposes one edge of a finite raw tree and each birth completes one teacher block, so every phase terminates after finitely many steps. At the end of a phase the metrics above serve as the physical terminals of the next phase; the first-use test is performed before continuing with the lowest cut. The final birth cannot carry a positive-degree remaining factor, since all $N$ slots are used.

A surviving block of height $h$ needs at least $h$ raw source nodes along every path: $h=1$ needs a branching node, the $m$ complete children of $h$ are blocks of height $h-1$, and since there is no remaining factor at active use, the parent birth lies strictly above those children. Any unary delay only increases the source height. If a source of depth at most $L$ produces a teacher block of height $L$, equality holds on every path, there is no unary delay, and every node has degree exactly $m$. So among odd profiles only the complete balanced $m$-ary profile survives; its expression keeps all raw coefficients and normalizations and is exactly the homogeneous comparator.

\subsection{The permanent mark path and the single-factor boundary of two modes}\label{fp:permanent-mark}
\paragraph{Do not contract the mark path early.}
Consider an exact two-mode state
\begin{equation}\label{fp:marked-state}
 \sum_{a,b}E_{ab}\,\ell_a\otimes Q_b.
\end{equation}
where $\ell$ is a complete current teacher terminal factor and $Q$ is a complete subtree of raw height at most $L-1$; all unused physical and spectator factors are kept. $E$ must come from a raw quadratic cut or its raw rectangular transport, not from an arbitrary adaptive matrix; its raw reconstructed coefficients satisfy the core estimate of Lemma~\ref{fp:core-closure}. In particular $\|E\|_{\rm HS}\le K_d$, and every nontrivial flattening is at most $\delta_d$. When a shorter source root leaves one literal factor and one complete child, this is $E_{ab}=d_aU_{ab}$, with $\|d\|_2\le K_d$, $\|d\|_\infty\le\kappa_d$, and $U$ a raw normalized Gaussian edge.

Keep the $a$ mode permanently, mark the unique path from $\ell$ to the final teacher root, and never close the teacher blocks on this path. The other teacher blocks are processed in increasing height; all their source slots lie in $Q$. Each reconstructed core still keeps $a$, and the current new color stops before all copies and feature maps. Hence the core order is at least two when one descendant position is selected, and at least three when two distinct descendant positions are selected.

Therefore the first-use test inside $Q$ can always use the source-flattening version; the norm of the full source cotensor for a fixed teacher row is at most $K_d\delta_d(\delta_d+\kappa_d)$. The remaining slots of all marked paths are contained in this cotensor, and after multiplying by \eqref{fp:teacher-slices} the term vanishes. The mixed-parent test uses \eqref{fp:two-mark-vector}--\eqref{fp:two-mark-coordinate}. A surviving unmarked birth still consumes complete children. It cannot occur at $a$: all its active slots lie in the unique child $Q$; if that child is nonlinear the marks meet earlier inside it, and if it has only one current slot it cannot supply $m$ slots.

Every unmarked birth of height $h$ needs raw height at least $h$ inside $Q$. If the height of $Q$ is less than $L-1$, the $m-1$ unmarked children of height $L-1$ of the final root cannot be completed. If the height of $Q$ equals $L-1$, these $m-1$ children must all be born at the unique raw root $b$ of $Q$. Let its complete metric be $M^{(j)}$, the final teacher matrix $A$, and set $G^{(j)}=AM^{(j)}$. By \eqref{fp:metric-schur}, the columns $\ell^2$ of each $G^{(j)}$ have bounded norm and maximal entry at most $K_d\kappa_d$.

Keep the marked teacher child fully physical. Fix a final teacher row $i$; it has bounded norm. After removing the $m-1$ known literal factors, the full source cotensor is
\[
 R_b=\sum_aE_{ab}\ell_a\otimes Q_b^{\rm res},\qquad
 \|R_b\|\le K_d\|E_{:b}\|_2,
 \quad\sum_b\|R_b\|\le K_d\sqrt d\|E\|_{\rm HS}.
\]
Here every cofactor is kept as is, the column map of $\ell$ is bounded, and the columns of $Q_b^{\rm res}$ have bounded norm. The maximum of the normalized final readout is $O(d^{-1/2})$; using the column $\ell^2$ bounds of two of the $G$ and the entry bounds of the remaining $m-3$, the full scalar pairing is at most
\begin{equation}\label{fp:marked-terminal}
 K_d\|E\|_{\rm HS}\kappa_d^{m-3}=o(1).
\end{equation}
This step uses the $m-1$ literal factors simultaneously and never multiplies a small scalar kernel directly against an arbitrary cotensor. It settles the single-factor/non-single-factor two-mode boundary at full remaining depth, as well as the boundary of one literal factor plus one complete child for shorter arbitrary roots.

\subsection{Piecewise-linear profiles and the coherent quadratic branch}
\paragraph{Strong induction on the source height.}
The ordinary polynomial $P$ contains only degree one and positive even degrees. Fix an even $m\ge6$. We prove simultaneously: for every height $t<L$, every root given by an arbitrary positive-degree polynomial $A$, and every ordered monomial choice $P$ at the nodes below, the pairing vanishes under every fixed slot assignment; and among ordinary $P$ profiles of height $L$ only the full $m$ profile survives. The induction allows different children to choose different monomials, so as to cover the exact ordered profiles after fusion.

Lowest cuts of order at least three have been handled; a cut of order one cannot have two distinct active parent positions. For order two, the case of one single-factor and one non-single-factor position is settled by \eqref{fp:marked-state}--\eqref{fp:marked-terminal}. If both are non-single-factor, both must be new modes of the same height: a higher active literal factor has no remaining factor after its first use, and a higher spectator literal output would be a third mode. The cut is therefore the raw quadratic upper coefficient $D$ before the feature maps below are applied.

For its off-diagonal part, permutation symmetry, bounded Hilbert second moments and Gaussian hypercontractivity at fixed degree give
\[
 \max_{a\ne b}|D_{ab}|\le K_d/d.
\]
Keeping both labels, the full source cotensor is bounded by
\[
 K_d\Bigl(\sum_{a\ne b}|\eta_{ia}\zeta_{ib}D_{ab}|^2\Bigr)^{1/2}
 \le K_d/d,
\]
and pairing with the teacher slices makes it vanish. This argument uses only coordinate/coordinate estimates, not a general second-order coordinate/vector gain.

For the diagonal, the raw quadratic coefficient estimate \eqref{fp:actual-diag} gives
\begin{equation}\label{fp:coherent-diagonal}
 \operatorname{Diag}D=\lambda_s I/\sqrt{k_s}+E_s,
 \qquad\mathbb E\|E_s\|_{\rm HS}^2=o(1).
\end{equation}
The error is first controlled with the complete bounded source feature maps. The main term identifies the two raw parent labels, multiplies their two selected root monomials, and keeps all ordered children and exactly the same teacher slot assignment. This gives a raw prefix profile of height $s<t$ and readout $1/\sqrt{k_s}$, multiplied by the exact deterministic coefficient $\lambda_s$. Its root may have either parity, so the strengthened induction class above is needed. Completed literal metrics can be expanded back into raw slot contractions; these are identities, so fusion produces no truncated subtrees. If $s=0$, only input degree two remains, which cannot contribute $N>2$.

A two-mode cut may also consist of a higher literal output and a new mode: it has no same-color diagonal deletion; in the active case it falls under the permanent-mark argument, and in the spectator case there is only one active parent mode. If both modes are higher pure literal outputs, then in an uncompleted state each represents at most $m^{L-1}$ raw slots, their total is less than $N$, and they cannot form the whole profile; if one already has height $L$, it alone uses all slots. Structural copies are likewise counted by the number of raw slots they represent and cannot be treated as new independent parents.

After $q$ births in the current phase, an ordinary even-degree node has $r-qm$ raw children left, an even number. An active literal factor with positive remainder therefore has at least two further ordinary factors, or other separately kept literal factors. Hence the product-cotensor estimate applies when the upper core has order one, and the source-flattening estimate when it has order at least two. The case of several pure literal factors is handled by Lemma~\ref{fp:core-scalar-literal}. An arbitrary shorter root need not have even degree; write its factor list after first use as $q$ literal factors and $k=r-qm\ge0$ raw children: if $q+k\ge3$ there are at least two extra factors; if $q+k=2$, the only possibilities are two literal factors, or one literal factor plus one complete raw child, the latter handled by the permanent-mark argument; if $q+k=1$, it is a complete literal factor with no remainder. This exhausts all root degrees.

Start the induction at $t=1$; all quadratic diagonals fall at strictly lower heights, and when there is no lower term of positive height, input degree two excludes it. After excluding mixed cuts and remaining factors at first use, the finite schedule makes every birth complete and remainder-free when active. Height $t<L$ cannot produce a final block of height $L$, so its pairing vanishes. Then treat the main profile of height $L$; all fused terms have already been excluded in the shorter induction. Equality of heights forces the whole tree to have no unary delay and every node to have degree $m$; the surviving term is exactly the raw homogeneous term.

\subsection{From deterministic events to arbitrary prior couplings}\label{fp:coupling-uniform}
\paragraph{Control the approximation errors of each prior marginal first.}
At fixed degree and depth, the numbers of ordered profiles, slot assignments, raw cuts, fusion heights and history shapes are independent of $d$. Impose all raw Gaussian events above on each marginal separately, and only then allow the cross metrics to depend on both networks. On the common good event every estimate is a deterministic $o(1)$. On the bad event use bounded marginal higher moments of fixed-degree coefficients and H\"older's inequality; the bad-event probability can be made of any required polynomial order. This control does not depend on the coupling and is therefore uniform over all couplings of the prior marginals.

Let $T_W$ be the full $m$ balanced homogeneous input chaos term corresponding to the activation, with complete scalar coefficients, and denote by $\mathcal Q_N(V)$ the finite raw polynomial family of \eqref{fp:actual-U}. The profile-wise results above give
\begin{equation}\label{fp:profile-pairing}
 \sup_{\pi\in\Gamma(\mathsf P_W,\mathsf P_V)}
 \mathbb E_\pi\left|\langle T_W,\mathcal Q_N(V)-T_V\rangle\right|
 \longrightarrow0.
\end{equation}
Moreover $T_W$ lies entirely in the $N$-th input chaos, so $\langle T_W,F_V\rangle=\langle T_W,\Pi_NF_V\rangle$. By \eqref{fp:actual-final}, the closure error of each marginal is controlled by Cauchy--Schwarz:
\[
 \mathbb E_\pi|\langle T_W,\Pi_NF_V-\mathcal Q_N(V)\rangle|
 \le(\mathbb E\|T_W\|^2)^{1/2}
      (\mathbb E\|\Pi_NF_V-\mathcal Q_N(V)\|^2)^{1/2}=o(1).
\]
Combining with \eqref{fp:profile-pairing} gives the required global function comparison
\begin{equation}\label{fp:global-function-pairing}
 \boxed{\displaystyle
 \sup_{\pi\in\Gamma(\mathsf P_W,\mathsf P_V)}
 \mathbb E_\pi|\langle T_W,F_V-T_V\rangle|\longrightarrow0.}
\end{equation}

\section{Global matching of homogeneous comparator networks}\label{hm:section}\label{app:matching}
\paragraph{Homogeneous networks and the global matching statement}
This section treats only homogeneous networks with activations $c_m z^m$, $c_m=(m!)^{-1/2}$. Fix $L$ and an integer $m\ge5$, let $k_0=d$, $cd\le k_s\le Cd$, and define
\[
 \begin{gathered}
 h^{(0)}=x,\qquad h^{(s)}=c_m(W^{(s)}h^{(s-1)}/\sqrt{k_{s-1}})^{\odot m},\\
 \mathcal H_W(x)=k_L^{-1/2}\sum_i h_i^{(L)}(x),\qquad
 p=\sum_{s=1}^L k_{s-1}k_s.
 \end{gathered}
\]
All weights of each network have independent standard Gaussian priors; the two networks may be coupled arbitrarily. Let $P_s(Y)$ be the downstream output after replacing $h^{(s)}$ by a formal variable $Y$, $q_s=m^{L-s}$, $C_s=C_{q_s}(P_s)$. Here $P(Y)=\langle C_r(P),Y^{\otimes r}\rangle$, and tensors carry the inherited Hilbert norm; if $P=\sum_{|\nu|=r}p_\nu Y^\nu$, then
\begin{equation}\label{hm:coeff-norm}
 \|C_r(P)\|^2=\sum_{|\nu|=r}\frac{\nu!}{r!}|p_\nu|^2.
\end{equation}
In particular $C_L=k_L^{-1/2}{\bf1}$.

\begin{theorem}[Finite-label global matching for homogeneous networks]\label{hm:global-matching}
There exist a family of deterministic exact symmetry maps $\mathcal G_d$ of the homogeneous network, and constants $K$ depending only on the fixed depth, degree and width ratio, such that
\begin{equation}\label{hm:main}
 \log|\mathcal G_d|=O(d\log d)=o(p),\qquad
 \frac1p\mathbb E\min_{g\in\mathcal G_d}\|W-gV\|_F^2
 \le K\mathbb E\|C_0(W)-C_0(V)\|^2+o_d(1).
\end{equation}
The remainder is uniform over all couplings with the prior marginals above. If $X\sim N(0,I_d)$, the right-hand side may also be replaced by
$K\mathbb E\|\mathcal H_W-\mathcal H_V\|_{L^2_X}^2+o_d(1)$.
\end{theorem}

\subsection{The common frame and downstream mode energies}
\paragraph{Common frame and downstream mode energies}
We use the multiplicity frame theorem (frame bound~\eqref{tg:frame}) in the following form. For unit row vectors $a_i$, $k=O(d)$ with maximal self-coherence $\eta=O(\sqrt{\log d/d})$ and fixed $r$, the pure tensors
\[
 \operatorname{Sym}\bigotimes_i a_i^{\otimes m\nu_i},\quad |\nu|=r,
\]
and the single-deletion tensors
\[
 \operatorname{Sym}\bigotimes_i a_i^{\otimes(m\nu_i-\mathbf1_{i=j})},
 \quad |\nu|=r,\quad \nu_j>0,
\]
each form a frame with constant lower and upper bounds (no repetition among distinct $(\nu,j)$). That is, the squared norm of any linear combination is equivalent to the sum of squared coefficients up to dimension-free constants. Allowing $1/2\le\|a_i\|\le2$ changes only the constants. A sufficient small-perturbation bound is
$K[\eta+\sum_{h=1}^r k^h\eta^{mh-2}]=o(1)$.
The global symmetrization here cannot be replaced by a colored tensor space. All other comparisons below are proved in this section.

Let $B_{s,i}=\partial_iP_s$. On the usual Gaussian concentration events for row norms, coherence and column square sums of each layer,
\begin{equation}\label{hm:derivative}
 a/k_s\le\|C(B_{s,i})\|^2\le b/k_s\quad\hbox{for all }s,i.
\end{equation}
To spell out the consequence of the common frame, the last layer is $B_{L,i}=k_L^{-1/2}$. Setting $w_j=W^{(s+1)}_{j:}/\sqrt{k_s}$, the chain rule gives
\[
 B_{s,i}(Y)=mc_m\sum_j w_{ji}(w_j\cdot Y)^{m-1}
 B_{s+1,j}\big(c_m(w\cdot Y)^{\odot m}\big).
\]
After expanding $B_{s+1,j}$, each term is exactly a single-deletion tensor of label $(\nu,j)$; at fixed degree \eqref{hm:coeff-norm} is equivalent to the ordinary monomial coefficient norm. The deletion frame therefore gives
\[
 \|C(B_{s,i})\|^2\asymp\sum_j w_{ji}^2\|C(B_{s+1,j})\|^2
 \asymp k_{s+1}^{-1}\sum_j w_{ji}^2\asymp k_s^{-1}.
\]
This is a deterministic frame estimate; it does not require the downstream derivative to be independent of the current row. If $H_s$ is the first-mode Gram matrix of $C_s$, the exact differential identity gives
\begin{equation}\label{hm:mode}
 (H_s)_{ii}=q_s^{-2}\|C(\partial_iP_s)\|^2,\qquad
 a'/k_s\le(H_s)_{ii}\le b'/k_s.
\end{equation}

\subsection{Control of dense coefficients and coordinate collisions}
\paragraph{Marginal events that survive adaptive deletion}
For any row matrix $A$ set $U_Ae_i=a_i^{\otimes m}$. Since
$U_A=(A^T)^{\otimes m}J$, where $Je_i=e_i^{\otimes m}$ is an isometric embedding,
$\|U_A\|_{\rm op}\le\|A\|_{\rm op}^m$. Each homogeneous substitution is a tensor power of $U_A$ followed by the contraction map $\operatorname{Sym}$. Recursing from $\|C_L\|=1$, $\|C_s\|$ is controlled by a fixed power product of operator norms of finitely many normalized Gaussian matrices. Fixed-order moments of Gaussian matrix operator norms are uniformly bounded, so
\begin{equation}\label{hm:moments}
 \sup_d\mathbb E\|C_s\|^u<\infty\quad(0<u<\infty).
\end{equation}
For a square-free monomial $Y_I=\prod_{i\in I}Y_i$, $|I|=q_s$, write its coefficient as $p_I$ and set $c_I=p_I/\sqrt{q_s!}$. By \eqref{hm:coeff-norm}, $\sum_Ic_I^2$ is the squared tensor norm of this sector. Column permutation invariance of the first upper matrix and \eqref{hm:moments} give
$\mathbb E c_I^2\le C\binom{k_s}{q_s}^{-1}\le Ck_s^{-q_s}$.
Each $c_I$ is a Gaussian polynomial of degree at most $D_s=m+\cdots+m^{L-s}$, so Gaussian polynomial hypercontractivity
$\|c_I\|_u\le(u-1)^{D_s/2}\|c_I\|_2$ ($u\ge2$) gives: taking $u=A\log d$, applying Markov at $e\|c_I\|_u$ and a union bound over at most $k_s^{q_s}$ terms, for any prescribed $M$,
\begin{equation}\label{hm:dense}
 \mathbb P\left\{\max_I|c_I|^2>
 C_M (\log d)^{D_s}k_s^{-q_s}\right\}\le d^{-M}.
\end{equation}
The last layer is a deterministic vector and is treated separately.

We also need the \emph{whole} downstream tensor to differ from its square-free truncation by $o(1)$. Here is a direct operator proof. Let $T=\operatorname{Sym}_{mr}U_A^{\otimes r}B$, $\|A\|_{\rm op}\le K$, $\max_{ij}|a_{ij}|\le\kappa$. Contract two specified input slots of $T$ with $e_j$: if both slots come from the same $m$ block, the resulting block map is
\[
 (A^T)^{\otimes(m-2)}J\operatorname{diag}(a_{ij}^2),
\]
with norm at most $K^{m-2}\kappa^2$; if they come from different blocks, each block has norm at most $K^{m-1}\kappa$. The other blocks have bounded operator norms, and at fixed degree symmetrization involves finitely many slot assignments. Hence
\[
 \|T(e_j,e_j,\cdot,\ldots,\cdot)\|\le C\kappa^2\|B\|.
\]
The squared norm of the coordinate-collision tensor is at most $\binom{mr}{2}$ times the sum of squared norms of all coordinate-diagonal slices, so
\begin{equation}\label{hm:collision}
 \|(I-\mathrm{SF})T\|^2\le C n\kappa^4\|B\|^2=o(1)\|B\|^2,
 \qquad \kappa=O(\sqrt{\log d/d}),\ n\asymp d.
\end{equation}
Applying this to the first downstream substitution of $C_s$ and using the bounded upper tensor norm gives
$\|(I-\mathrm{SF})C_s\|=o(1)$. This does not yield a coordinatewise mode lower bound for the square-free truncation; below we first restore the full tensor and then use \eqref{hm:mode}.

All these marginal events can be chosen simultaneously with failure probability tending to zero. Row norms, operator norms, coherence, maximal coordinates and column square sums are invariant under signed row and column permutations; the downstream tensor is only reindexed with coordinate sign changes, so its norm, maximal absolute coefficient, mode diagonal bounds and collision sectors are invariant too. Hence these events still hold after an adaptively chosen internal alignment; there is no need to claim that the transformed network is still independent Gaussian.

\subsection{Exact conversion between global symmetric and grouped tensors}
\paragraph{Keep radius factors and factorial normalizations}
At layer $s$ write $w_i=W^{(s)}_{i:}/\sqrt{k_{s-1}}=r_i a_i$, $\|a_i\|=1$,
$R=\operatorname{diag}(r_i^m)$, $\overline C_s=R^{\otimes q_s}C_s$. The exact substitution is
\begin{equation}\label{hm:substitution}
 C_{s-1}=c_m^{q_s}\operatorname{Sym}_{mq_s}U_A^{\otimes q_s}\overline C_s.
\end{equation}
A uniform radius error $\max_i|r_i-1|\le\delta_d\to0$ implies
$\|\overline C_s-C_s\|\le C\delta_d\|C_s\|$.
Mode energies, maximal coefficients and collision estimates still hold, since each tensor coordinate is multiplied by a fixed number of $1+o(1)$ radius factors.

Fix $q=q_s$ and for $|I|=q$ define
\[
 Q_I=\sqrt{q!}\operatorname{Sym}_q\bigotimes_{i\in I}a_i^{\otimes m},\qquad
 E_I=\sqrt{(mq)!/(m!)^q}\operatorname{Sym}_{mq}\bigotimes_{i\in I}a_i^{\otimes m}.
\]
$Q_I$ lies in the grouped space and $E_I$ in the actual global symmetric space. A square-free coordinate basis vector becomes $Q_I$ under the grouped lift and then $\sqrt{\lambda_{m,q}}E_I$ under \eqref{hm:substitution}, where
\begin{equation}\label{hm:lambda}
 \lambda_{m,q}=q!/(mq)!.
\end{equation}
The two normalizations make the cross-network pairing terms that do not split the $m$ slots exactly equal.

\subsection{Adaptive mixed-slot comparison}
\paragraph{Coefficient deletion and two-sided Schur estimates}
Given two arrays of unit rows $A,B$ with self-coherence at most $\eta_d$ and operator norm at most $K$, write $\rho_{ij}=\langle a_i,b_j\rangle$. Join pairs with $|\rho_{ij}|>\tau$ to form a bipartite graph; squared row and column sums are at most $K^2$, so degrees are at most $D=\lceil K^2/\tau^2\rceil$. On each side, join two vertices sharing a highly correlated neighbor; the resulting graph has maximal degree at most $D^2$. A $q$-element set containing no edge of this graph is called good. For good $I$ and any vertex on the other side, at most one correlation in $I$ has absolute value exceeding $\tau$.

If $v_1+\cdots+v_q=m\ge3$, then
\begin{equation}\label{hm:mixed-column}
 \sum_j\prod_{a=1}^q|\rho_{i_a j}|^{v_a}\le
 \begin{cases}K^2,&v\text{ has exactly one nonzero entry},\\C\tau,&v\text{ has at least two nonzero entries and }I\text{ is good}.
 \end{cases}
\end{equation}
For the second line, partition according to the position $a$ of the largest correlation in the positive support. All other correlations are at most $\tau$. If $v_a\ge2$, keep the square of the largest term; at least one remaining power gives $\tau$. If $v_a=1$, keep one power of it and of another positive-support term; the remaining $m-2\ge1$ powers give $\tau$, and the kept product is summed by Cauchy--Schwarz. Finitely many partitions complete the proof.

Each entry of the actual cross-Gram is given by a $q\times q$ nonnegative integer matrix $(e_{ab})$ with all row and column sums equal to $m$. Precisely,
\[
 \langle E_I^A,E_J^B\rangle
 =\sum_{\substack{e_{ab}\ge0\\\sum_b e_{ab}=m,\ \sum_a e_{ab}=m}}
 \frac{(m!)^q}{\prod_{a,b}e_{ab}!}
 \prod_{a,b}\rho_{i_a j_b}^{e_{ab}}.
\]
Indeed, a fixed pairing count matrix $e$ has $(m!)^{2q}/\prod_{a,b}e_{ab}!$ slot arrangements; the symmetric tensor inner product first divides by $(mq)!$, and the product of the two normalizations multiplies by $(mq)!/(m!)^q$, giving the coefficient above. The permutation profiles of $e=mP$ give exactly the grouped cross-Gram; every other profile has at least one mixed column. Fix a good $I$ and sum absolute values over good $J$; after relaxing to all ordered $J$ the sum factorizes over columns, and \eqref{hm:mixed-column} gives $C\tau$. Exchanging the two sides gives absolute column sums $C\tau$. Schur's inequality then gives
\begin{equation}\label{hm:cross-schur}
 \|K_{\rm act}-K_{\rm grp}\|_{\rm op}\le C\tau
\end{equation}
on the good sets of both sides.

Any graph of maximal degree $B$ has at most $C_qBk^{q-1}$ $q$-element sets containing an edge, so \emph{simultaneously for all such graphs}
\[
 \sum_{I\ {\mathrm{bad}}}|c_I|^2\le C_qBk^{q-1}\max_I|c_I|^2.
\]
Taking $\tau=d^{-\zeta}$, $0<\zeta<1/4$, \eqref{hm:dense} bounds the deletion energy by
$C d^{4\zeta-1}(\log d)^{D_s}=o(1)$.
The graphs may depend on both complete networks and even on the coefficients; conditional independence is never applied to an adaptive graph here.

The self grouped Gram is close to the identity because $k\eta_d^m=o(1)$. The self actual square-free Gram is also close to the identity: non-identity pairings with the same support have error $O(\eta_d)$; if the new support adds $h$ labels, there are at least $mh$ pairings of unequal labels, and the absolute row sum is at most
$C[\eta_d+\sum_{h=1}^q k^h\eta_d^{mh}]=o(1)$.
Hence deleting coefficients costs a $o(1)$ norm error for both the actual and the grouped tensors. The collision sector of the full tensor is recovered by \eqref{hm:collision} and the bounded lift operator. Combining \eqref{hm:cross-schur} with the self-Gram comparison gives
\begin{equation}\label{hm:sym-group}
 \|C_{s-1}^A-C_{s-1}^B\|^2
 =\lambda_{m,q}\|G_A-G_B\|^2+o(1),\qquad
 G_A=U_A^{\otimes q}\overline C_s^A,\quad G_B=U_B^{\otimes q}\overline C_s^B.
\end{equation}
All errors are uniform on the joint typical event. For $q=1$ no slot splits and the identity is exact, $\lambda_{m,1}=1/m!$.

\subsection{Projection loss of raw rows and matching without square-root loss}
\paragraph{Projection loss controlled by the mode diagonal lower bound}
Let $U=U_A,V=U_B,M=U^*V$, $\rho_i=\max_j|\rho_{ij}|$. For a fixed threshold $t>0$ the highly correlated set has size at most $K^2/t^2$; the self Gram of this subset has operator norm at most $1+K^2\eta_d/t^2$, so
\[
 \sum_j|M_{ij}|\le \rho_i^{m-2}(1+K^2\eta_d/t^2)+K^2t^{m-2}.
\]
Taking $t=\eta_d^{1/m}$ (any positive null sequence when $\eta_d=0$) and transposing the same argument gives a uniform $e_d\to0$ with
\begin{equation}\label{hm:cross-power}
 \sum_j|M_{ij}|\le\rho_i^{m-2}+e_d,\qquad
 \sum_i|M_{ij}|\le1+e_d.
\end{equation}
Weighted Cauchy--Schwarz gives
$MM^*\preceq(1+e_d)\operatorname{diag}(\rho_i^{m-2}+e_d)$: concretely, for any $x$, sum
$(\sum_iM_{ij}x_i)^2\le(\sum_i|M_{ij}|)\sum_i|M_{ij}|x_i^2$ over $j$.
Since $U^*U,V^*V=I+o_{\rm op}(1)$, the projection formula onto $\operatorname{ran}V$ gives
\[
 U^*(I-P_{\operatorname{ran}V})U
 =U^*U-M(V^*V)^{-1}M^*
 \succeq\operatorname{diag}(1-\rho_i^{m-2})-o(1)I.
\]
Apply this orthogonal projection to the first slot of $G_A-G_B$; the Gram of the remaining slots is $I+o(1)$. Hence for the first mode $H_A$ of the full coefficient tensor,
\begin{equation}\label{hm:projection-loss}
 \|G_A-G_B\|^2\ge\sum_i(1-\rho_i^{m-2})(H_A)_{ii}-o(1).
\end{equation}
Combining with \eqref{hm:mode} gives
\begin{equation}\label{hm:projective-match}
 k_s^{-1}\sum_i(1-\rho_i^{m-2})\le C\|G_A-G_B\|^2+o(1).
\end{equation}

Fix $\rho_0\in(2^{-1/2},1)$. Sufficiently small self-coherence ensures that no row has absolute correlation greater than $\rho_0$ with two rows of the other side, since otherwise the Gram upper bound of the two rows would imply $2\rho_0^2\le1+\eta_d$. So the above-threshold matches form a partial bijection; complete it arbitrarily to a permutation and orient each row so that the diagonal correlation is nonnegative. The fraction of teacher rows not covered by the partial bijection is at most
$[k_s(1-\rho_0^{m-2})]^{-1}\sum_i(1-\rho_i^{m-2})$.
For matched rows, $1-|\rho_{i,\pi(i)}|^m\le C_m(1-\rho_i^{m-2})$; the loss for other rows is at most one. Setting
$\ell_i=1-\langle a_i,b_i\rangle^m\ge0$ gives
\begin{equation}\label{hm:row-loss}
 k_s^{-1}\sum_i\ell_i+k_s^{-1}\sum_i\|a_i-b_i\|^2
 \le C\|G_A-G_B\|^2+o(1).
\end{equation}
The fixed threshold is used only to construct the bijection; the error bound stays linear and no square root is taken.

\subsection{Diagonal lift variation and recovery of formal coefficients}
\paragraph{Coordinatewise energy keeps the matching error linear}
After alignment $M_{ii}=1-\ell_i$. By \eqref{hm:cross-power} the absolute row and column sums are at most $1+e_d$. For any real vector $x$,
\[
 \langle x,(2I-M-M^*)x\rangle
 \le2\sum_i\ell_ix_i^2+\sum_{i\ne j}|M_{ij}|(x_i^2+x_j^2)
 \le4\sum_i\ell_ix_i^2+2e_d\|x\|^2.
\]
Hence, with $D=U_A-U_B$, the $o(1)$ error of the self Gram gives
\begin{equation}\label{hm:lift-loewner}
 0\preceq D^*D\preceq4\operatorname{diag}(\ell_i)+o(1)I.
\end{equation}
For any symmetric coefficient tensor $C$ of order $q$, assume $\|C\|\le C_0$ with first mode $H_{ii}\le C_0/k_s$. Telescope $U_A^{\otimes q}-U_B^{\otimes q}$ slot by slot; the squared triangle inequality gives one factor $q$, and the same mode energy of each slot gives another $q$. The Gram of the other lift factors is at most $(1+o(1))I$, so
\begin{equation}\label{hm:lift-linear}
 \|(U_A^{\otimes q}-U_B^{\otimes q})C\|^2
 \le q^2(1+o(1))\operatorname{tr}(D^*DH)
 \le C_q k_s^{-1}\sum_i\ell_i+o(1).
\end{equation}
This step uses the diagonal energy upper bound to keep finer control than a global operator-difference estimate.
Applying it to the aligned student $\overline C_s^B$, together with
\[
 U_A^{\otimes q}(\overline C_s^A-\overline C_s^B)
 =G_A-G_B+(U_B^{\otimes q}-U_A^{\otimes q})\overline C_s^B
\]
and the lower bound on $U_A^*U_A=I+o(1)$, gives
\begin{equation}\label{hm:recover}
 \|C_s^A-C_s^B\|^2\le C\|G_A-G_B\|^2+o(1)
 \le C\|C_{s-1}^A-C_{s-1}^B\|^2+o(1).
\end{equation}
Removing the radius factors adds only a $O(\delta_d^2)$ squared error, already included in the remainder.

\subsection{Sign propagation and the fixed readout constraint}
\paragraph{Controlling sign errors at an odd last layer}
Write the signed row permutation of an internal layer $s<L$ as $S_s=D_sP_s$, where $D_s$ has diagonal entries $\pm1$. The hidden features become
$R_sh^{(s)}$, $R_s=D_s^{m\bmod2}P_s$. Replacing the next-layer matrix by $V^{(s+1)}R_s^{-1}$ at the same time leaves the whole homogeneous function exactly unchanged. Odd $m$ propagates signs to the columns of the next layer, while even $m$ propagates only the permutation; formal downstream coefficients are reindexed and change sign accordingly, and the typical events above are preserved.

Last-layer permutations preserve the all-ones readout. For even $m$, last-layer row signs also preserve the readout, but for odd $m$ last-layer sign flips are not allowed. In this case rows are first oriented only in the analytic representation: if the matching sign is $\varsigma_i$ and the student radius is $t_i$, represent the original student tensor with oriented rows; its last-layer coefficient is
$\overline c_i^B=\varsigma_i t_i^m/\sqrt{k_L}$, and the teacher's is
$\overline c_i^A=r_i^m/\sqrt{k_L}$.
Putting the signs into the coefficients leaves the original function unchanged, and the mode diagonal bounds still hold. Applying \eqref{hm:lift-linear}--\eqref{hm:recover} to $q=1$ gives
\[
 k_L^{-1}\sum_i(r_i^m-\varsigma_i t_i^m)^2
 \le C\|C_{L-1}^A-C_{L-1}^B\|^2+o(1).
\]
The radii tend to one uniformly, so each negative sign contributes a fixed positive cost, and the fraction of negative signs is controlled linearly by the right-hand side. Only the last-layer permutation is actually applied: rows with positive sign are controlled by \eqref{hm:row-loss}, rows with negative sign have distance bounded by a constant, and their number is controlled by the bound above. Hence the average row distance after the last-layer permutation still satisfies the same linear estimate.

\subsection{Induction on depth and arbitrary prior couplings}
\paragraph{Completing the proof of the theorem}
Let $E_0=\|C_0(W)-C_0(V)\|^2$. Choose signed matchings in the internal layers from the bottom up, propagate them to the columns of the next layer, and denote by $E_s$ the formal tensor error after the current alignment. Internal symmetry transformations preserve the current $C_{s-1}$, and later alignments do not change the lower layers already completed. Hence \eqref{hm:sym-group}, \eqref{hm:row-loss}, \eqref{hm:recover} give at each step
\[
 E_s\le CE_{s-1}+e_d,\qquad
 \frac1{k_s}\sum_i\|w_{s,i}-(gv)_{s,i}\|^2\le CE_{s-1}+e_d,
 \qquad e_d\to0.
\]
The radius error is absorbed by
$\|r_ia_i-t_ib_i\|^2\le C\|a_i-b_i\|^2+C\delta_d^2$
The last layer uses the previous subsection, and induction over the fixed depth shows that the average normalized row distance of every layer is at most $C_{L,m}E_0+o(1)$.
Recovering the raw weights requires multiplying by $k_{s-1}$, so the weighted sum with $k_{s-1}k_s/p$ gives the deterministic version of \eqref{hm:main} on the joint typical event.

The family of maps is
\[
 (gV)^{(1)}=S_1V^{(1)},\qquad (gV)^{(s)}=S_sV^{(s)}R_{s-1}^{-1}\quad(s\ge2),
\]
For odd $m$, $S_L$ is restricted to permutations; the other $S_s$ may be signed permutations. Hence
\[
 \log|\mathcal G_d|\le\sum_s\{\log(k_s!)+k_s\log2\}=O(d\log d).
\]
Fixing a label order breaks all ties and yields a measurable matching choice. Under any coupling of the prior marginals, a union bound over the two separate typical events makes the joint failure probability vanish; the two networks need not be independent. On the bad event take the identity alignment,
\[
 p^{-1}\min_g\|W-gV\|^2\le2p^{-1}(\|W\|^2+\|V\|^2).
\]
The squared norms of normalized Gaussian weights have uniformly bounded second moments; Cauchy--Schwarz turns the bad-event expectation into $o(1)$, uniformly in the coupling. Taking expectations proves \eqref{hm:main}. Finally, Gaussian chaos of top input degree $N=m^L$ satisfies
\[
 N!\|C_0(W)-C_0(V)\|^2
 =\|\Pi_N(\mathcal H_W-\mathcal H_V)\|_{L^2_X}^2
 \le\|\mathcal H_W-\mathcal H_V\|_{L^2_X}^2,
\]
which proves the function-distance version. \hfill$\square$


\paragraph{Parameter comparison for the actual functions.}
The comparison~\eqref{fp:global-function-pairing} and the polarization identity of the main text turn control of the actual function distance into control of the homogeneous coefficient distance, to which Theorem~\ref{hm:global-matching} applies. Row signs of the even-degree comparator are used only for finite-label parameter alignment and need not preserve the values of the actual activation.

\section{Derivation of the predictive KL information lower bound}\label{app:information}
\paragraph{The coefficient $1/2$ comes from an exact scalar integral.}
This appendix concerns the Gaussian channel only; $f=F_W(X)$ in all densities is the true output. Choose a reference density $p_0=N(0,\Delta)$ with average reference KL $\E f^2/(2\Delta)<\infty$; write the density of $p_u(y)=N(u,\Delta)$. For $0<s<1$ set
\[
q_s=(1-s)p_0+sp_u,\quad
S_s=\frac{p_u-p_0}{q_s},\quad
\alpha_s=\frac{p_0p'_u}{q_s^2},\quad h_s=\frac{sp'_u}{q_s}.
\]
Fix $Z=(\mathcal D_n,X)$. The second derivative of the conditional mutual information $F_Z(s)=I(f;Y_s\mid Z)$ is
\[
F_Z''(s)=\E[\operatorname{Var}(S_s(f,Y_s)\mid Z,Y_s)\mid Z].
\]
obtained by writing $F_Z$ as a difference of two KLs relative to the reference density and differentiating. On any closed subinterval $[a,b]\subset(0,1)$, $|\partial_sq_s|/q_s$ is uniformly bounded and
\[
\int\frac{(p_u-p_0)^2}{q_s}\,dy\le s^{-1}+(1-s)^{-1},
\]
so two derivatives may be taken. To check the zero endpoint, set $\Phi(v)=v\log v-v+1$, $v=p_f/p_0$. Convexity gives
\[
0\le\frac{\Phi(1+s(v-1))}{s}\le\Phi(v),\qquad
\frac{\Phi(1+s(v-1))}{s}\longrightarrow0.
\]
Since $\E\KL(p_f\|p_0)=\E f^2/(2\Delta)<\infty$, dominated convergence gives $F_Z(s)/s\to0$, so $F_Z(0)=F_Z'(0+)=0$. At the one endpoint, revealing the independent mixture coin and data processing give $F_Z(s)\le sF_Z(1)$; the joint and product laws converge in total variation, and lower semicontinuity of relative entropy gives the reverse limit. Integrating the nonnegative second derivative and using Tonelli gives
\begin{equation}\label{eq:mixture}
\epsilon_{\log,d,n}=\int_0^1(1-s)\E\operatorname{Var}(S_s\mid \mathcal D_n,X,Y_s)\,ds.
\end{equation}

\subsection{Gaussian integration by parts}
\paragraph{The test field may depend on the observed auxiliary response.}
Let $a=\nabla_Wf$, $S_D=\nabla_W\log p(\mathcal D_n\mid W)$,
$\delta b=\operatorname{div}b-W\cdot b$. For a field $b_s(W,X,Y_s)$ that does not depend on the training samples, define
\[
G_s=\delta b_s+b_s^TS_D+h_s b_s^Ta.
\]
Integrating by parts in the parameters with the observations fixed, and setting
$e_s=\E[S_s\mid \mathcal D_n,X,Y_s]-S_s$, gives
\[
\E[e_sG_s]=\E[b_s^T\alpha_sa].
\]
Conditionally on $W,X,Y_s$, the training score has mean zero and covariance $n\I_W$, so
\[
\E G_s^2=n\E[b_s^T\I_Wb_s]+\E(\delta b_s+h_sb_s^Ta)^2.
\]
Integrate both identities under $(1-s)ds$ and denote the target on the right, the training quadratic term and the auxiliary quadratic term by $\mathcal T,\mathcal Q,\mathcal J$. By \eqref{eq:mixture} and Cauchy--Schwarz,
\begin{equation}\label{eq:tqj}
\epsilon_{\log,d,n}\ge\frac{\mathcal T^2}{n\mathcal Q+\mathcal J}.
\end{equation}
Here the posterior estimator is a function of the observed variables during integration by parts, so its derivative with respect to the parameter integration variable vanishes.

\subsection{The regularized field and the exact half coefficient}
\paragraph{Regularize the true population Fisher matrix.}
Let $M=(\I_W+rI/p)^{-1}$, $c=Ma$, $b_s=\alpha_sc$. For fixed $P=p_u(y)>0,B=p_0(y)>0$,
\[
\int_0^1\frac{(1-s)B^2}{(B+s(P-B))^3}\,ds=\frac1{2P}.
\]
Integrating over $y$ gives
\[
\int_0^1(1-s)\int q_s\alpha_s^2\,dy\,ds
=\frac12\int\frac{(p'_u)^2}{p_u}\,dy=\frac1{2\Delta}.
\]
Hence
\begin{equation}\begin{aligned}
\mathcal T&=\tfrac12\E\tr(M\I_W)=\tfrac p2(1-D_d(r)),\\
\mathcal Q&=\tfrac12\E\tr[(M\I_W)^2]\le p/2.
\end{aligned}\end{equation}

We also need to control the field derivative, keeping the parameter dependence of $M$. Set $v=a/\sqrt\Delta$, $\phi=M^{1/2}v$; then $\E_X\phi\phi^T\preceq I$. For $w_i=M^{1/2}\partial_iv$, Bessel's inequality gives
\[
\norm{\E_Xw_i\phi^T}_{\rm HS}^2\le\E_X\norm{w_i}^2.
\]
Since $\partial_i\I=\E[(\partial_iv)v^T+v(\partial_iv)^T]$ and
$\partial_ic=M\partial_ia-M(\partial_i\I)Ma$, summing gives
\[
\E\norm{Dc}_{\rm HS}^2\le C(p/r)^2\E\norm{\nabla^2f}_{\rm HS}^2.
\]
The Gaussian divergence satisfies
$\E(\delta g)^2\le\E(\norm g^2+\norm{Dg}_{\rm HS}^2)$.
Absorb the auxiliary density into $g_s=\sqrt{q_s}b_s=\psi_sc$, where $\psi_s=\sqrt{q_s}\alpha_s$. The exact product rule is
\[
\delta b_s+h_sb_s^Ta=q_s^{-1/2}\delta g_s+\tfrac12h_sb_s^Ta.
\]
Besides the half-coefficient integral above, we need
\begin{align*}
\int_0^1\frac{(1-s)sB^2}{q_s^4}\,ds&=\frac1{6P^2},\\
\int_0^1\frac{(1-s)s^2B^2}{q_s^5}\,ds&=\frac1{12P^3}.
\end{align*}
These follow by differentiating the first integral once and twice in $P$. Hence
\[
\int_0^1(1-s)\int q_s\alpha_s^2h_s^2\,dy\,ds
=\frac1{12}\int\frac{(p'_u)^4}{p_u^3}\,dy.
\]
Since $\partial_u\psi_s=B[p''_uq_s^{-3/2}-(3s/2)(p'_u)^2q_s^{-5/2}]$, the three identities above and the sum-of-squares inequality give
\[
\int_0^1(1-s)\int(\partial_u\psi_s)^2\,dy\,ds
\le\int\frac{(p''_u)^2}{p_u}\,dy
+\frac38\int\frac{(p'_u)^4}{p_u^3}\,dy.
\]
For the Gaussian channel the two integrals on the right are $2/\Delta^2$ and $3/\Delta^2$. Moreover
\[
\E_X\norm{c}^2/\Delta=\tr(M^2\I)\le p^2/(4r).
\]
By $Dg_s=\psi_sDc+(\partial_u\psi_s)ca^T$,
$\norm M\le p/r$, the fourth moment of the gradient and the $Dc$ bound above, $\E\norm{Dg_s}_{\rm HS}^2$ after the outer integral is at most $Cp^2/r^2$, and $\E\norm{g_s}^2$ is at most $Cp^2/r$. The divergence bound and the final $h_s$ term together give
$\mathcal J\le Cp^2/r^2$. Substituting into \eqref{eq:tqj} gives \eqref{eq:lower-smooth}. Lemma~\ref{an:sobolev} justifies the integration by parts by smooth compactly supported truncation followed by a limit in the Gaussian Sobolev norm.

\subsection{Piecewise-linear case: smoothing only the test gradient}
\paragraph{Keep the true likelihood; smooth only the test gradient.}
Take $a_{\dsm}=\nabla f_{\dsm}$ from \eqref{eq:smoothing} and define
\[
\I_{\dsm}=\E_Xa_{\dsm} a_{\dsm}^T/\Delta,
\quad M_{\dsm}=(\I_{\dsm}+rI/p)^{-1},
\quad b_{s,\dsm}=\alpha_s(f,y)M_{\dsm} a_{\dsm}.
\]
In $\alpha_s,q_s,h_s$ it is still the actual $f$. The same derivative computation gives
$\mathcal J_{\dsm}\le Cp^2/(r^2\dsm)$.
The gradient difference and the regularized trace bound give
\begin{equation}\begin{aligned}
|\mathcal T_{\dsm}-\tfrac p2(1-D_{d,\dsm}(r))|
&\le Cp\sqrt{\dsm/r},\\
\mathcal Q_{\dsm}&\le\tfrac p2(1+C\sqrt{\dsm/r})^2.
\end{aligned}\end{equation}
By the trace-norm inequality for rank-one kernels,
$\E\norm{\I-\I_{\dsm}}_*\le C\sqrt{\dsm}$; the resolvent identity and the two inverse norms $\le p/r$ give
\[
|D_{d,\dsm}(r)-D_d(r)|\le C\sqrt{\dsm}/r.
\]
Substituting into \eqref{eq:tqj} and taking the positive part of the target lower bound gives \eqref{eq:lower-pl}. The $1/\dsm$ in the denominator and the two smoothing losses in the numerator come from the estimates above. This applies the test-field method to nonsmooth networks while keeping the true likelihood.


\end{document}
